\documentclass[11pt]{article}

\usepackage[margin=1in]{geometry}
\usepackage{amsmath,amssymb,amsthm,mathtools}
\usepackage{booktabs,tabularx}
\usepackage[authoryear]{natbib}
\usepackage[colorlinks=true,allcolors=blue]{hyperref}
\usepackage{microtype}

\newtheorem{theorem}{Theorem}[section]
\newtheorem{proposition}[theorem]{Proposition}
\newtheorem{corollary}[theorem]{Corollary}
\newtheorem{lemma}[theorem]{Lemma}
\theoremstyle{definition}
\newtheorem{definition}[theorem]{Definition}
\theoremstyle{remark}
\newtheorem{remark}[theorem]{Remark}

\newcommand{\Acal}{\mathcal{A}}
\newcommand{\Ccal}{\mathcal{C}}
\newcommand{\Ocal}{\mathcal{O}}
\newcommand{\Pcal}{\mathcal{P}}
\newcommand{\Hcal}{\mathcal{H}}
\newcommand{\Fcal}{\mathcal{F}}
\newcommand{\Gcal}{\mathcal{G}}
\newcommand{\Qcal}{\mathcal{Q}}
\newcommand{\Rcal}{\mathcal{R}}
\newcommand{\E}{\mathbb{E}}
\newcommand{\Prob}{\mathbb{P}}
\newcommand{\KL}{\mathrm{KL}}
\newcommand{\I}{\mathrm{I}}
\newcommand{\Hh}{\mathrm{H}}
\newcommand{\argmin}{\mathop{\mathrm{arg\,min}}}
\newcommand{\argmax}{\mathop{\mathrm{arg\,max}}}
\newcommand{\givena}{\,\|\,}
\newcommand{\ellp}{\ell}
\newcommand{\wt}{\widetilde}
\newcommand{\barxi}{\bar{\xi}}
\newcommand{\one}{\mathbf{1}}
\newcommand{\problembox}[1]{%
\begin{center}
\fbox{\parbox{0.93\textwidth}{#1}}
\end{center}}

\title{Sufficient Conditions for Semantic Convergence in Interactive Learning}
\author{}
\date{}

\begin{document}

\maketitle

\begin{abstract}
We seek sufficient conditions under which an interactive learner's posterior concentrates on the \emph{interventional semantic class} $[p^\star]$ of an unknown computable environment --- the programs that induce the same action-conditioned law as $p^\star$ under every admissible intervention, and the strongest target any history-based learner can reach. The nested chain $\{p^\star\}\subseteq[p^\star]\subseteq R_\pi(p^\star)\subseteq R_{h_t}(p^\star)$ collects the indistinguishability classes of four \emph{observers} at increasing observational power, and each learning rule in the paper is an observer-promotion attempt. The paper studies the \emph{algorithmic free energy}: a Gibbs-form variational functional whose reference measure is the Kolmogorov-complexity-weighted universal interactive prior of \citet{Solomonoff1964a,Solomonoff1964b,Kolmogorov1965,Hutter2005}. Its belief-side minimization is exactly Bayesian updating with the KL divergence forced by convex duality, and its zero-cost uniform-preference action term recovers the information-gain criterion shared by Bayesian experimental design, compression-based curiosity, and active-inference curiosity. That generic information-gain route can still fail: a substituted finite-support near-miss freezes the posterior on $[p^\star]$ at prior mass for every finite horizon. The positive theorem therefore targets the right observer geometry. Under the universal prior, Bayes/Gibbs belief, semantic Bhattacharyya separation, and divergent cumulative policy mass on separating actions, the posterior on $[p^\star]$ converges almost surely to one; under bounded log-likelihood ratios, explicit finite-time high-probability lower bounds follow. The support condition is sharp, and it is attained by the class-against-complement selector, the exact Gibbs kernel lift, a finite-latent surrogate under explicit soundness assumptions, and divergent full-support exploration. Corollary~\ref{cor:compact-action-kernel} extends the exact kernel route to compact metric action spaces under action-continuity and a full-support Borel reference measure.

A public Lean 4 repository at \url{https://github.com/Populustremuloides/semantic_convergence} contains Lean counterparts for all 106 manuscript items over the explicit concrete foundation used in the repo, together with generated coverage, proof-shape, and per-declaration axiom audits. On the belief and variational side, \path{SemanticConvergence.def_afe}, \path{SemanticConvergence.def_two_observer_functional}, and \path{SemanticConvergence.def_kernel_functional} implement the repaired generalized-KL / Gibbs objectives, while \path{SemanticConvergence.lem_variational}, \path{SemanticConvergence.lem_kl_necessity}, \path{SemanticConvergence.prop_two_observer_minimizer}, and \path{SemanticConvergence.prop_kernel_functional_minimizer} certify the corresponding exact variational identities and minimizer theorems. Around Section~\ref{sec:main}, the Lean route now separates the soft theorem from the stronger zero-out special case: \path{SemanticConvergence.HellingerConvergenceSpine} packages the paper's Bhattacharyya martingale hypotheses, \path{SemanticConvergence.thm_separating_support_convergence_hellinger_spine} proves almost-sure posterior-share convergence from that spine, and \path{SemanticConvergence.thm_kernel_functional_minimizer_convergence_of_hellinger_spine} composes exact kernel-functional minimization with that convergence spine. The existing support/rate declarations \path{SemanticConvergence.thm_separating_support_convergence}, \path{SemanticConvergence.thm_separating_support_rate}, and \path{SemanticConvergence.cor_separating_support_finite_time} remain as a concrete zero-out support route with finite-horizon bounds. Thus the Lean development now verifies the central promise at the Hellinger-spine interface; the remaining first-principles exactness work is to construct that spine internally from the Bayes/Gibbs update, semantic separation, and realized policy support, rather than supplying the spine as a named hypothesis.
\end{abstract}

\section{Introduction}
\label{sec:intro}

Prediction admits many successes that explanation does not. An interactive learner may reproduce a realized transcript, forecast the next observation along the policy it happens to follow, and reduce uncertainty about the latent program in general, and still be wrong about what the environment would do under an intervention it never tried. In an interactive problem, that residual ambiguity is the question worth asking: whether the learner's internal inference converges to the environment law that is correct under every admissible intervention.

Four targets organize the answer for computable interactive environments. The \emph{history-compatible set} $R_{h_t}(p^\star)$ consists of programs that reproduce the realized transcript; the \emph{policy-predictable set} $R_\pi(p^\star)$, of those that induce the same law as $p^\star$ under the policy the learner actually follows; the \emph{interventional semantic class} $[p^\star]$, of those that induce the same action-conditioned law under every admissible intervention; and the innermost target, the singleton $\{p^\star\}$, is the truth. These form a nested hierarchy,
\[
\{p^\star\}
\;\subseteq\;
[p^\star]
\;\subseteq\;
R_\pi(p^\star)
\;\subseteq\;
R_{h_t}(p^\star)
\;\subseteq\;
\Pcal,
\]
and all three inclusions between the truth and the full compatibility set are strict in general (Theorem~\ref{thm:strict-hierarchy}). An observable-quotient theorem (Theorem~\ref{thm:factor-through-quotient}) makes $[p^\star]$ the theorem-level ceiling: any target a history-based learner can reliably recover is constant on $[p^\star]$, so concentration on $\{p^\star\}$ itself is impossible whenever a semantic class contains more than one program. The four sets are the indistinguishability classes of four \emph{observers} at increasing observational power (Definition~\ref{def:observer}), from reading only the realized transcript up to reading the source code; the paper's theorems turn on whether a learning rule \emph{promotes} the realized observer toward the behavioral one.

Three existing traditions supply natural ingredients for a learner aimed at that ceiling. Solomonoff induction and AIXI supply a universal prior over computable interactive environments, weighted by Kolmogorov complexity \citep{Solomonoff1964a,Solomonoff1964b,Kolmogorov1965,Hutter2005}, but no variational functional. The free energy principle and active inference supply the Gibbs variational form and a risk-plus-epistemic-value action rule \citep{Friston2010,FristonEtAl2017Curiosity,ParrFriston2019}, but with a stipulated parametric prior whose justification has been debated on philosophical grounds \citep{ColombWright2021}. Bayesian experimental design, information-theoretic bounded rationality, and compression-based curiosity supply closely related information-seeking action criteria, again without a universal-prior reference measure \citep{Lindley1956,ChalonerVerdinelli1995,MacKay1992,OrtegaBraun2013,Schmidhuber1991,Schmidhuber2010}. No existing framework couples these three ingredients into a single variational object. We formulate the \emph{algorithmic free energy} --- the Gibbs-form variational functional whose reference measure is the Kolmogorov-complexity-weighted universal interactive prior rather than a stipulated parametric prior (Definition~\ref{def:afe}) --- and on it the three traditions unify on both sides of the inference problem. The Bayes/Gibbs posterior is its unique minimizer (Lemma~\ref{lem:variational}), and the KL-based divergence is forced within the natural convex class (Lemma~\ref{lem:kl-necessity}); minimizing the algorithmic free energy is not an approximation of Bayesian updating but is Bayesian updating exactly. Under zero action cost and uniform preferences, the resulting expected free energy coincides with the information-gain criterion of Bayesian experimental design, compression-based curiosity, and active-inference curiosity (Corollary~\ref{cor:efe-specialization}).

The unifications do not close the problem. The same Bayes/Gibbs/expected-free-energy architecture can fail to concentrate its posterior on $[p^\star]$ (Theorem~\ref{thm:afe-near-miss}): for every horizon $T$ and every nontrivial prior mass $\alpha_0\in(0,1/2)$ there is a problem instance, using a substituted finite-support prior, on which the posterior on $[p^\star]$ is frozen at $\alpha_0$ for all $t\le T$. The theorem extends the AIXI near-miss tradition \citep{LeikeHutter2015,Leike2016,LeikeEtAl2016} and locates the failure specifically at the level of the interventional semantic class. Because expected-free-energy action is exactly the unifying generic information-gain criterion, the failure is a failure of generic information-seeking itself, not of any peculiar implementation choice. In observer language, the near-miss construction makes the AFE principle's realized policy induce an observer strictly coarser than the behavioral observer, and the frozen-posterior conclusion is the Bayesian shadow of that strict refinement (Theorem~\ref{thm:observer-promotion-failure}).

\problembox{\textbf{Problem.} Give sufficient conditions on the learner's prior, belief rule, and action rule under which the posterior mass on $[p^\star]$ converges almost surely to one.}

The paper's answer is a master support theorem: under the universal prior, Bayes/Gibbs belief, semantic separation, and divergent cumulative policy mass on separating actions, the posterior on $[p^\star]$ converges almost surely to one (Theorem~\ref{thm:separating-support-convergence}, Section~\ref{sec:main}). The exact minimizing class-action kernel of the kernel lift is the strongest exact variational realization of that condition under the standing finite-action setup (Theorem~\ref{thm:kernel-semantic-convergence}); the class-against-complement rule is the canonical selector realization (Theorem~\ref{thm:semantic-convergence}); and Corollary~\ref{cor:compact-action-kernel} extends the exact kernel route to compact metric action spaces under an explicit action-continuity hypothesis and a full-support Borel reference measure. The quantitative companion Theorem~\ref{thm:separating-support-rate} gives explicit finite-time posterior lower bounds under bounded log-likelihood ratios. The condition is sharp: zero separating support makes uniform recovery impossible (Corollary~\ref{cor:support-necessary}), and summable separating support still admits failure with positive probability (Theorem~\ref{thm:summable-support-insufficient}). A divergent cumulative exploration budget delivers the same guarantee on its own (Theorem~\ref{thm:exploration-floor-behavioral}).

Section~\ref{sec:setup} fixes the interactive setting. Section~\ref{sec:hierarchy} constructs the four-set hierarchy, proves the observable-quotient theorem, and witnesses the three strict inclusions. Section~\ref{sec:functional} introduces the two-observer functional $\mathcal J_t^{\omega_B,\omega_A}$ together with its exact action-realizing kernel lift $\mathcal J_{t,\mathrm{ker}}^{\omega_B,\omega_A}$, and proves the meeting-point theorem. Section~\ref{sec:belief} identifies the belief-side minimizer at the canonical belief observer. Section~\ref{sec:semantic-action} identifies the exact kernel action law and the canonical class-against-complement selector realization at the canonical action observer, and proves that the separation geometry transfers through decodable observation channels. Section~\ref{sec:main} proves the master support theorem for semantic recovery, its explicit finite-time quantitative companion, the exact kernel and selector realizations, the compact-action extension, and the matched converse theorems showing what fails when separating support is absent or only summable. Section~\ref{sec:self-ref} replaces the uncomputable behavioral observer by a finite-time proxy and recovers the main theorem's guarantee under eventual isolation (deterministic splitting or divergent exploration). Section~\ref{sec:near-miss} reads the algorithmic free energy principle as the boundary example of the functional outside the action-admissible range, and proves the near-miss. Section~\ref{sec:implications} records a differentiable surrogate loss (Section~\ref{sec:amortized-surrogate}) that inherits the master theorem's guarantee over a finite latent class, together with several single-problem consequences. Sections~\ref{sec:discussion} and~\ref{sec:conclusion} draw consequences.

\section{Formal Setup}
\label{sec:setup}

Let $\Acal$ be a finite action alphabet and $\Ocal$ a countable observation alphabet. Write $\Delta(X)$ for the set of probability measures on a measurable space $X$, and equip every finite or countable set appearing below with its discrete $\sigma$-algebra. For each $t\ge 0$ write
\[
\Hcal_t:=(\Acal\times\Ocal)^t,
\qquad
\Hcal_0:=\{\varepsilon\},
\qquad
\Hcal:=\bigsqcup_{t\ge 0}\Hcal_t,
\]
with the product and disjoint-union $\sigma$-algebras. Time is discrete. At step $t$ the learner emits an action $a_t\in\Acal$ and receives an observation $o_t\in\Ocal$. Histories are
\[
h_t=(a_1,o_1,\dots,a_t,o_t),
\qquad
h_0=\varepsilon.
\]
A history $h_t$ is an element of $\Hcal_t$. A policy is a sequence of measurable stochastic kernels
\[
\pi_t:\Hcal_{t-1}\rightsquigarrow \Acal,
\qquad
h_{t-1}\longmapsto \pi_t(a\mid h_{t-1})\text{ measurable for each }a\in\Acal.
\]
For $S\subseteq\Acal$ write
\[
\pi_t(S\mid h_{t-1}):=\sum_{a\in S}\pi_t(a\mid h_{t-1}).
\]
Actions are interventions and observations are evidence. Under an environment $p$ and a policy $\pi$, write $(A_t,O_t)_{t\ge 1}$ for the resulting action-observation process and
\[
H_t:=(A_1,O_1,\dots,A_t,O_t)
\]
for the random history at time $t$.

Fix a universal prefix machine $U$. Let $\Pcal$ be the countable set of programs in the prefix-free domain of $U$ that induce interactive environment models. Each $p\in\Pcal$ defines an action-conditioned chronological law
\[
\mu_p(o_{1:t}\givena a_{1:t})
=
\prod_{i=1}^{t}\mu_p(o_i\mid h_{i-1},a_i),
\qquad
\sum_{o\in\Ocal}\mu_p(o\mid h_{t-1},a_t)=1.
\]
For $h_t=(a_1,o_1,\dots,a_t,o_t)\in\Hcal_t$, the induced finite-horizon law under $(p,\pi)$ is
\[
\Prob_{p,\pi}(H_t=h_t)
:=
\mu_p(o_{1:t}\givena a_{1:t})\prod_{i=1}^t \pi_i(a_i\mid h_{i-1}),
\]
which defines a probability measure $\Prob_{p,\pi}(H_t\in\cdot)\in\Delta(\Hcal_t)$. The consistent family of finite-horizon laws determines the full trajectory law $\Prob_{p,\pi}$ on the product space $(\Acal\times\Ocal)^{\mathbb N}$. Call a history-action pair $(h_t,a)$ \emph{reachable under $p$} if $\mu_p(o_{1:t}\givena a_{1:t})>0$. Throughout, kernel equality $\mu_p=\mu_{p'}$ denotes agreement at pairs reachable under at least one of $p,p'$: $\mu_p(\cdot\mid h,a)=\mu_{p'}(\cdot\mid h,a)$ whenever $(h,a)$ is reachable under $p$ or under $p'$; jointly unreachable pairs have no observable consequence under any policy. This reachable-pair equality defines an equivalence relation on $\Pcal$; write
\[
p\equiv p'
\iff
\mu_p=\mu_{p'},
\qquad
\Ccal:=\Pcal/\!\equiv,
\qquad
[\cdot]:\Pcal\to\Ccal
\]
for the semantic quotient set and quotient map.

For a prior $w$ on $\Pcal$ and a history $h_t$ with positive $w$-induced marginal likelihood, the Bayes posterior over programs is
\[
q_t(p\mid h_t)
:=
\frac{w(p)\mu_p(o_{1:t}\givena a_{1:t})}
{\sum_{r\in\Pcal} w(r)\mu_r(o_{1:t}\givena a_{1:t})}.
\]
For any subset $\Rcal\subseteq\Pcal$ write
\[
q_t(\Rcal\mid h_t):=\sum_{p\in\Rcal} q_t(p\mid h_t).
\]
For a semantic class $c\in\Ccal$, write
\[
q_t(c\mid h_t):=q_t([\cdot]^{-1}(c)\mid h_t)=\sum_{p\in c} q_t(p\mid h_t).
\]
The prior $w$ pushes forward to $\Ccal$ as
\[
\bar w(c):=w([\cdot]^{-1}(c))=\sum_{p\in c}w(p).
\]
For a semantic class $c=[p]\in\Ccal$, write $\mu_c:=\mu_p$ for the common action-conditioned law induced by any representative of $c$; this is well-defined by definition of $\equiv$.
For a continuation $(a,o)$ after $h_t$ with positive predictive normalizing constant,
\[
q_{t+1}(p\mid h_t,a,o)
:=
\frac{q_t(p\mid h_t)\mu_p(o\mid h_t,a)}
{\sum_{r\in\Pcal} q_t(r\mid h_t)\mu_r(o\mid h_t,a)}.
\]
For a semantic class $c\in\Ccal$, write
\[
q_{t+1}(c\mid h_t,a,o):=q_{t+1}([\cdot]^{-1}(c)\mid h_t,a,o)=\sum_{p\in c} q_{t+1}(p\mid h_t,a,o).
\]

\section{The Knowability Hierarchy}
\label{sec:hierarchy}

This section replaces a list of informal learning objectives with four nested sets of programs and proves that each inclusion is strict. The strictness results are the structural spine of the paper.

\subsection{Four nested sets}

\begin{definition}[History-compatible set]
\label{def:history-compat}
For a target program $p^\star\in\Pcal$ and a history $h_t\in\Hcal_t$, the \emph{history-compatible set} is
\[
R_{h_t}(p^\star)
:=
\bigl\{p\in\Pcal\,:\,\mu_p(o_{1:t}\givena a_{1:t})=\mu_{p^\star}(o_{1:t}\givena a_{1:t})\bigr\}.
\]
\end{definition}

\begin{definition}[Policy-predictable set]
\label{def:policy-pred}
For a target program $p^\star\in\Pcal$ and a policy $\pi$, the \emph{policy-predictable set} is
\[
R_\pi(p^\star)
:=
\bigl\{p\in\Pcal\,:\,\Prob_{p,\pi}(H_T\in\cdot)=\Prob_{p^\star,\pi}(H_T\in\cdot)\text{ for every }T\bigr\}.
\]
\end{definition}

\begin{definition}[Interventional semantic class]
\label{def:int-sem-class}
For a target program $p^\star\in\Pcal$, the \emph{interventional semantic class} of $p^\star$ is its $\equiv$-class
\[
[p^\star]
=
[\cdot](p^\star)
=
\bigl\{p\in\Pcal\,:\,p\equiv p^\star\bigr\}
=
\bigl\{p\in\Pcal\,:\,\mu_p=\mu_{p^\star}\bigr\}.
\]
\end{definition}

Thus $\Ccal$ is the quotient set of programs by semantic equivalence, and two programs lie in the same class exactly when they agree on the entire action-conditioned environment law.

\begin{definition}[Observer]
\label{def:observer}
An \emph{observer} is a pair $\omega=(\Omega_\omega,V_\omega)$ of a measurable space $\Omega_\omega$ and a measurable view map $V_\omega:\Pcal\to\Omega_\omega$. Observer $\omega$ is \emph{refined by} $\omega'$, written $\omega\le\omega'$, if there is a measurable factorization $\phi:\Omega_{\omega'}\to\Omega_\omega$ with $V_\omega=\phi\circ V_{\omega'}$. The \emph{$\omega$-indistinguishability class} of $p^\star$ is the fiber $V_\omega^{-1}(V_\omega(p^\star))$: the programs an adversary could substitute for $p^\star$ without the $\omega$-observer detecting the swap.
\end{definition}

Four canonical observers structure the hierarchy, each equipped with the image $\sigma$-algebra making its view map measurable: the \emph{syntactic observer} $\omega_{\mathrm{syn}}$ with $\Omega_{\omega_{\mathrm{syn}}}:=\Pcal$ and $V_{\omega_{\mathrm{syn}}}(p):=p$, reading the source code; the \emph{behavioral observer} $\omega_{\mathrm{behav}}$ with $\Omega_{\omega_{\mathrm{behav}}}:=\{\mu_p:p\in\Pcal\}$ and $V_{\omega_{\mathrm{behav}}}(p):=\mu_p$, running every policy and seeing every induced law; the \emph{policy observer} $\omega_\pi$ with $\Omega_{\omega_\pi}:=\{\Prob_{p,\pi}:p\in\Pcal\}$ and $V_{\omega_\pi}(p):=\Prob_{p,\pi}$, seeing only the law under a fixed policy; and the \emph{history observer} $\omega_{h_t}$ with $\Omega_{\omega_{h_t}}:=\{\mu_p(o_{1:t}\givena a_{1:t}):p\in\Pcal\}\subseteq[0,1]$ and $V_{\omega_{h_t}}(p):=\mu_p(o_{1:t}\givena a_{1:t})$, reading only the realized transcript. Their indistinguishability classes of $p^\star$ are exactly $\{p^\star\}$, $[p^\star]$, $R_\pi(p^\star)$, and $R_{h_t}(p^\star)$.

\begin{remark}
\label{rem:four-observers}
The hierarchy is the partition-refinement chain of four observers at increasing observational power: as the observer is promoted, the indistinguishability class shrinks. The nesting, strict-separation, and observable-quotient results below are the three consequences of this single reading. The near-miss pair (Theorems~\ref{thm:afe-near-miss} and~\ref{thm:observer-promotion-failure}) and the main theorem (Theorem~\ref{thm:semantic-convergence}) turn on whether a learning rule \emph{promotes} the realized policy observer $\omega_\pi$ to the behavioral observer $\omega_{\mathrm{behav}}$: the former show that a natural rule does not, the latter gives a sufficient condition under which a stricter rule does.
\end{remark}

\begin{lemma}[Nesting]
\label{lem:nesting}
For any $p^\star\in\Pcal$, any policy $\pi$, and any history $h_t$ with $\Prob_{p^\star,\pi}(H_t=h_t)>0$,
\[
\{p^\star\}\subseteq [p^\star]\subseteq R_\pi(p^\star)\subseteq R_{h_t}(p^\star)\subseteq \Pcal.
\]
\end{lemma}

\begin{proof}
A swap the stronger observer cannot detect, the weaker observer cannot detect. If $p\in[p^\star]$ then $\mu_p=\mu_{p^\star}$, so the induced $\pi$-law matches $p^\star$'s, giving $[p^\star]\subseteq R_\pi(p^\star)$. If $p\in R_\pi(p^\star)$ then $\Prob_{p,\pi}(H_t=h_t)=\Prob_{p^\star,\pi}(H_t=h_t)$; combining with the factorization
\[
\Prob_{p,\pi}(H_t=h_t)=\mu_p(o_{1:t}\givena a_{1:t})\prod_{i=1}^t \pi_i(a_i\mid h_{i-1}),
\]
and dividing through by the policy factor (positive by assumption), one gets $\mu_p(o_{1:t}\givena a_{1:t})=\mu_{p^\star}(o_{1:t}\givena a_{1:t})$, hence $p\in R_{h_t}(p^\star)$. The outer inclusion is immediate.
\end{proof}

\begin{proposition}[Refinement chain]
\label{prop:refinement-chain}
Under the hypotheses of Lemma~\ref{lem:nesting},
\[
\omega_{h_t}\le\omega_\pi\le\omega_{\mathrm{behav}}\le\omega_{\mathrm{syn}}.
\]
\end{proposition}

\begin{proof}
The factorization $p\mapsto\mu_p$ witnesses $\omega_{\mathrm{behav}}\le\omega_{\mathrm{syn}}$. The history-factorization identity of Lemma~\ref{lem:nesting} exhibits $\Prob_{p,\pi}$ as a fixed measurable function of $\mu_p$, giving $\omega_\pi\le\omega_{\mathrm{behav}}$. Positivity of the policy factor along $h_t$ lets one recover $\mu_p(o_{1:t}\givena a_{1:t})=\Prob_{p,\pi}(H_t=h_t)/\prod_i\pi_i(a_i\mid h_{i-1})$, giving $\omega_{h_t}\le\omega_\pi$.
\end{proof}

\subsection{Observable quotient: why $[p^\star]$ is the target}

\begin{lemma}[Interventional semantic classes are the observable quotient]
\label{lem:observable-quotient}
If $p\equiv p'$, then for every policy $\pi$ and every horizon $T$, the induced law of $H_T$ under $(p,\pi)$ equals the induced law of $H_T$ under $(p',\pi)$.
\end{lemma}

\begin{proof}
Fix $h_T$. If $\Prob_{p,\pi}(H_T=h_T)>0$ or $\Prob_{p',\pi}(H_T=h_T)>0$, then every prefix $(h_{t-1},a_t)$ visited along $h_T$ is reachable under $p$ or $p'$, so semantic equivalence gives $\mu_p(o_t\mid h_{t-1},a_t)=\mu_{p'}(o_t\mid h_{t-1},a_t)$ for each $t$. Using the factorization
\[
\Prob_{p,\pi}(H_T=h_T)
=
\prod_{t=1}^{T}\pi_t(a_t\mid h_{t-1})\,\mu_p(o_t\mid h_{t-1},a_t),
\]
and the analogous identity for $p'$, the two probabilities are equal. If both probabilities are zero, equality is trivial. Hence the laws coincide.
\end{proof}

\begin{definition}[History-recoverable target]
\label{def:history-recoverable}
Let $(\mathcal{T},d)$ be a metric space equipped with its Borel $\sigma$-algebra, and let $\tau:\Pcal\to\mathcal{T}$ be a target map. The target $\tau$ is \emph{history-recoverable} if there exist a history-based policy $\pi$ and measurable estimators $\hat\tau_t:\Hcal_t\to\mathcal{T}$ such that for every $p\in\Pcal$,
\[
\hat\tau_t(H_t)\longrightarrow \tau(p)
\qquad
\text{in probability under the law generated by } (p,\pi).
\]
\end{definition}

\begin{theorem}[Every history-recoverable target factors through the semantic quotient]
\label{thm:factor-through-quotient}
Let $(\mathcal{T},d)$ be a metric space and let $\tau:\Pcal\to\mathcal{T}$ be history-recoverable. Then $p\equiv p'\Rightarrow\tau(p)=\tau(p')$. Equivalently, there exists a unique map $\bar\tau:\Ccal\to\mathcal{T}$ such that $\tau=\bar\tau\circ[\cdot]$.
\end{theorem}

\begin{proof}
Fix $p\equiv p'$. Let $\pi$ and $\hat\tau_t$ witness history-recoverability, and set $Y_t:=\hat\tau_t(H_t)$. By Lemma~\ref{lem:observable-quotient}, the law of $H_t$ under $(p,\pi)$ equals the law of $H_t$ under $(p',\pi)$ for every $t$. Since $Y_t$ is a measurable function of $H_t$, the law of $Y_t$ is the same under $(p,\pi)$ and $(p',\pi)$. Definition~\ref{def:history-recoverable} gives $Y_t\to\tau(p)$ in probability under $(p,\pi)$ and $Y_t\to\tau(p')$ in probability under $(p',\pi)$. Weak limits in a metric space are unique, so $\delta_{\tau(p)}=\delta_{\tau(p')}$ and $\tau(p)=\tau(p')$. Define $\bar\tau([p]):=\tau(p)$.
\end{proof}

\begin{remark}
Theorem~\ref{thm:factor-through-quotient} is the universality of the semantic quotient expressed in the adversarial reading of Remark~\ref{rem:four-observers}: no history-based estimator can resolve distinctions finer than the full-policy observer detects. Concentration on a target finer than $[p^\star]$ --- in particular, on $\{p^\star\}$ --- asks the observer to see through what the adversary can hide by substituting a semantically equivalent twin.
\end{remark}

\subsection{Strict separation of the hierarchy}

\begin{lemma}[{Fit gap: $R_\pi(p^\star)\subsetneq R_{h_t}(p^\star)$}]
\label{lem:fit-gap}
Assume $|\Ocal|\ge 2$, fix any target program $p^\star\in\Pcal$, any policy $\pi$, and any history $h_t$ with $\Prob_{p^\star,\pi}(H_t=h_t)>0$ and any next action $a$ with $\pi_{t+1}(a\mid h_t)>0$ such that $\mu_{p^\star}(\cdot\mid h_t,a)$ is not a point mass at a single observation in $\Ocal$. Then there exists $p\in R_{h_t}(p^\star)\setminus R_\pi(p^\star)$.
\end{lemma}

\begin{proof}
Let $x\in\Ocal$ be any observation with $\mu_{p^\star}(x\mid h_t,a)<1$. Construct a computable environment $p$ that agrees with $p^\star$ on $\mu(o_{1:t}\givena a_{1:t})$ for the fixed interaction prefix $h_t$ (achievable by copying $p^\star$'s history factorization on that prefix) and that emits $x$ deterministically at $(h_t,a)$, extending arbitrarily elsewhere. Then $\mu_p(o_{1:t}\givena a_{1:t})=\mu_{p^\star}(o_{1:t}\givena a_{1:t})$, so $p\in R_{h_t}(p^\star)$; and $\mu_p(\cdot\mid h_t,a)\neq \mu_{p^\star}(\cdot\mid h_t,a)$, so under $\pi$ the induced law of $H_{t+1}$ on the event $\{H_t=h_t,A_{t+1}=a\}$ differs between $p$ and $p^\star$, so $p\notin R_\pi(p^\star)$.
\end{proof}

\begin{theorem}[{Policy gap: $[p^\star]\subsetneq R_\pi(p^\star)$}]
\label{thm:policy-gap}
Assume $|\Acal|\ge 2$. Fix a target program $p^\star\in\Pcal$, and let $\pi$ be any policy. Suppose there exist a history $h_{t_0}$ with $\Prob_{p^\star,\pi}(H_{t_0}=h_{t_0})>0$, an action $b\in\Acal$ with $\pi_{t_0+1}(b\mid h_{t_0})=0$, and an observation $y\in\Ocal$ with $\mu_{p^\star}(y\mid h_{t_0},b)<1$. Then there exists $p'\in R_\pi(p^\star)\setminus [p^\star]$.
\end{theorem}

\begin{proof}
Construct a computable environment $p'$ that agrees with $p^\star$ at every conditioning pair $(h,a)\neq(h_{t_0},b)$ and that emits $y$ deterministically at $(h_{t_0},b)$. Such a $p'$ exists: if $p^\star$ has a program in $\Pcal$, adding a fixed conditional override at a single pair yields another program in $\Pcal$.

Because $\pi_{t_0+1}(b\mid h_{t_0})=0$, any history that visits the pair $(h_{t_0},b)$ has probability zero under every program in $\Pcal$ paired with $\pi$. For a history $h_T$ that does not visit $(h_{t_0},b)$, every conditioning pair along $h_T$ is unchanged, so
\[
\Prob_{p',\pi}(H_T=h_T)
=
\prod_{t=1}^T \pi_t(a_t\mid h_{t-1})\mu_{p'}(o_t\mid h_{t-1},a_t)
=
\prod_{t=1}^T \pi_t(a_t\mid h_{t-1})\mu_{p^\star}(o_t\mid h_{t-1},a_t)
=
\Prob_{p^\star,\pi}(H_T=h_T),
\]
because every factor $\mu_{p'}(o_t\mid h_{t-1},a_t)$ along such a history agrees with $\mu_{p^\star}(o_t\mid h_{t-1},a_t)$. Thus the two laws agree on every history, and $p'\in R_\pi(p^\star)$.

On the other hand, $\mu_{p'}(\cdot\mid h_{t_0},b)=\delta_y\neq \mu_{p^\star}(\cdot\mid h_{t_0},b)$, so $\mu_{p'}\neq\mu_{p^\star}$ and $p'\notin [p^\star]$.
\end{proof}

\begin{remark}
Theorem~\ref{thm:policy-gap} is the central separation in the hierarchy. The constructed $p'$ is the adversary's move: it agrees with $p^\star$ on every pair the $\pi$-fixed observer sees and disagrees only at the pair that observer never queries. Closing this gap is therefore a matter of policy, not of belief; no deterministic or support-limited policy that leaves a semantically informative reachable pair unqueried can close it.
\end{remark}

\begin{lemma}[{Syntactic gap: $\{p^\star\}\subsetneq [p^\star]$, typically}]
\label{lem:syntactic-gap}
If $\Pcal$ contains two distinct programs $p^\star, p^{\star\prime}$ with $\mu_{p^\star}=\mu_{p^{\star\prime}}$, then $\{p^\star\}\subsetneq [p^\star]$.
\end{lemma}

\begin{proof}
Immediate from the definitions.
\end{proof}

\begin{remark}
The hypothesis of Lemma~\ref{lem:syntactic-gap} is automatic in the Solomonoff setting. Any computable environment admits infinitely many syntactic encodings: append a dead code segment, reroute through a universal interpreter, or permute the program's internal representation. No nontrivial semantic class is a singleton in $\Pcal$.
\end{remark}

\begin{theorem}[Strict knowability hierarchy]
\label{thm:strict-hierarchy}
Assume $|\Ocal|\ge 2$, $|\Acal|\ge 2$, and that the hypothesis of Lemma~\ref{lem:syntactic-gap} holds. For any $p^\star\in\Pcal$ and any policy $\pi$ whose support fails at some reachable history $h_{t_0}$ and action $b$ with $\mu_{p^\star}(\cdot\mid h_{t_0},b)$ non-degenerate, and any realized history $h_t$ with $\Prob_{p^\star,\pi}(H_t=h_t)>0$ for which there exists a next action $a$ satisfying $\pi_{t+1}(a\mid h_t)>0$ and $\mu_{p^\star}(\cdot\mid h_t,a)$ non-degenerate,
\[
\{p^\star\}\subsetneq [p^\star]\subsetneq R_\pi(p^\star)\subsetneq R_{h_t}(p^\star)\subseteq \Pcal.
\]
\end{theorem}

\begin{proof}
The three strict inclusions are furnished by Lemma~\ref{lem:syntactic-gap}, Theorem~\ref{thm:policy-gap}, and Lemma~\ref{lem:fit-gap} respectively. Nesting is Lemma~\ref{lem:nesting}.
\end{proof}

\begin{remark}
Table~\ref{tab:hierarchy} records the role of each set in the paper.
\end{remark}

\begin{table}[ht]
\centering
\small
\caption{The four nested sets read as indistinguishability classes at increasing observational power, with the results that witness strict separation.}
\label{tab:hierarchy}
\begin{tabularx}{\textwidth}{@{}>{\raggedright\arraybackslash}p{0.20\textwidth}>{\raggedright\arraybackslash}p{0.34\textwidth}>{\raggedright\arraybackslash}X>{\raggedright\arraybackslash}p{0.16\textwidth}@{}}
\toprule
\textbf{Set} & \textbf{Observer that cannot detect a swap from $p^\star$} & \textbf{Strictly contains} & \textbf{Witnessed by} \\
\midrule
$\{p^\star\}$ & Reads the source code & Nothing & --- \\
$[p^\star]$ & Runs any policy and sees its induced law & $\{p^\star\}$ whenever $[p^\star]$ has $>1$ element & Lemma~\ref{lem:syntactic-gap} \\
$R_\pi(p^\star)$ & Sees the full law of histories under fixed $\pi$ & $[p^\star]$ when $\pi$ leaves some pair unqueried & Theorem~\ref{thm:policy-gap} \\
$R_{h_t}(p^\star)$ & Sees only the realized transcript $h_t$ & $R_\pi(p^\star)$ when some supported next action at $h_t$ is non-degenerate & Lemma~\ref{lem:fit-gap} \\
\bottomrule
\end{tabularx}
\end{table}

The rest of the paper is organized by which set the learner's posterior is trying to concentrate on. The belief side of Section~\ref{sec:belief} delivers concentration on $R_{h_t}(p^\star)$ in a predictive-merging sense. Sections~\ref{sec:semantic-action} and~\ref{sec:main} give the sufficient conditions under which the posterior does concentrate on $[p^\star]$. Section~\ref{sec:near-miss} then reads the algorithmic-free-energy principle as the wrong-observer boundary instance of the two-observer functional and shows that, outside the action-admissible range, the same Bayes/Gibbs architecture need not concentrate on $[p^\star]$ under a substituted finite-support prior: its posterior may stay bounded away from one.

\section{The Two-Observer Functional and the Meeting-Point Theorem}
\label{sec:functional}

The problem box of Section~\ref{sec:intro} asks for a triple --- prior, belief rule, action rule --- under which the posterior mass on $[p^\star]$ converges almost surely to one. The four-set hierarchy $\{p^\star\}\subseteq[p^\star]\subseteq R_\pi(p^\star)\subseteq R_{h_t}(p^\star)$ of Section~\ref{sec:hierarchy} identifies the target $[p^\star]$ as the indistinguishability class of the behavioral observer $\omega_{\mathrm{behav}}$ and as the finest set any history-based learner can hope to recover (Theorem~\ref{thm:factor-through-quotient}). This section packages the paper's answer to the box as a single variational centerpiece $\mathcal J_t^{\omega_B,\omega_A}$ together with its exact action-realizing kernel lift $\mathcal J_{t,\mathrm{ker}}^{\omega_B,\omega_A}$: the former isolates the belief and class-targeting geometry, while the latter absorbs the realized action into the variational object itself. The section's punchline is the \emph{meeting-point theorem} (Theorem~\ref{thm:meeting-point}): belief admissibility is a half-lattice above $\omega_{\mathrm{behav}}$ and action admissibility is a half-lattice below, so the two ranges meet at the single observer $\omega_{\mathrm{behav}}$, with the canonical quotient-free belief lift at $\omega_{\mathrm{syn}}$. The two-observer indexing is therefore not a stylistic convenience; it is forced by the refinement lattice of observers that organizes the four-set hierarchy.

The results of this section use only the observer structure of Section~\ref{sec:hierarchy} and the universal prior $w$ of Definition~\ref{def:universal-prior}; the belief-side and action-side constructions they license are identified in Sections~\ref{sec:belief} and~\ref{sec:semantic-action} respectively, and the convergence claim is settled in Section~\ref{sec:main}.

\subsection{Quotient data at an observer}
\label{sec:quotient-data}

Fix an observer $\omega$ with equivalence $\sim_\omega$ on $\Pcal$, and write $\Ccal^\omega:=\Pcal/\!\sim_\omega$ for the set of $\omega$-classes. The prior pushes forward to
\[
\bar w^\omega(c):=\sum_{p\in c}w(p),\qquad c\in\Ccal^\omega,
\]
and the Bayes posterior on $\Pcal$ marginalizes to the class posterior
\[
q_t^\star(c\mid h_t):=\sum_{p\in c}q_t^\star(p\mid h_t).
\]
When $\omega\ge\omega_{\mathrm{behav}}$, every $\omega$-class $c$ is contained in a single behavioral class, so all $p\in c$ share the same interactive law; the class likelihood $\mu_c:=\mu_p$ for any $p\in c$ is then well-defined. When $\omega<\omega_{\mathrm{behav}}$, an $\omega$-class may contain programs with distinct interactive laws, and $\mu_c$ is multi-valued at the class level.

Regardless of where $\omega$ sits relative to $\omega_{\mathrm{behav}}$, the posterior-weighted class-predictive law
\[
\mu_c^{q_t^\star}(o\mid h_t,a):=\frac{1}{q_t^\star(c\mid h_t)}\sum_{p\in c}q_t^\star(p\mid h_t)\,\mu_p(o\mid h_t,a)
\]
is well-defined for every $c$ with $q_t^\star(c\mid h_t)>0$; and so is its complement
\[
\Qcal_t^{-c,\omega}(o\mid h_t,a):=\frac{1}{1-q_t^\star(c\mid h_t)}\sum_{c'\in\Ccal^\omega\setminus\{c\}}q_t^\star(c'\mid h_t)\,\mu_{c'}^{q_t^\star}(o\mid h_t,a).
\]
When $\omega\ge\omega_{\mathrm{behav}}$, $\mu_c^{q_t^\star}=\mu_c$ and the posterior weights in the numerator cancel; we write $\mu_c$ without the posterior superscript in this regime.

\begin{definition}[Bhattacharyya separation at observer $\omega$]
\label{def:bhat-omega}
The \emph{$\omega$-class Bhattacharyya separation} of class $c\in\Ccal^\omega$ under action $a\in\Acal$ at history $h_t$ with $0<q_t^\star(c\mid h_t)<1$ is
\[
B_t^\omega(c,a;h_t):=-\log\sum_{o\in\Ocal}\sqrt{\mu_c^{q_t^\star}(o\mid h_t,a)\,\Qcal_t^{-c,\omega}(o\mid h_t,a)}\;\in\;[0,\infty].
\]
At $\omega=\omega_{\mathrm{behav}}$ we write $B_t:=B_t^{\omega_{\mathrm{behav}}}$; this is the separation used by the class-against-complement action primitive of Section~\ref{sec:semantic-action}.
\end{definition}

Finally, any policy $\pi$ whose action at step $t+1$ is generated by first sampling a class $C\in\Ccal^{\omega_A}$ and then selecting an action targeting that class induces a \emph{class-targeting distribution} $\nu_t^\pi(\cdot\mid h_t)$ on $\Ccal^{\omega_A}$. The class-against-complement action rule of Definition~\ref{def:semantic-rule} fits this template with $\omega_A=\omega_{\mathrm{behav}}$; for now it suffices to know that $\nu_t^\pi$ is well-defined for the policies Section~\ref{sec:semantic-action} will construct.

\subsection{The two-observer functional}
\label{sec:two-observer-functional}

Fix positive scalars $\beta,\gamma>0$ and two observers $\omega_B,\omega_A$.

\begin{definition}[Raw policy-coupled scaffold]
\label{def:raw-two-observer-functional}
For a distribution $q$ on $\Pcal$ and a policy $\pi$, define
\[
\mathcal J_{t,\mathrm{raw}}^{\omega_B,\omega_A}[q,\pi;h_t]
\;:=\;
\underbrace{\E_q[\ellp_t(p;h_t)]+\KL(q\givena w)}_{\mathcal F_t[q;h_t]}
\;-\;\beta\,\mathcal A_t^{\omega_A}[\pi;h_t,q]
\;+\;\gamma\,\KL\!\bigl(\nu_t^\pi\,\|\,\bar w^{\omega_A}\bigr),
\]
where $\ellp_t(p;h_t):=-\log\mu_p(o_{1:t}\mid a_{1:t})$ is the history codelength, $\mathcal F_t$ is the algorithmic free energy referenced by the belief observer $\omega_B$ through its pushforward $\bar q^{\omega_B}$ (Section~\ref{sec:belief}), and
\[
\mathcal A_t^{\omega_A}[\pi;h_t,q]:=\E_{C\sim\nu_t^\pi}\bigl[B_t^{\omega_A}(C,A_\pi;h_t)\bigr]
\]
is the expected $\omega_A$-class Bhattacharyya separation under $\pi$'s class-targeting distribution on $\Ccal^{\omega_A}$.
\end{definition}

A parallel \emph{mutual-information variant} replaces the action term:
\[
\mathcal J_{t,\mathrm{raw}}^{\omega_B,\omega_A,\mathrm{MI}}[q,\pi;h_t]
:=\mathcal F_t[q;h_t]-\beta\,\I_q\!\bigl(C^{\omega_A};O\mid A_\pi,h_t\bigr),
\]
where $C^{\omega_A}$ is the random variable coding the $\omega_A$-class of the true program under $q$. At $\omega_A=\omega_{\mathrm{syn}}$, $C^{\omega_A}=P$ and this recovers the standard expected-free-energy information term. The two variants differ by a swap of one R\'enyi divergence for another: the Bhattacharyya action term is $\tfrac{1}{2}D_{1/2}$, the MI term is $D_1=\KL$ under the conventional active-inference decomposition.

\begin{definition}[Two-observer variational functional]
\label{def:two-observer-functional}
For $c\in\Ccal^{\omega_A}$ and history $h_t$, set
\[
\bar B_t^{\omega_A}(c,a;h_t):=\min\!\bigl\{1,B_t^{\omega_A}(c,a;h_t)\bigr\},
\qquad
G_t^{\omega_A}(c;h_t):=\sup_{a\in\Acal}\bar B_t^{\omega_A}(c,a;h_t)\in[0,1].
\]
For a distribution $\nu$ on $\Ccal^{\omega_A}$, define the \emph{two-observer functional}
\[
\boxed{\;
\mathcal J_t^{\omega_B,\omega_A}[q,\nu;h_t]
:=\;
\mathcal F_t[q;h_t]
-\beta\,\E_{C\sim\nu}\!\bigl[G_t^{\omega_A}(C;h_t)\bigr]
+\gamma\,\KL\!\bigl(\nu\,\|\,\bar w^{\omega_A}\bigr).
\;}
\]
\end{definition}

\begin{proposition}[Exact global minimizer of the two-observer functional]
\label{prop:two-observer-minimizer}
For fixed $h_t$ and observers $(\omega_B,\omega_A)$, the exact Bayes/Gibbs posterior $q_t^\star(\cdot\mid h_t)$ uniquely minimizes $\mathcal J_t^{\omega_B,\omega_A}[q,\nu;h_t]$ in $q$. For this $q_t^\star$ fixed, the unique minimizer in $\nu$ is
\[
\nu_t^{\mathrm G,\omega_A}(c\mid h_t)
:=
\frac{\bar w^{\omega_A}(c)\exp\!\bigl((\beta/\gamma)G_t^{\omega_A}(c;h_t)\bigr)}
{\sum_{d\in\Ccal^{\omega_A}}\bar w^{\omega_A}(d)\exp\!\bigl((\beta/\gamma)G_t^{\omega_A}(d;h_t)\bigr)}.
\]
Hence $\bigl(q_t^\star,\nu_t^{\mathrm G,\omega_A}\bigr)$ is the exact global minimizer of $\mathcal J_t^{\omega_B,\omega_A}[q,\nu;h_t]$.
\end{proposition}

\begin{proof}
The $q$-dependence of $\mathcal J_t^{\omega_B,\omega_A}$ is exactly the belief term $\mathcal F_t[q;h_t]$, whose unique minimizer is the Bayes/Gibbs posterior by Lemma~\ref{lem:variational}. For the $\nu$-part, abbreviate $G(c):=G_t^{\omega_A}(c;h_t)$ and define
\[
Z_G:=\sum_{d\in\Ccal^{\omega_A}}\bar w^{\omega_A}(d)e^{(\beta/\gamma)G(d)}.
\]
Because $G(d)\in[0,1]$ and $\sum_d\bar w^{\omega_A}(d)=1$, one has $1\le Z_G\le e^{\beta/\gamma}<\infty$. Let
\[
\nu_G(c):=\frac{\bar w^{\omega_A}(c)e^{(\beta/\gamma)G(c)}}{Z_G}.
\]
Then, for every $\nu$,
\begin{align*}
\gamma\,\KL(\nu\|\bar w^{\omega_A})-\beta\,\E_\nu[G]
&=\gamma\sum_c \nu(c)\log\frac{\nu(c)}{\bar w^{\omega_A}(c)e^{(\beta/\gamma)G(c)}}\\
&=\gamma\,\KL(\nu\|\nu_G)-\gamma\log Z_G.
\end{align*}
The right-hand side is uniquely minimized at $\nu=\nu_G$. Combining the two one-sided minimizations gives the exact global minimizer pair.
\end{proof}

\begin{remark}[Why the raw scaffold is auxiliary]
\label{rem:raw-scaffold-auxiliary}
The raw policy functional $\mathcal J_{t,\mathrm{raw}}^{\omega_B,\omega_A}[q,\pi;h_t]$ is still the geometric scaffold of the paper, but its action term is coupled to both $q$ and the realized policy $\pi$, so an exact joint minimizer need not coincide with the canonical Bayes/Gibbs-plus-separation pair. Definition~\ref{def:two-observer-functional} separates the belief side from the class-targeting action side while preserving the same observer-indexed separation geometry. Proposition~\ref{prop:two-observer-minimizer} is the exact global-minimizer statement for $\mathcal J_t^{\omega_B,\omega_A}$.
\end{remark}

\begin{definition}[Kernel lift of the two-observer functional]
\label{def:kernel-functional}
Fix a full-support reference distribution $\lambda\in\Delta(\Acal)$. For a distribution $\kappa$ on $\Ccal^{\omega_A}\times\Acal$, define
\[
\mathcal J_{t,\mathrm{ker}}^{\omega_B,\omega_A}[q,\kappa;h_t]
:=
\mathcal F_t[q;h_t]
-\beta\,\E_{(C,A)\sim\kappa}\!\bigl[\bar B_t^{\omega_A}(C,A;h_t)\bigr]
+\gamma\,\KL\!\bigl(\kappa\,\|\,\bar w^{\omega_A}\otimes\lambda\bigr).
\]
\end{definition}

\begin{proposition}[Exact global minimizer of the kernel lift]
\label{prop:kernel-functional-minimizer}
For fixed $h_t$, observers $(\omega_B,\omega_A)$, and reference distribution $\lambda$, the exact Bayes/Gibbs posterior $q_t^\star(\cdot\mid h_t)$ uniquely minimizes $\mathcal J_{t,\mathrm{ker}}^{\omega_B,\omega_A}[q,\kappa;h_t]$ in $q$. For this $q_t^\star$ fixed, the unique minimizer in $\kappa$ is
\[
\kappa_t^{\mathrm G,\omega_A}(c,a\mid h_t)
:=
\frac{\bar w^{\omega_A}(c)\lambda(a)\exp\!\bigl((\beta/\gamma)\bar B_t^{\omega_A}(c,a;h_t)\bigr)}
{\sum_{d\in\Ccal^{\omega_A}}\sum_{b\in\Acal}\bar w^{\omega_A}(d)\lambda(b)\exp\!\bigl((\beta/\gamma)\bar B_t^{\omega_A}(d,b;h_t)\bigr)}.
\]
Hence $\bigl(q_t^\star,\kappa_t^{\mathrm G,\omega_A}\bigr)$ is the exact global minimizer of $\mathcal J_{t,\mathrm{ker}}^{\omega_B,\omega_A}[q,\kappa;h_t]$.
\end{proposition}

\begin{proof}
The $q$-dependence is again exactly $\mathcal F_t[q;h_t]$, whose unique minimizer is $q_t^\star$ by Lemma~\ref{lem:variational}. For the $\kappa$-part, abbreviate $\bar B(c,a):=\bar B_t^{\omega_A}(c,a;h_t)$ and define
\[
Z_{\mathrm{ker}}
:=
\sum_{d\in\Ccal^{\omega_A}}\sum_{b\in\Acal}\bar w^{\omega_A}(d)\lambda(b)e^{(\beta/\gamma)\bar B(d,b)}.
\]
Since $\bar B(d,b)\in[0,1]$ and $\bar w^{\omega_A}\otimes\lambda$ is a probability measure, one has $1\le Z_{\mathrm{ker}}\le e^{\beta/\gamma}<\infty$. Let
\[
\kappa_G(c,a):=
\frac{\bar w^{\omega_A}(c)\lambda(a)e^{(\beta/\gamma)\bar B(c,a)}}{Z_{\mathrm{ker}}}.
\]
Then, for every $\kappa$,
\begin{align*}
\gamma\,\KL\!\bigl(\kappa\,\|\,\bar w^{\omega_A}\otimes\lambda\bigr)-\beta\,\E_\kappa[\bar B]
&=\gamma\sum_{c,a}\kappa(c,a)\log\frac{\kappa(c,a)}{\bar w^{\omega_A}(c)\lambda(a)e^{(\beta/\gamma)\bar B(c,a)}}\\
&=\gamma\,\KL(\kappa\|\kappa_G)-\gamma\log Z_{\mathrm{ker}}.
\end{align*}
The right-hand side is uniquely minimized at $\kappa=\kappa_G$. Combining the two one-sided minimizations gives the exact global minimizer pair.
\end{proof}

% NOTE [compact-kernel exactness target]
% This proposition is tracked by compact_kernel_exactness_lock.md. The matching Lean declaration
% now uses the measure-theoretic compact-kernel setup; downstream compact support and separation
% corollaries remain under the compact exactness lock until they are rethreaded through it.
\begin{proposition}[Exact global minimizer of the compact-action kernel lift]
\label{prop:kernel-functional-minimizer-compact}
If one relaxes the standing finite-action setup and instead lets $(\Acal,d_\Acal)$ be a compact metric space equipped with a full-support Borel probability measure $\lambda$, read Definition~\ref{def:kernel-functional} measure-theoretically: $\kappa$ ranges over Borel probability measures on $\Ccal^{\omega_A}\times\Acal$, the expectation term is the corresponding integral, and the KL term is relative to the product measure $\bar w^{\omega_A}\otimes\lambda$. Assume further that for every fixed history $h_t$ and observer class $c\in\Ccal^{\omega_A}$, the map
\[
a\longmapsto \bar B_t^{\omega_A}(c,a;h_t)
\]
is Borel measurable. Then for fixed $h_t$, observers $(\omega_B,\omega_A)$, and reference measure $\lambda$, the exact Bayes/Gibbs posterior $q_t^\star(\cdot\mid h_t)$ uniquely minimizes $\mathcal J_{t,\mathrm{ker}}^{\omega_B,\omega_A}[q,\kappa;h_t]$ in $q$. For this $q_t^\star$ fixed, the unique minimizer in $\kappa$ is the Gibbs kernel $\kappa_{t,\mathrm{cts}}^{\mathrm G,\omega_A}(\cdot,\cdot\mid h_t)$ absolutely continuous with respect to $\bar w^{\omega_A}\otimes\lambda$ and density
\[
\frac{d\kappa_{t,\mathrm{cts}}^{\mathrm G,\omega_A}}{d(\bar w^{\omega_A}\otimes\lambda)}(c,a\mid h_t)
:=
\frac{\exp\!\bigl((\beta/\gamma)\bar B_t^{\omega_A}(c,a;h_t)\bigr)}
{Z_{\mathrm{ker}}^{\mathrm{cts}}(h_t)},
\]
where
\[
Z_{\mathrm{ker}}^{\mathrm{cts}}(h_t)
:=
\sum_{d\in\Ccal^{\omega_A}}\bar w^{\omega_A}(d)\int_{\Acal}
\exp\!\bigl((\beta/\gamma)\bar B_t^{\omega_A}(d,a;h_t)\bigr)\,d\lambda(a).
\]
Hence $\bigl(q_t^\star,\kappa_{t,\mathrm{cts}}^{\mathrm G,\omega_A}\bigr)$ is the exact global minimizer of the measure-theoretic kernel lift.
\end{proposition}

\begin{proof}
The $q$-dependence is still exactly $\mathcal F_t[q;h_t]$, so Lemma~\ref{lem:variational} gives the unique minimizer $q_t^\star$. For the $\kappa$-part, abbreviate $\bar B(c,a):=\bar B_t^{\omega_A}(c,a;h_t)$ and $\mu_0:=\bar w^{\omega_A}\otimes\lambda$. Because $0\le \bar B(c,a)\le 1$, the normalizing constant satisfies
\[
1\le Z_{\mathrm{ker}}^{\mathrm{cts}}(h_t)=\int e^{(\beta/\gamma)\bar B}\,d\mu_0\le e^{\beta/\gamma}<\infty.
\]
Let $\kappa_G:=\kappa_{t,\mathrm{cts}}^{\mathrm G,\omega_A}(\cdot,\cdot\mid h_t)$. Then
\[
\frac{d\kappa_G}{d\mu_0}(c,a)=\frac{e^{(\beta/\gamma)\bar B(c,a)}}{Z_{\mathrm{ker}}^{\mathrm{cts}}(h_t)}.
\]
If $\kappa$ is not absolutely continuous with respect to $\mu_0$, then $\KL(\kappa\|\mu_0)=+\infty$, so it cannot minimize the action term. For $\kappa\ll\mu_0$, the Radon--Nikodym identity gives
\begin{align*}
\gamma\,\KL(\kappa\|\mu_0)-\beta\int \bar B\,d\kappa
&=\gamma\int \log\!\left(\frac{d\kappa/d\mu_0}{e^{(\beta/\gamma)\bar B}}\right)\,d\kappa\\
&=\gamma\int \log\!\left(\frac{d\kappa/d\mu_0}{d\kappa_G/d\mu_0}\right)\,d\kappa-\gamma\log Z_{\mathrm{ker}}^{\mathrm{cts}}(h_t)\\
&=\gamma\,\KL(\kappa\|\kappa_G)-\gamma\log Z_{\mathrm{ker}}^{\mathrm{cts}}(h_t).
\end{align*}
The right-hand side is uniquely minimized at $\kappa=\kappa_G$. Combining the belief-side and action-side minimizations yields the exact global minimizer pair.
\end{proof}

\begin{remark}[What the kernel lift changes]
\label{rem:kernel-functional-scope}
Definition~\ref{def:kernel-functional} absorbs the selector into the variational object itself: action is sampled directly from the exact minimizing class-action kernel $\kappa_t^{\mathrm G,\omega_A}$. Its price is that convergence now depends on the reference mass that $\lambda$ assigns to separating actions. In a finite action alphabet this is automatic for uniform $\lambda$; in a countably infinite action alphabet it is an extra hypothesis, not a free consequence of $\sup_a B_t^{\omega_A}(c,a;h_t)>0$.
\par
Proposition~\ref{prop:kernel-functional-minimizer-compact} gives a second exact route: for compact metric action spaces, the same Gibbs-kernel minimizer exists measure-theoretically, and a continuity argument can recover a positive $\lambda$-mass floor on separating neighborhoods even without finiteness. Section~\ref{sec:main} records this as a compact-action extension of the main kernel theorem.
\end{remark}

\begin{definition}[Meeting-point specialization]
\label{def:meeting-point-shorthand}
Write
\[
\mathcal J_t^{\omega_{\mathrm{behav}}}:=\mathcal J_t^{\omega_{\mathrm{syn}},\omega_{\mathrm{behav}}},
\qquad
\mathcal J_{t,\mathrm{raw}}^{\omega_{\mathrm{behav}}}:=\mathcal J_{t,\mathrm{raw}}^{\omega_{\mathrm{syn}},\omega_{\mathrm{behav}}},
\qquad
\mathcal J_{t,\mathrm{ker}}^{\omega_{\mathrm{behav}}}:=\mathcal J_{t,\mathrm{ker}}^{\omega_{\mathrm{syn}},\omega_{\mathrm{behav}}}
\]
for the specializations with the canonical belief observer $\omega_B=\omega_{\mathrm{syn}}$ (the quotient-free top of the refinement lattice) and action observer $\omega_A=\omega_{\mathrm{behav}}$ (the meeting point identified in Theorem~\ref{thm:meeting-point}).
\end{definition}

The subsections that follow establish the two-sided admissibility structure of $\mathcal J_t^{\omega_B,\omega_A}$: the belief observer $\omega_B$ must lie above $\omega_{\mathrm{behav}}$ in the refinement lattice, and the action observer $\omega_A$ must lie at or below. These two conditions pick out $\omega_{\mathrm{behav}}$ as the unique single-observer admissible choice and single out the pair $(\omega_{\mathrm{syn}},\omega_{\mathrm{behav}})$ as the canonical two-observer choice.

\subsection{Belief admissibility: a half-lattice above $\omega_{\mathrm{behav}}$}
\label{sec:belief-admissibility}

Let $q_t^{\omega_B}(\cdot\mid h_t)$ denote Bayes applied at the level of $\omega_B$-classes, using the pushforward prior $\bar w^{\omega_B}$ and the class likelihood $\mu_c$ (when well-defined). Two observations bound where $\omega_B$ can sit.

\begin{proposition}[Belief-observer invariance above $\omega_{\mathrm{behav}}$]
\label{prop:belief-invariance-above}
Let $\omega_B$ be any observer with $\omega_B\ge\omega_{\mathrm{behav}}$. Then for every behavioral class $c_{\mathrm{behav}}\in\Pcal/\!\sim_{\mathrm{behav}}$ and every history $h_t$,
\[
\sum_{\substack{c\in\Ccal^{\omega_B}\\ c\subseteq c_{\mathrm{behav}}}}q_t^{\omega_B}(c\mid h_t)\;=\;q_t^{\omega_{\mathrm{behav}}}(c_{\mathrm{behav}}\mid h_t).
\]
In particular, $\omega_B$-level Bayes and $\omega_{\mathrm{behav}}$-level Bayes produce identical $\omega_{\mathrm{behav}}$-class posterior dynamics.
\end{proposition}

\begin{proof}
Because $\omega_B\ge\omega_{\mathrm{behav}}$, every $\omega_B$-class $c$ lies inside a single behavioral class, so $\mu_c=\mu_{c_{\mathrm{behav}}}$ for each $c\subseteq c_{\mathrm{behav}}$. Direct computation:
\begin{align*}
\sum_{c\subseteq c_{\mathrm{behav}}}q_t^{\omega_B}(c\mid h_t)
&=\sum_{c\subseteq c_{\mathrm{behav}}}\frac{\bar w^{\omega_B}(c)\,\mu_c(o_{1:t}\mid a_{1:t})}{\barxi(o_{1:t}\mid a_{1:t})}\\
&=\frac{\mu_{c_{\mathrm{behav}}}(o_{1:t}\mid a_{1:t})}{\barxi(o_{1:t}\mid a_{1:t})}\sum_{c\subseteq c_{\mathrm{behav}}}\bar w^{\omega_B}(c)
=\frac{\bar w^{\omega_{\mathrm{behav}}}(c_{\mathrm{behav}})\,\mu_{c_{\mathrm{behav}}}(o_{1:t}\mid a_{1:t})}{\barxi(o_{1:t}\mid a_{1:t})},
\end{align*}
which equals $q_t^{\omega_{\mathrm{behav}}}(c_{\mathrm{behav}}\mid h_t)$. The $\omega_B$-level and program-level mixtures coincide with $\barxi$ because pushforward preserves the mixture identity.
\end{proof}

\begin{proposition}[Belief-observer ill-typing below $\omega_{\mathrm{behav}}$]
\label{prop:belief-illtyped-below}
If $\omega_B<\omega_{\mathrm{behav}}$, there exists an $\omega_B$-class $c\in\Ccal^{\omega_B}$ containing two programs $p_1,p_2\in c$ with $\mu_{p_1}\neq\mu_{p_2}$. Consequently the class likelihood $\mu_c$ is not well-defined as a single distribution, and $\omega_B$-level Bayes is ill-typed without auxiliary within-class sub-posterior structure.
\end{proposition}

\begin{proof}
$\omega_B<\omega_{\mathrm{behav}}$ means $\omega_{\mathrm{behav}}\not\le\omega_B$ in the refinement preorder (Section~\ref{sec:hierarchy}), i.e., $p_1\sim_{\omega_B}p_2$ does not imply $p_1\sim_{\mathrm{behav}}p_2$. Pick witnesses $p_1,p_2$ with $p_1\sim_{\omega_B}p_2$ but $\mu_{p_1}\neq\mu_{p_2}$; they belong to the same $\omega_B$-class but induce distinct interactive laws. The assignment $c\mapsto\mu_c$ is then multi-valued on $c=[p_1]_{\omega_B}$.
\end{proof}

Proposition~\ref{prop:belief-invariance-above} says: above the behavioral observer, \emph{any} belief observer produces identical $\omega_{\mathrm{behav}}$-class dynamics, so the choice $\omega_B\in[\omega_{\mathrm{behav}},\omega_{\mathrm{syn}}]$ is a lift rather than a real parameter. Proposition~\ref{prop:belief-illtyped-below} says: below the behavioral observer, class-level Bayes is not a single procedure. Belief admissibility is therefore the half-lattice $\{\omega:\omega\ge\omega_{\mathrm{behav}}\}$, closed at $\omega_{\mathrm{behav}}$ and open-ended above.

\subsection{Action admissibility: a half-lattice below $\omega_{\mathrm{behav}}$}
\label{sec:action-admissibility}

The symmetric bound comes from the action side. Here the constraint is that separating two $\omega_A$-classes demands classes with distinct laws; above $\omega_{\mathrm{behav}}$, the lattice delivers pairs of classes that are behaviorally indistinguishable.

\begin{proposition}[Action-observer cap]
\label{prop:action-cap}
Let $\omega_A$ be an observer with $\omega_A>\omega_{\mathrm{behav}}$, and let $c_1,c_2\in\Ccal^{\omega_A}$ be two distinct $\omega_A$-classes contained in a single behavioral class, with $q_t^\star(c_1\mid h_t),q_t^\star(c_2\mid h_t)>0$. Set $\alpha:=q_t^\star(c_2\mid h_t)/(1-q_t^\star(c_1\mid h_t))\in(0,1]$. Then for every action $a\in\Acal$,
\[
B_t^{\omega_A}(c_1,a;h_t)\;\le\;\tfrac{1}{2}\log(1/\alpha)
\;=\;\tfrac{1}{2}\log\!\frac{1-q_t^\star(c_1\mid h_t)}{q_t^\star(c_2\mid h_t)},
\]
a bound independent of $a$.
\end{proposition}

\begin{proof}
Since $c_1,c_2$ share a behavioral class, $\mu_{c_1}(o\mid h_t,a)=\mu_{c_2}(o\mid h_t,a)=:\mu(o)$ for every $o$ and $a$. Writing $\Rcal$ for the posterior mixture over $\omega_A$-classes other than $c_1$ and $c_2$,
\[
\Qcal_t^{-c_1,\omega_A}(o\mid h_t,a)=\alpha\,\mu(o)+(1-\alpha)\,\Rcal(o).
\]
The pointwise bound $\sqrt{\alpha\mu+(1-\alpha)\Rcal}\ge\sqrt{\alpha\mu}$ gives
\[
\sum_o\sqrt{\mu_{c_1}(o)\,\Qcal_t^{-c_1,\omega_A}(o)}=\sum_o\sqrt{\mu(o)\bigl(\alpha\mu(o)+(1-\alpha)\Rcal(o)\bigr)}\ge\sqrt{\alpha}\sum_o\mu(o)=\sqrt{\alpha}.
\]
Taking $-\log$ yields $B_t^{\omega_A}(c_1,a;h_t)\le\tfrac{1}{2}\log(1/\alpha)$, independent of $a$.
\end{proof}

\begin{corollary}[Behavioral twins have frozen posterior ratio]
\label{cor:twins-frozen-ratio}
Let $\omega_A>\omega_{\mathrm{behav}}$ and let $c_1,c_2\in\Ccal^{\omega_A}$ be distinct $\omega_A$-classes contained in a single behavioral class. Then for every $t$ and every history $h_t$ with $q_t^\star(c_1\mid h_t),q_t^\star(c_2\mid h_t)>0$,
\[
\frac{q_t^\star(c_1\mid h_t)}{q_t^\star(c_2\mid h_t)}=\frac{\bar w^{\omega_A}(c_1)}{\bar w^{\omega_A}(c_2)}.
\]
In particular, no $\omega_A$-class posterior $q_t^\star(c_i\mid H_t)$ can tend to one almost surely whenever a behavioral twin has positive prior mass.
\end{corollary}

\begin{proof}
Bayes for classes with identical likelihoods $\mu_{c_1}=\mu_{c_2}$ gives
\[
q_t^\star(c_i\mid h_t)=\frac{\bar w^{\omega_A}(c_i)\,\mu_{c_i}(o_{1:t}\mid a_{1:t})}{\barxi(o_{1:t}\mid a_{1:t})},
\]
and the ratio at $c_1,c_2$ reduces to $\bar w^{\omega_A}(c_1)/\bar w^{\omega_A}(c_2)$ at every history.
\end{proof}

The cap theorem says: above the behavioral observer, no action can separate behavioral twins, so the Bhattacharyya separation geometry underlying $\mathcal J_t^{\omega_B,\omega_A}$ cannot drive their posteriors apart; Corollary~\ref{cor:twins-frozen-ratio} is the sharpened version --- the ratio is pinned to the prior ratio, uniformly in $t$. Action admissibility is therefore the symmetric half-lattice $\{\omega:\omega\le\omega_{\mathrm{behav}}\}$, closed at $\omega_{\mathrm{behav}}$ and open-ended below.

\subsection{The meeting-point theorem}
\label{sec:meeting-point}

\begin{theorem}[Meeting point]
\label{thm:meeting-point}
In the $\mathcal J_t^{\omega_B,\omega_A}$ framework, belief admissibility and action admissibility are half-lattices bounded on opposite sides of the refinement order, and they meet at a single observer:
\begin{itemize}
\item[(a)] \textbf{Belief lower bound.} $\omega_B\ge\omega_{\mathrm{behav}}$ is required for $\omega_B$-level Bayes to be well-defined as a single procedure on $\omega_B$-classes (Proposition~\ref{prop:belief-illtyped-below}).
\item[(b)] \textbf{Belief upper freedom.} For every $\omega_B\ge\omega_{\mathrm{behav}}$, $\omega_B$-level Bayes produces identical $\omega_{\mathrm{behav}}$-class posterior trajectories to $\omega_{\mathrm{behav}}$-level Bayes (Proposition~\ref{prop:belief-invariance-above}).
\item[(c)] \textbf{Action upper bound.} $\omega_A\le\omega_{\mathrm{behav}}$ is required for the Bhattacharyya separation geometry underlying the action side to separate distinct $\omega_A$-classes uniformly in the policy (Proposition~\ref{prop:action-cap} and Corollary~\ref{cor:twins-frozen-ratio}).
\item[(d)] \textbf{Action lower freedom.} For every $\omega_A\le\omega_{\mathrm{behav}}$, the corresponding $\omega_A$-indexed class-against-complement scheme supports almost-sure convergence of $q_t^\star([p^\star]_{\omega_A}\mid H_t)\to 1$ under the appropriate separation and rule conditions (Proposition~\ref{prop:goal-dialed}).
\end{itemize}
Consequently $\omega_{\mathrm{behav}}$ is the unique observer at which belief admissibility meets action admissibility: the minimum admissible belief observer equals the maximum useful action observer.
\end{theorem}

\begin{proof}
(a) and (b) are Propositions~\ref{prop:belief-illtyped-below} and~\ref{prop:belief-invariance-above}; (c) is Proposition~\ref{prop:action-cap} specialized to any $\omega_A>\omega_{\mathrm{behav}}$, which by Corollary~\ref{cor:twins-frozen-ratio} freezes the posterior ratio of any pair of behavioral twins independently of the policy; (d) is Proposition~\ref{prop:goal-dialed} below. The belief-admissibility range $\{\omega:\omega\ge\omega_{\mathrm{behav}}\}$ and the action-admissibility range $\{\omega:\omega\le\omega_{\mathrm{behav}}\}$ are the two principal half-lattices of the refinement order bounded at $\omega_{\mathrm{behav}}$; their intersection is the singleton $\{\omega_{\mathrm{behav}}\}$ and uniqueness is the lattice identity $\{\omega_{\mathrm{behav}}\}=\{\omega\ge\omega_{\mathrm{behav}}\}\cap\{\omega\le\omega_{\mathrm{behav}}\}$.
\end{proof}

\begin{corollary}[Canonical two-observer choice]
\label{cor:canonical-pair}
If the framework requires a single-observer specification, $\omega_{\mathrm{behav}}$ is forced. If the framework allows distinct belief and action observers, the canonical pair is $(\omega_{\mathrm{syn}},\omega_{\mathrm{behav}})$: $\omega_{\mathrm{syn}}$ is the quotient-free top of the refinement lattice and $\omega_{\mathrm{behav}}$ is the meeting point. By Proposition~\ref{prop:belief-invariance-above}, the family $\{\mathcal J_t^{\omega_B,\omega_{\mathrm{behav}}}:\omega_B\ge\omega_{\mathrm{behav}}\}$ produces identical $\omega_{\mathrm{behav}}$-class dynamics; $\mathcal J_t^{\omega_{\mathrm{behav}}}=\mathcal J_t^{\omega_{\mathrm{syn}},\omega_{\mathrm{behav}}}$ of Definition~\ref{def:meeting-point-shorthand} is its canonical representative.
\end{corollary}

The meeting-point theorem is the structural content of the two-observer design. The functional's two observer indices are not stylistic: the refinement lattice has a single fixed point where the belief-side admissibility range (closed above $\omega_{\mathrm{behav}}$) meets the action-side usefulness range (closed below $\omega_{\mathrm{behav}}$), and the observer indices record where each half of the learning problem must live. The object $\omega_{\mathrm{behav}}$ is a lattice-theoretic invariant of the paper's framework, not a design choice.

\begin{remark}[Leibniz identification]
\label{rem:leibniz}
Two programs $p_1,p_2$ satisfy $p_1\sim_{\mathrm{behav}}p_2$ iff they are indiscernible under every interactive probing --- the interactive-learning analogue of Leibniz's identity of indiscernibles on $\Pcal$. In these terms, $[p^\star]=[p^\star]_{\mathrm{behav}}$ is the Leibniz invariant of $p^\star$, and Theorem~\ref{thm:meeting-point} identifies this invariant as a lattice fixed point rather than a pragmatic stipulation: the two admissibility ranges already force it before any external criterion of ``behavioral equivalence'' is imposed.
\end{remark}

\subsection{Goal-dialed convergence}
\label{sec:goal-dialed}

The meeting-point theorem fixes a canonical observer; the family $\{\mathcal J_t^{\omega_B,\omega_A}:\omega_A\le\omega_{\mathrm{behav}}\}$ nevertheless has content as a \emph{dial}. Each $\omega_A\le\omega_{\mathrm{behav}}$ corresponds to a coarser learning target $[p^\star]_{\omega_A}\supseteq[p^\star]$, and the convergence argument of Section~\ref{sec:main} specializes to the dialed-down target with minimal change.

\begin{proposition}[Goal-dialed convergence]
\label{prop:goal-dialed}
Fix an observer $\omega_A$ with $\omega_A\le\omega_{\mathrm{behav}}$ and a target program $p^\star\in\Pcal$, and let $[p^\star]_{\omega_A}:=\{p\in\Pcal:p\sim_{\omega_A}p^\star\}$. Suppose:
\begin{itemize}
\item[\textup{(C1$^{\omega_A}$)}] the universal prior of Definition~\ref{def:universal-prior} is in force, so $\bar w^{\omega_A}(c)>0$ for every $c\in\Ccal^{\omega_A}$;
\item[\textup{(C2)}] belief is the Bayes/Gibbs posterior (Lemma~\ref{lem:variational});
\item[\textup{(C3$^{\omega_A}$)}] for every $c\in\Ccal^{\omega_A}$ with $0<q_t^\star(c\mid h_t)<1$, $\sup_{a\in\Acal}B_t^{\omega_A}(c,a;h_t)\ge\eta^{\omega_A}(c)>0$;
\item[\textup{(C4$^{\omega_A}$)}] action follows the class-against-complement rule of Section~\ref{sec:semantic-action} with $\Ccal$ replaced by $\Ccal^{\omega_A}$ throughout.
\end{itemize}
Then $q_t^\star([p^\star]_{\omega_A}\mid H_t)\to 1$ almost surely.
\end{proposition}

\begin{proof}[Proof sketch]
The main-theorem argument of Section~\ref{sec:main} (Theorem~\ref{thm:semantic-convergence}) uses five ingredients: positivity of the pushforward prior on each class; the separation condition (C3); a conditional Borel--Cantelli argument; a Bhattacharyya contraction lemma; and a promotion-supporting lower bound on class targeting. Ingredient~1 transfers by (C1$^{\omega_A}$), since the universal prior assigns positive mass to every computable program and each $\omega_A$-class contains at least one. Ingredient~2 is (C3$^{\omega_A}$). Ingredients~3 and~4 are observer-agnostic: Borel--Cantelli does not depend on the $\omega$ index, and the Bhattacharyya contraction uses only the R\'enyi-$1/2$ identity in the likelihood-ratio calculus. Ingredient~5 follows from Proposition~\ref{prop:semantic-is-promotion-supporting} applied with $\Ccal^{\omega_A}$ in place of $\Ccal$. The place where the choice $\omega_A=\omega_{\mathrm{behav}}$ was load-bearing in Theorem~\ref{thm:semantic-convergence} was the implicit exclusion of $\omega_A>\omega_{\mathrm{behav}}$ (Proposition~\ref{prop:action-cap}); under $\omega_A\le\omega_{\mathrm{behav}}$ no cap applies. Details of the transfer are given alongside Theorem~\ref{thm:semantic-convergence}.
\end{proof}

Proposition~\ref{prop:goal-dialed} reads the family $\{\mathcal J_t^{\omega_B,\omega_A}:\omega_A\le\omega_{\mathrm{behav}}\}$ as an observer dial: the corresponding $\omega_A$-indexed class-against-complement scheme drives the posterior onto the $\omega_A$-class of $p^\star$, with $\omega_A=\omega_{\mathrm{behav}}$ the finest allowable setting (the knowability ceiling of Theorem~\ref{thm:factor-through-quotient}) and coarser $\omega_A$'s targeting coarser properties of programs.

\subsection{How the rest of the paper reads under the functional}
\label{sec:functional-forward-plan}

The remainder of the paper identifies the canonical Bayes/Gibbs-plus-separation structure singled out by the observer analysis, realizes the exact global minimizers of $\mathcal J_t^{\omega_B,\omega_A}$ and its action-realizing lift $\mathcal J_{t,\mathrm{ker}}^{\omega_B,\omega_A}$, and then validates the convergence conclusion.

Section~\ref{sec:belief} identifies the \emph{belief-side minimizer} of $\mathcal J_t^{\omega_B,\omega_A}$ at the canonical belief observer $\omega_B=\omega_{\mathrm{syn}}$: the Bayes/Gibbs posterior under the universal prior. The belief-side story is the content of Lemma~\ref{lem:variational} and Lemma~\ref{lem:kl-necessity}, re-read as the $q$-minimization of $\mathcal F_t$.

Section~\ref{sec:semantic-action} identifies the \emph{canonical action-side realizations} at the canonical action observer $\omega_A=\omega_{\mathrm{behav}}$: the exact class-action Gibbs kernel $\kappa_t^{\mathrm G,\omega_{\mathrm{behav}}}$ of Proposition~\ref{prop:kernel-functional-minimizer}, and the Gibbs class-targeting minimizer $\nu_t^{\mathrm G,\omega_{\mathrm{behav}}}$ of Proposition~\ref{prop:two-observer-minimizer} together with a canonical score-witness selector that turns class scores into actual interventions. Propositions~\ref{prop:kernel-promotion-support} and~\ref{prop:semantic-is-promotion-supporting} supply the two promotion-supporting lower bounds used in the convergence arguments. Corollaries~\ref{cor:noise-transfer} and~\ref{cor:noise-left-invertible-history-independent} transfer the separation geometry through decodable and left-invertible observation channels.

Section~\ref{sec:main} establishes the \emph{convergence payoff}: the master support theorem and its explicit finite-time quantitative companion, the selector-free kernel theorem, the canonical selector theorem, the compact-action extension, and the matched converse theorems showing sharpness of the support condition.

Section~\ref{sec:self-ref} replaces the uncomputable behavioral observer by a finite-time policy proxy and recovers the main theorem's guarantee under eventual isolation (deterministic splitting or divergent exploration).

Section~\ref{sec:near-miss} reads the algorithmic-free-energy principle as the \emph{wrong-observer} boundary specialization of the same two-observer framework: under Bayes/Gibbs belief it pairs the top-of-lattice observer $\omega_A=\omega_{\mathrm{syn}}$ with the MI action bonus, which by Proposition~\ref{prop:action-cap} lies above the meeting point and therefore cannot separate behavioral twins; the near-miss is the boundary example outside the admissible range, not a defeat of Bayesian updating.

\section{Belief at $\omega_{\mathrm{syn}}$: Universal Bayesian Belief}
\label{sec:belief}

Section~\ref{sec:functional} fixed the canonical belief observer at the top of the refinement lattice: $\omega_B=\omega_{\mathrm{syn}}$ (Corollary~\ref{cor:canonical-pair}). By Proposition~\ref{prop:belief-invariance-above}, this choice is invariant --- any $\omega_B\in[\omega_{\mathrm{behav}},\omega_{\mathrm{syn}}]$ produces the same $\omega_{\mathrm{behav}}$-class posterior dynamics --- so nothing is lost by working at the quotient-free top. What is gained is a clean setting in which to identify the belief-side minimizer of $\mathcal J_t^{\omega_{\mathrm{syn}},\omega_A}$: Bayes/Gibbs on programs, under the universal interactive prior, with its algorithmic free energy as the belief term.

This section makes that object explicit. Definition~\ref{def:universal-prior} fixes the universal prior of Solomonoff and Hutter. Definition~\ref{def:afe} identifies the \emph{algorithmic free energy} $\Fcal_t$ with the belief term of $\mathcal J_t^{\omega_{\mathrm{syn}},\omega_A}$ at the canonical belief observer; Lemma~\ref{lem:variational} shows that its unique minimizer over $q$ is the exact Bayes/Gibbs posterior; Lemma~\ref{lem:kl-necessity} shows the KL-based form is forced within the natural convex class; Lemma~\ref{lem:merging} shows that the resulting posterior predictive law merges almost surely into the history-compatible set $R_{h_t}(p^\star)$ --- the outermost of the four sets of Section~\ref{sec:hierarchy}. The belief side alone therefore fixes the posterior at the coarsest observer in the four-set hierarchy, $\omega_{h_t}$; the gap from $R_{h_t}(p^\star)$ down to $[p^\star]$ is what the action side of Sections~\ref{sec:semantic-action} and~\ref{sec:main} closes.

\subsection{Universal prior}

\begin{definition}[Universal interactive prior]
\label{def:universal-prior}
For $p\in\Pcal$, let $\ellp(p)=|p|$ be the self-delimiting code length. Define raw weights $\wt w(p)=2^{-\ellp(p)}$, normalizing constant $Z=\sum_{p\in\Pcal}\wt w(p)\le 1$, normalized prior
\[
w(p)=\frac{\wt w(p)}{Z},
\]
and universal interactive mixture
\[
\barxi(o_{1:t}\givena a_{1:t})
:=
\sum_{p\in\Pcal} w(p)\mu_p(o_{1:t}\givena a_{1:t}).
\]
\end{definition}

\begin{lemma}[Machine invariance]
\label{lem:prior-invariance}
Let $U$ and $V$ be universal prefix machines, with corresponding normalized universal interactive mixtures $\barxi_U$ and $\barxi_V$. Then there are constants $c_{U,V}, c_{V,U}>0$ such that
\[
\barxi_U(o_{1:t}\givena a_{1:t})\ge c_{U,V}\barxi_V(o_{1:t}\givena a_{1:t}),
\qquad
\barxi_V(o_{1:t}\givena a_{1:t})\ge c_{V,U}\barxi_U(o_{1:t}\givena a_{1:t})
\]
for every $t$ and every compatible interaction prefix. Consequently the semantic complexity $K_U(\nu):=\min\{|p|:\mu_p=\nu\}$ of any computable interactive law $\nu$ is invariant up to an additive constant across universal machines.
\end{lemma}

\begin{proof}
Let $Z_U:=\sum_{q\in\Pcal_U}2^{-|q|}$ and $Z_V:=\sum_{p\in\Pcal_V}2^{-|p|}$. Because $U$ is universal, there is a fixed self-delimiting translator $\tau_{V\to U}$ such that for every $p\in\Pcal_V$, the concatenated $U$-program $\tau_{V\to U}p$ lies in $\Pcal_U$ and induces the same interactive law. Hence
\[
\barxi_U(o_{1:t}\givena a_{1:t})
\ge
\frac{1}{Z_U}\sum_{p\in\Pcal_V}2^{-|\tau_{V\to U}p|}\mu_p^V(o_{1:t}\givena a_{1:t})
=
2^{-|\tau_{V\to U}|}\frac{Z_V}{Z_U}\barxi_V(o_{1:t}\givena a_{1:t}).
\]
The reverse inequality is symmetric. For semantic complexity, if $p_V$ is a shortest $V$-program realizing $\nu$, then $\tau_{V\to U}p_V$ is a $U$-program realizing $\nu$, so $K_U(\nu)\le K_V(\nu)+|\tau_{V\to U}|$. Swapping $U$ and $V$ gives $K_V(\nu)\le K_U(\nu)+|\tau_{U\to V}|$, hence additive invariance.
\end{proof}

\begin{lemma}[{Positive prior mass on $[p^\star]$ is necessary}]
\label{lem:prior-necessity}
Let $w'$ be any prior on $\Pcal$, with induced mixture $\xi'(o_{1:t}\givena a_{1:t}):=\sum_p w'(p)\mu_p(o_{1:t}\givena a_{1:t})$. If $\xi'$ fails to dominate some computable interactive law $\nu$, and $c_\nu\in\Ccal$ is its semantic class, then $w'(c_\nu)=0$, and along any history generated by $\nu$ the posterior under $w'$ satisfies $q_t^{w'}(c_\nu\mid h_t)=0$ whenever defined.
\end{lemma}

\begin{proof}
If $w'(c_\nu)>0$, there exists $p\in c_\nu$ with $w'(p)>0$, and $\xi'(o_{1:t}\givena a_{1:t})\ge w'(p)\nu(o_{1:t}\givena a_{1:t})$, so $\xi'$ dominates $\nu$. Contradiction. Posterior mass on $c_\nu$ is a numerator sum over a measure-zero set, so it is zero.
\end{proof}

\subsection{Algorithmic free energy as the belief term of $\mathcal J_t$ at $\omega_B=\omega_{\mathrm{syn}}$}

For $p\in\Pcal$ and history $h_t$, define the cumulative history-fit loss
\[
L_t(p;h_t):=-\log\mu_p(o_{1:t}\givena a_{1:t}),
\]
with the convention $L_t=+\infty$ if the likelihood vanishes. Consistent with the notation of Definition~\ref{def:two-observer-functional}, $L_t(p;h_t)=\ellp_t(p;h_t)$.

\begin{definition}[Algorithmic free energy]
\label{def:afe}
For a posterior $q\in\Delta(\Pcal)$ with $\E_q[L_t]<\infty$ and $\KL(q\|w)<\infty$, the \emph{algorithmic free energy} at history $h_t$ is
\[
\Fcal_t[q;h_t]:=\E_q[L_t(p;h_t)]+\KL(q\|w).
\]
By construction, $\Fcal_t[q;h_t]$ is the belief term of the two-observer functional $\mathcal J_t^{\omega_{\mathrm{syn}},\omega_A}[q,\nu;h_t]$ at the canonical belief observer $\omega_B=\omega_{\mathrm{syn}}$ (Definition~\ref{def:two-observer-functional}): it is the part that depends on $q$ alone and not on $\nu$.
\end{definition}

\begin{lemma}[Variational identity]
\label{lem:variational}
Assume $\barxi(o_{1:t}\givena a_{1:t})>0$, and define the Bayes/Gibbs posterior
\[
q_t^\star(p\mid h_t):=\frac{w(p)\mu_p(o_{1:t}\givena a_{1:t})}{\barxi(o_{1:t}\givena a_{1:t})}.
\]
Then for every admissible posterior $q$,
\[
\Fcal_t[q;h_t]
=
\KL\!\bigl(q\,\|\,q_t^\star(\cdot\mid h_t)\bigr)-\log\barxi(o_{1:t}\givena a_{1:t}),
\]
and $\Fcal_t$ is uniquely minimized at $q_t^\star$.
\end{lemma}

\begin{proof}
By definition,
\[
\Fcal_t[q;h_t]=\sum_p q(p)\log\frac{q(p)}{w(p)\mu_p(o_{1:t}\givena a_{1:t})}
=\KL\!\bigl(q\,\|\,q_t^\star(\cdot\mid h_t)\bigr)-\log\barxi(o_{1:t}\givena a_{1:t}).
\]
Uniqueness of the minimizer follows from strict positivity of the KL divergence away from $q_t^\star$.
\end{proof}

\begin{remark}[Belief-side minimizer of $\mathcal J_t^{\omega_{\mathrm{syn}},\omega_A}$]
\label{rem:afe-bayes-unification}
Because $\Fcal_t$ is the $q$-dependent part of $\mathcal J_t^{\omega_{\mathrm{syn}},\omega_A}[q,\nu;h_t]$, Lemma~\ref{lem:variational} identifies the belief-side minimizer of the two-observer functional at $\omega_B=\omega_{\mathrm{syn}}$ with the exact Bayes/Gibbs posterior $q_t^\star$, uniformly in the action-side choice of $\nu$ and $\omega_A$. Put differently, fixing $\omega_B=\omega_{\mathrm{syn}}$ turns the $q$-minimization of $\mathcal J_t^{\omega_B,\omega_A}$ into Bayes under the universal interactive prior; no approximation is involved and no action-side choice disturbs this identification.

The algorithmic free energy of Definition~\ref{def:afe} has not, to our knowledge, appeared in this exact form. The variational free energy of active inference takes a Gibbs form but uses a stipulated parametric prior that upper-bounds model evidence only approximately \citep{Friston2010,ParrFriston2019}; Solomonoff--Hutter Bayesian prediction uses the Kolmogorov-complexity-weighted universal prior but does not state a free-energy functional \citep{Solomonoff1964a,Solomonoff1964b,Hutter2005}. Coupling the two --- the Gibbs form with the universal prior as reference measure --- collapses the distinction. Under the universal prior the algorithmic free energy is not a bound on the evidence, it is the negative log-evidence up to an additive constant, and its minimizer is the exact Bayes/Gibbs posterior. The variational structure of active inference is preserved; the stipulative quality of the prior is removed.
\end{remark}

\begin{lemma}[KL is forced by the exact Gibbs variational formula]
\label{lem:kl-necessity}
% NOTE [variational-core exactness target]
% This manuscript statement is the exact target frozen in
% `variational_core_exactness_lock.md`. The Lean declaration
% `SemanticConvergence.lem_kl_necessity` proves the countable-simplex specialization and keeps the
% posterior-gap zero criterion under a separate helper name.
Let $D(\cdot\|w)$ be a proper, convex, lower-semicontinuous functional on $\ell^1(\Pcal)$, extended by $+\infty$ off $\Delta(\Pcal)$. Assume that for every bounded loss profile $L\in\ell^\infty(\Pcal)$,
\[
\inf_{q\in\Delta(\Pcal)}\bigl\{\E_q[L]+D(q\|w)\bigr\}
=
-\log\sum_{p\in\Pcal} w(p)e^{-L(p)}.
\]
Then $D(q\|w)=\KL(q\|w)$ for every $q\in\Delta(\Pcal)$.
\end{lemma}

\begin{proof}
For bounded $f\in\ell^\infty(\Pcal)$ define $H(f):=\log\sum_p w(p)e^{f(p)}$. With $f=-L$, the hypothesis is equivalent to $H(f)=\sup_q\{\E_q[f]-D(q\|w)\}$. Fenchel--Moreau on the pair $(\ell^1,\ell^\infty)$ gives $D(q\|w)=\sup_f\{\E_q[f]-H(f)\}$.

If $\KL(q\|w)=\infty$, the upper bound $D(q\|w)\le \KL(q\|w)$ is immediate. Otherwise $q\ll w$. For $q_f(p):=w(p)e^{f(p)}/\sum_r w(r)e^{f(r)}$,
\[
\E_q[f]-H(f)=\KL(q\|w)-\KL(q\|q_f)\le \KL(q\|w),
\]
so $D(q\|w)\le\KL(q\|w)$. For the reverse, first suppose $q\not\ll w$, and let $A:=\{p\in\Pcal:q(p)>0,\ w(p)=0\}$. Then $q(A)>0$. For $f_n:=n\,\one_A$, one has $H(f_n)=\log\sum_p w(p)e^{f_n(p)}=0$ and $\E_q[f_n]=nq(A)\to\infty$, so $D(q\|w)=\infty=\KL(q\|w)$. It remains to treat the case $q\ll w$. Set $r(p):=q(p)/w(p)$ on $\{w(p)>0\}$ and $f_n(p):=\max\{-n,\min\{\log r(p),n\}\}$. Each $f_n$ is bounded. Dominated convergence gives $H(f_n)\to 0$. If $\KL(q\|w)<\infty$, $f_n(p)\to\log r(p)$ pointwise; the positive part is $q$-integrable by finiteness of $\KL$, and the bound $\sum_{r(p)<1}w(p)(-r\log r)\le 1/e$ controls the negative part, so $\E_q[f_n]\to\KL(q\|w)$. If $\KL=\infty$, monotone convergence on $f_n^+$ gives $\E_q[f_n^+]\to\infty$, and the same bound on the negative part shows $\sup_n\E_q[f_n^-]<\infty$, hence $\E_q[f_n]\to\infty$. Either way $D(q\|w)\ge\sup_n\{\E_q[f_n]-H(f_n)\}=\KL(q\|w)$.
\end{proof}

\begin{remark}
Lemma~\ref{lem:kl-necessity} singles out the KL-based free energy within the natural convex class.
\end{remark}

\subsection{Predictive merging into $R_{h_t}(p^\star)$}

\begin{lemma}[Predictive merging]
\label{lem:merging}
Fix a target program $p^\star\in\Pcal$. If the interaction is generated by $p^\star$ under any policy in the sense of Section~\ref{sec:setup}, then under the universal prior and Bayes/Gibbs posterior
\[
\sum_{t=0}^{\infty}
\KL\!\bigl(\mu_{p^\star}(\cdot\mid h_t,a_{t+1})\,\|\,\Qcal_t^O(\cdot\mid h_t,a_{t+1})\bigr)
<\infty
\qquad
\text{almost surely,}
\]
where
\[
\Qcal_t^O(o\mid h_t,a):=\sum_{p\in\Pcal} q_t^\star(p\mid h_t)\mu_p(o\mid h_t,a).
\]
In particular the one-step predictive KL tends to zero almost surely.
\end{lemma}

\begin{proof}
Mixture dominance gives $\barxi(o_{1:T}\givena a_{1:T})\ge w(p^\star)\mu_{p^\star}(o_{1:T}\givena a_{1:T})$, so $\log(\mu_{p^\star}/\barxi)\le -\log w(p^\star)$. Sequential factoring and conditional expectation under the true law yield
\[
\sum_{t=0}^{T-1}\E\!\left[\KL\!\bigl(\mu_{p^\star}(\cdot\mid h_t,a_{t+1})\,\|\,\Qcal_t^O(\cdot\mid h_t,a_{t+1})\bigr)\right]
\le -\log w(p^\star).
\]
Monotone convergence gives finiteness of the infinite sum almost surely, hence summand convergence to zero.
\end{proof}

\begin{remark}[Belief alone reaches the outermost set]
\label{rem:belief-reaches-Rht}
Predictive merging places the Bayes/Gibbs posterior, in a one-step predictive sense, inside the outermost set $R_{h_t}(p^\star)$ of the four-set hierarchy of Section~\ref{sec:hierarchy}: the posterior predictive law matches the true next-step law along the realized interaction. This is the Solomonoff-interactive analogue of the classical merging-of-opinions phenomenon \citep{BlackwellDubins1962}. In observer terms, the belief-side minimizer of $\mathcal J_t^{\omega_{\mathrm{syn}},\omega_A}$ under the universal prior reaches the history observer $\omega_{h_t}$; the two remaining strict inclusions $R_{h_t}(p^\star)\supsetneq R_\pi(p^\star)\supsetneq[p^\star]$ of Theorem~\ref{thm:strict-hierarchy} are not closed by belief alone. Closing them requires an action-side rule that lifts the realized observer past $\omega_\pi$ toward $\omega_{\mathrm{behav}}$ --- the content of Sections~\ref{sec:semantic-action} and~\ref{sec:main}.
\end{remark}


\section{Action at $\omega_{\mathrm{behav}}$: Class-Against-Complement}
\label{sec:semantic-action}

Theorem~\ref{thm:meeting-point} fixed the canonical action observer at the meeting point $\omega_A=\omega_{\mathrm{behav}}$: any observer above the meeting point caps the action bonus independent of $a$ (Proposition~\ref{prop:action-cap}), and any observer below loosens the target (Proposition~\ref{prop:goal-dialed}). What remains is to exhibit action-side realizations matched to the separation objective at $\omega_A=\omega_{\mathrm{behav}}$. This section does so at two levels. The strongest exact realization is the class-action Gibbs kernel of the kernel lift $\mathcal J_{t,\mathrm{ker}}^{\omega_{\mathrm{behav}}}$; the canonical selector realization is the class-against-complement rule attached to $\mathcal J_t^{\omega_{\mathrm{behav}}}$. The relevant comparison at $\omega_A=\omega_{\mathrm{behav}}$ is the one the boundary example of Section~\ref{sec:near-miss} will spell out the failure of: once the target is an $\omega_{\mathrm{behav}}$-class $c$, the relevant rival is not any single alternative program; it is all posterior mass outside $c$ pooled together. An action that fails to separate $c$ from that aggregate rival may be highly informative about the latent program --- as the algorithmic free energy principle's mutual-information action will be --- and still supply no evidence for convergence to $c$. The \emph{class-against-complement} primitive of Definition~\ref{def:semantic-rule} is built to drive the capped Bhattacharyya class score entering $\mathcal J_t^{\omega_{\mathrm{behav}}}$ high while maintaining a Gibbs promotion floor, while Proposition~\ref{prop:kernel-functional-minimizer} gives the exact class-action Gibbs minimizer of the kernel lift.

For a semantic class $c\in\Ccal$, the notation $\mu_c$ fixed in Section~\ref{sec:setup} denotes the common action-conditioned law induced by any program in $c$.

\begin{definition}[Class-complement predictive law]
\label{def:class-complement}
Fix a history $h_t$, a class $c\in\Ccal$, and an action $a$. If $0<q_t^\star(c\mid h_t)<1$, define the class-complement predictive law $\Qcal_t^{-c}(\cdot\mid h_t,a)\in\Delta(\Ocal)$ by
\[
\Qcal_t^{-c}(o\mid h_t,a)
:=
\frac{\sum_{d\neq c} q_t^\star(d\mid h_t)\mu_d(o\mid h_t,a)}
{1-q_t^\star(c\mid h_t)}.
\]
\end{definition}

\begin{remark}[Identification with the action-observer complement]
\label{rem:class-complement-is-omega-behav}
Definition~\ref{def:class-complement} is the specialization of the $\omega_A$-class complement $\Qcal_t^{-c,\omega_A}$ of Section~\ref{sec:quotient-data} to $\omega_A=\omega_{\mathrm{behav}}$. Above the meeting point, all behavioral twins share $\mu_p$, so the posterior-weighted superscript collapses and $\mu_c$ is the unambiguous class likelihood of Proposition~\ref{prop:belief-invariance-above}; the sum over $d\neq c$ in Definition~\ref{def:class-complement} is the same sum over $c'\in\Ccal^{\omega_{\mathrm{behav}}}\setminus\{c\}$ that appears in $\Qcal_t^{-c,\omega_{\mathrm{behav}}}$. The class score entering the action side of $\mathcal J_t^{\omega_B,\omega_{\mathrm{behav}}}$ at the canonical action observer is therefore computed exactly from the pair $(\mu_c,\Qcal_t^{-c})$ of Definition~\ref{def:class-complement}.
\end{remark}

\begin{definition}[Semantic gain]
\label{def:semantic-gain}
The \emph{semantic gain} of action $a$ for class $c$ at $h_t$ is
\[
S_t(c,a;h_t)
:=
\begin{cases}
\KL\!\bigl(\mu_c(\cdot\mid h_t,a)\,\|\,\Qcal_t^{-c}(\cdot\mid h_t,a)\bigr), & 0<q_t^\star(c\mid h_t)<1,\\
0, & q_t^\star(c\mid h_t)\in\{0,1\}.
\end{cases}
\]
\end{definition}

\begin{lemma}[Posterior-odds identity]
\label{lem:odds-identity}
Fix $c\in\Ccal$ with $0<q_t^\star(c\mid h_t)<1$ and an action $a$. Let the latent truth lie in $c$. Then
\[
S_t(c,a;h_t)
=
\E_{o\sim\mu_c(\cdot\mid h_t,a)}\!\left[
\log\frac{q_{t+1}^\star(c\mid h_t,a,o)}{1-q_{t+1}^\star(c\mid h_t,a,o)}
-\log\frac{q_t^\star(c\mid h_t)}{1-q_t^\star(c\mid h_t)}
\right].
\]
\end{lemma}

\begin{proof}
Let $\alpha_t:=q_t^\star(c\mid h_t)$. Class-level Bayes gives
\[
\frac{q_{t+1}^\star(c\mid h_t,a,o)}{1-q_{t+1}^\star(c\mid h_t,a,o)}
=
\frac{\alpha_t}{1-\alpha_t}\cdot\frac{\mu_c(o\mid h_t,a)}{\Qcal_t^{-c}(o\mid h_t,a)}.
\]
Taking logarithms and expectation under $\mu_c(\cdot\mid h_t,a)$ yields the claim.
\end{proof}

\begin{definition}[Semantic separation]
\label{def:semantic-separation}
The \emph{semantic separation} of action $a$ for class $c$ at $h_t$ is
\[
B_t(c,a;h_t)
:=
\begin{cases}
-\log\sum_{o\in\Ocal}\sqrt{\mu_c(o\mid h_t,a)\Qcal_t^{-c}(o\mid h_t,a)}, & 0<q_t^\star(c\mid h_t)<1,\\
0, & q_t^\star(c\mid h_t)\in\{0,1\}.
\end{cases}
\]
\end{definition}

\begin{remark}[Identification with the action-observer separation]
\label{rem:B-is-B-omega-behav}
$B_t(c,a;h_t)=B_t^{\omega_{\mathrm{behav}}}(c,a;h_t)$ in the sense of Definition~\ref{def:bhat-omega}: Remark~\ref{rem:class-complement-is-omega-behav} identifies the pair $(\mu_c,\Qcal_t^{-c})$ with the pair $(\mu_c^{q_t^\star},\Qcal_t^{-c,\omega_{\mathrm{behav}}})$ of Section~\ref{sec:quotient-data}, and the two definitions take $-\log$ of the same Hellinger affinity. The notational shorthand $B_t:=B_t^{\omega_{\mathrm{behav}}}$ fixed in Definition~\ref{def:bhat-omega} is in force throughout this section.
\end{remark}

\begin{lemma}[Zero criterion]
\label{lem:zero-criterion}
For $c$ and $a$ with $0<q_t^\star(c\mid h_t)<1$,
\[
S_t(c,a;h_t)=0
\iff
B_t(c,a;h_t)=0
\iff
\mu_c(\cdot\mid h_t,a)=\Qcal_t^{-c}(\cdot\mid h_t,a).
\]
\end{lemma}

\begin{proof}
KL vanishes exactly at coincidence of laws. The Hellinger affinity equals one exactly at coincidence. The zero set is preserved under $-\log$.
\end{proof}

\begin{proposition}[Chernoff--Bhattacharyya correspondence]
\label{prop:chernoff-correspondence}
Fix $c\in\Ccal$ and $a\in\Acal$ at a history $h_t$ with $0<q_t^\star(c\mid h_t)<1$. Writing $D_\alpha$ for the R\'enyi divergence of order $\alpha$:
\begin{itemize}
\item[\textup{(i)}] $B_t(c,a;h_t)=\tfrac{1}{2}D_{1/2}\!\bigl(\mu_c(\cdot\mid h_t,a)\,\|\,\Qcal_t^{-c}(\cdot\mid h_t,a)\bigr)$;
\item[\textup{(ii)}] $\KL\!\bigl(\mu_c(\cdot\mid h_t,a)\,\|\,\Qcal_t^{-c}(\cdot\mid h_t,a)\bigr)\ge 2\,B_t(c,a;h_t)$;
\item[\textup{(iii)}] for every event $E\subseteq\Ocal$, the one-step Bayes error probability from a single observation under action $a$ is bounded by
\[
\Prob\!\bigl(\text{$c$ misclassified}\mid h_t,a\bigr)\;\le\;\sqrt{q_t^\star(c\mid h_t)(1-q_t^\star(c\mid h_t))}\;e^{-B_t(c,a;h_t)}.
\]
\end{itemize}
\end{proposition}

\begin{proof}
For $\alpha\in(0,1)$,
\[
D_\alpha(P\|Q)=\tfrac{1}{\alpha-1}\log\sum_o P(o)^\alpha Q(o)^{1-\alpha},
\]
so $D_{1/2}(P\|Q)=-2\log\sum_o\sqrt{P(o)Q(o)}$, which is exactly $2B_t(c,a;h_t)$ applied to $(\mu_c,\Qcal_t^{-c})$. This is (i). For (ii), R\'enyi divergence is nondecreasing in $\alpha$ on $[0,1]$ (see \citealp[Thm.~3]{vanErvenHarremoes2014}), so $D_1\ge D_{1/2}$, and $D_1=\KL$ by definition; substituting (i) gives $\KL\ge 2B_t$. For (iii), the Bayes-optimal test between $\mu_c$ and $\Qcal_t^{-c}$ under prior weights $(q_t^\star(c\mid h_t),1-q_t^\star(c\mid h_t))$ has error probability
\[
\sum_{o\in\Ocal}\min\!\bigl\{q_t^\star(c\mid h_t)\mu_c(o),(1-q_t^\star(c\mid h_t))\Qcal_t^{-c}(o)\bigr\}
\le
\sqrt{q_t^\star(c\mid h_t)(1-q_t^\star(c\mid h_t))}\sum_o\sqrt{\mu_c(o)\Qcal_t^{-c}(o)},
\]
using $\min(x,y)\le\sqrt{xy}$ on each summand (Chernoff's bound at $s=1/2$). The inner sum equals $e^{-B_t}$ by Definition~\ref{def:semantic-separation}.
\end{proof}

The correspondence settles three book-keeping issues. First, (i) identifies $B_t$ as a R\'enyi divergence rather than an ad~hoc bound, placing the separation geometry underlying the action side of $\mathcal J_t^{\omega_B,\omega_A}$ inside the same divergence family as the KL of $\mathcal F_t$. Second, (ii) converts any KL-type separation assumption into a Bhattacharyya separation assumption at no worse than a factor of two, which is the only inequality the convergence proof will need. Third, (iii) exhibits $B_t$ as a single-step Chernoff bound on the Bayes error, so the separation condition of Definition~\ref{def:sep-condition} is a quantitative statement about how well an action distinguishes class $c$ from its class-complement on that one observation. The capped score $G_t(c;h_t)=\sup_a\min\{1,B_t(c,a;h_t)\}$ that enters $\mathcal J_t^{\omega_{\mathrm{behav}}}$ is precisely the bounded version of this same separation geometry.

\begin{definition}[Semantic action rule]
\label{def:semantic-rule}
Let $\bar w:\Ccal\to[0,1]$ be the pushforward prior from Section~\ref{sec:setup}; explicitly, $\bar w(c)=\sum_{p\in c}w(p)$. Define the capped class score
\[
\bar B_t(c,a;h_t):=\min\!\bigl\{1,B_t(c,a;h_t)\bigr\},
\qquad
G_t(c;h_t):=\sup_{a\in\Acal}\bar B_t(c,a;h_t)\in[0,1].
\]
For each class $c$ and history $h_t$, choose a measurable selector $a_t^c(h_t)$ with
\[
\tfrac{1}{2}G_t(c;h_t)
\le
\bar B_t(c,a_t^c(h_t);h_t)
\]
when $0<q_t^\star(c\mid h_t)<1$, and arbitrary otherwise. Define the class-targeting distribution
\[
\nu_t^{\mathrm G}(c\mid h_t)
:=
\frac{\bar w(c)\exp\!\bigl((\beta/\gamma)G_t(c;h_t)\bigr)}
{\sum_{d\in\Ccal}\bar w(d)\exp\!\bigl((\beta/\gamma)G_t(d;h_t)\bigr)}.
\]
The \emph{semantic action rule} samples $C_t\sim\nu_t^{\mathrm G}(\cdot\mid h_t)$ and plays $a_{t+1}=a_t^{C_t}(h_t)$.
\end{definition}

Such a selector exists because $\Acal$ is finite, so one may fix an enumeration and choose the first action satisfying the displayed inequality. By Proposition~\ref{prop:two-observer-minimizer}, $\nu_t^{\mathrm G}$ is exactly the $\omega_A=\omega_{\mathrm{behav}}$ action-side minimizer of $\mathcal J_t^{\omega_{\mathrm{behav}}}$.

\begin{definition}[Promotion-supporting action rule]
\label{def:promotion-supporting}
An action rule is \emph{promotion-supporting} if there exists $\delta>0$ such that for every class $c\in\Ccal$ and every history $h_t$ with $0<q_t^\star(c\mid h_t)<1$,
\[
\Prob\bigl(A_{t+1}=a_t^c(h_t)\mid H_t=h_t\bigr)\ge \delta\bar w(c),
\]
where $a_t^c(h_t)$ is the $c$-separating selector of Definition~\ref{def:semantic-rule}.
\end{definition}

\begin{proposition}[Semantic rule is promotion-supporting]
\label{prop:semantic-is-promotion-supporting}
The semantic action rule of Definition~\ref{def:semantic-rule} is promotion-supporting with
\[
\delta_0:=e^{-\beta/\gamma}.
\]
In particular, for every target program $p^\star\in\Pcal$, the true class $c^\star:=[p^\star]$ satisfies
\[
\Prob(C_t=c^\star\mid H_t)\ge\delta_0\bar w(c^\star)>0
\]
almost surely for every $t$.
\end{proposition}

\begin{proof}
Fix $h_t$ and $c\in\Ccal$ with $0<q_t^\star(c\mid h_t)<1$.
\[
\Prob(A_{t+1}=a_t^c(h_t)\mid H_t=h_t)\ge\Prob(C_t=c\mid H_t=h_t)=\nu_t^{\mathrm G}(c\mid h_t).
\]
Since $G_t(d;h_t)\in[0,1]$ for every $d$,
\[
\sum_{d\in\Ccal}\bar w(d)\exp\!\bigl((\beta/\gamma)G_t(d;h_t)\bigr)\le e^{\beta/\gamma}\sum_{d\in\Ccal}\bar w(d)=e^{\beta/\gamma},
\]
so
\[
\nu_t^{\mathrm G}(c\mid h_t)\ge e^{-\beta/\gamma}\bar w(c)=\delta_0\bar w(c).
\]
The universal prior assigns positive mass to every computable program, so $\bar w(c)>0$ for every $c\in\Ccal$.
\end{proof}

\begin{definition}[Kernel semantic action rule]
\label{def:kernel-semantic-rule}
Fix a full-support reference distribution $\lambda\in\Delta(\Acal)$. The \emph{kernel semantic action rule} at $\omega_A=\omega_{\mathrm{behav}}$ samples
\[
(C_t,A_{t+1})\sim \kappa_t^{\mathrm G,\omega_{\mathrm{behav}}}(\cdot,\cdot\mid h_t),
\]
where $\kappa_t^{\mathrm G,\omega_{\mathrm{behav}}}$ is the exact minimizer of Proposition~\ref{prop:kernel-functional-minimizer}. Thus action is drawn directly from the minimizing class-action kernel, with no external selector.
\end{definition}

\begin{proposition}[Kernel rule has reference-mass promotion support]
\label{prop:kernel-promotion-support}
Under the kernel semantic action rule of Definition~\ref{def:kernel-semantic-rule}, for every history $h_t$, every class $c\in\Ccal$, and every subset $S\subseteq\Acal$,
\[
\Prob\bigl(C_t=c,\ A_{t+1}\in S\,\big|\,H_t=h_t\bigr)\ge \delta_0\,\bar w(c)\,\lambda(S),
\qquad
\delta_0:=e^{-\beta/\gamma},
\]
where $\lambda(S):=\sum_{a\in S}\lambda(a)$.
\end{proposition}

\begin{proof}
By Definition~\ref{def:kernel-semantic-rule},
\[
\Prob(C_t=c,\ A_{t+1}\in S\mid H_t=h_t)
=
\sum_{a\in S}\kappa_t^{\mathrm G,\omega_{\mathrm{behav}}}(c,a\mid h_t).
\]
Since $\bar B_t(c,a;h_t)\in[0,1]$ for every pair $(c,a)$, the denominator in Proposition~\ref{prop:kernel-functional-minimizer} is at most
\[
e^{\beta/\gamma}\sum_{d\in\Ccal}\sum_{b\in\Acal}\bar w(d)\lambda(b)=e^{\beta/\gamma}.
\]
Therefore
\[
\sum_{a\in S}\kappa_t^{\mathrm G,\omega_{\mathrm{behav}}}(c,a\mid h_t)
\ge
e^{-\beta/\gamma}\bar w(c)\sum_{a\in S}\lambda(a)
=
\delta_0\,\bar w(c)\,\lambda(S).
\]
\end{proof}

\begin{proposition}[Compact-action kernel rule has reference-mass promotion support]
\label{prop:kernel-promotion-support-compact}
Assume the compact-action setting of Proposition~\ref{prop:kernel-functional-minimizer-compact} and sample
\[
(C_t,A_{t+1})\sim \kappa_{t,\mathrm{cts}}^{\mathrm G,\omega_{\mathrm{behav}}}(\cdot,\cdot\mid h_t).
\]
Then for every history $h_t$, every class $c\in\Ccal$, and every Borel set $S\subseteq\Acal$,
\[
\Prob\bigl(C_t=c,\ A_{t+1}\in S\,\big|\,H_t=h_t\bigr)\ge \delta_0\,\bar w(c)\,\lambda(S),
\qquad
\delta_0:=e^{-\beta/\gamma}.
\]
\end{proposition}

\begin{proof}
By Proposition~\ref{prop:kernel-functional-minimizer-compact},
\[
\Prob(C_t=c,\ A_{t+1}\in S\mid H_t=h_t)
=
\kappa_{t,\mathrm{cts}}^{\mathrm G,\omega_{\mathrm{behav}}}(\{c\}\times S\mid h_t)
=
\frac{\bar w(c)\int_S e^{(\beta/\gamma)\bar B_t(c,a;h_t)}\,d\lambda(a)}
{Z_{\mathrm{ker}}^{\mathrm{cts}}(h_t)}.
\]
Since $\bar B_t(c,a;h_t)\in[0,1]$ for every $(c,a)$, Proposition~\ref{prop:kernel-functional-minimizer-compact} gives
\[
Z_{\mathrm{ker}}^{\mathrm{cts}}(h_t)\le e^{\beta/\gamma}.
\]
Therefore
\[
\Prob(C_t=c,\ A_{t+1}\in S\mid H_t=h_t)
\ge
e^{-\beta/\gamma}\bar w(c)\int_S 1\,d\lambda
=
\delta_0\,\bar w(c)\,\lambda(S).
\]
\end{proof}

\begin{definition}[Decodable observation channel]
\label{def:decodable-channel}
Let $\widetilde\Ocal$ be a countable observation alphabet. A computable Markov kernel $K:\Ocal\to\Delta(\widetilde\Ocal)$ is \emph{decodable} if there exists a computable map $D:\widetilde\Ocal\to\Ocal$ such that
\[
\sum_{\widetilde o:\,D(\widetilde o)=o}K(\widetilde o\mid o)=1
\qquad\text{for every }o\in\Ocal.
\]
Equivalently, if $\widetilde O\sim K(\cdot\mid o)$, then $D(\widetilde O)=o$ almost surely. For a corrupted history
\[
\widetilde h_t=(a_1,\widetilde o_1,\dots,a_t,\widetilde o_t),
\]
write
\[
D(\widetilde h_t):=(a_1,D(\widetilde o_1),\dots,a_t,D(\widetilde o_t)).
\]
\end{definition}

\begin{proposition}[Pointwise noise-channel monotonicity]
\label{prop:noise-immunity}
Let $K:\Ocal\to\Delta(\widetilde\Ocal)$ be a Markov kernel --- a noise channel that replaces each observation $o$ with a randomized $\widetilde o\sim K(\cdot\mid o)$, independently across $t$. Write $\widetilde\mu_c(\widetilde o\mid h_t,a):=\sum_o K(\widetilde o\mid o)\mu_c(o\mid h_t,a)$ for the post-channel class likelihood and $\widetilde\Qcal_t^{-c}$ for the post-channel class-complement likelihood, and let $\widetilde B_t(c,a;h_t)$ be the Bhattacharyya separation computed from $(\widetilde\mu_c,\widetilde\Qcal_t^{-c})$. Then
\[
\widetilde B_t(c,a;h_t)\;\le\;B_t(c,a;h_t)
\qquad\text{for every $c\in\Ccal$, $a\in\Acal$, and history $h_t$.}
\]
Consequently, the automatic conclusion under channel corruption is only this pointwise data-processing inequality at a fixed latent history $h_t$, not a full transfer theorem for learning from corrupted observation histories.
\end{proposition}

\begin{proof}
By the Cauchy--Schwarz inequality, for every $\widetilde o$,
\[
\sqrt{\widetilde\mu_c(\widetilde o)\,\widetilde\Qcal_t^{-c}(\widetilde o)}
=\sqrt{\Bigl(\!\sum_o K(\widetilde o\mid o)\mu_c(o)\Bigr)\!\Bigl(\!\sum_{o'}K(\widetilde o\mid o')\Qcal_t^{-c}(o')\Bigr)}
\ge\sum_o K(\widetilde o\mid o)\sqrt{\mu_c(o)\Qcal_t^{-c}(o)},
\]
where the dependence on $h_t$ and $a$ is suppressed for readability. Summing over $\widetilde o$ and using $\sum_{\widetilde o}K(\widetilde o\mid o)=1$,
\[
\sum_{\widetilde o}\sqrt{\widetilde\mu_c(\widetilde o)\,\widetilde\Qcal_t^{-c}(\widetilde o)}\ge\sum_o\sqrt{\mu_c(o)\Qcal_t^{-c}(o)}.
\]
Taking $-\log$ reverses the inequality: $\widetilde B_t\le B_t$. No weaker positivity claim is automatic: a nonconstant channel can still collapse $\mu_c$ and $\Qcal_t^{-c}$ to the same output law. Moreover, in the interactive setting a learner that conditions on the corrupted history need not induce the pointwise posterior obtained by simply replacing $\mu_c$ with $\widetilde\mu_c$ at a fixed latent history, so a full convergence transfer would require a separate augmented-state analysis.
\end{proof}

\begin{definition}[Left-invertible finite observation channel]
\label{def:left-invertible-channel}
Assume $\Ocal$ and $\widetilde\Ocal$ are finite. For a channel $K:\Ocal\to\Delta(\widetilde\Ocal)$, write
\[
(T_Kx)(\widetilde o):=\sum_{o\in\Ocal}K(\widetilde o\mid o)x(o),
\qquad
x\in\Delta(\Ocal).
\]
The channel is \emph{left-invertible} if the linear map $T_K:\mathbb R^{\Ocal}\to\mathbb R^{\widetilde\Ocal}$ is injective, equivalently if the channel matrix has full column rank.
\end{definition}

\begin{proposition}[Uniform positive separation under left-invertible finite noise]
\label{prop:noise-left-invertible}
Assume $\Ocal$ and $\widetilde\Ocal$ are finite and let $K:\Ocal\to\Delta(\widetilde\Ocal)$ be a left-invertible channel. For every $\eta>0$ such that
\[
\mathsf S_\eta
:=
\Bigl\{
(x,y)\in\Delta(\Ocal)^2:\ 
-\log \sum_{o\in\Ocal}\sqrt{x(o)y(o)}\ge \eta
\Bigr\}
\]
is nonempty, define
\[
\widetilde\eta_K(\eta)
:=
\min_{(x,y)\in\mathsf S_\eta}
\left(
-\log \sum_{\widetilde o\in\widetilde\Ocal}\sqrt{(T_Kx)(\widetilde o)(T_Ky)(\widetilde o)}
\right).
\]
Then $\widetilde\eta_K(\eta)>0$. Consequently, whenever
\[
B_t(c,a;h_t)\ge \eta(c),
\]
one also has
\[
\widetilde B_t(c,a;h_t)\ge \widetilde\eta_K(\eta(c)).
\]
\end{proposition}

\begin{proof}
Fix $\eta>0$. The set
\[
\mathsf S_\eta
:=
\Bigl\{
(x,y)\in\Delta(\Ocal)^2:\ 
-\log \sum_{o\in\Ocal}\sqrt{x(o)y(o)}\ge \eta
\Bigr\}
\]
is compact because $\Delta(\Ocal)^2$ is compact and the Bhattacharyya separation is continuous. The map
\[
(x,y)\longmapsto \sum_{\widetilde o\in\widetilde\Ocal}\sqrt{(T_Kx)(\widetilde o)(T_Ky)(\widetilde o)}
\]
is continuous on $\mathsf S_\eta$, so its maximum on $\mathsf S_\eta$ is attained; equivalently, the minimum defining $\widetilde\eta_K(\eta)$ is attained.

To show positivity, it is enough to show that the displayed post-channel Hellinger affinity is strictly less than $1$ on $\mathsf S_\eta$. Suppose toward a contradiction that its maximum on $\mathsf S_\eta$ were $1$. Then for some $(x,y)\in\mathsf S_\eta$,
\[
\sum_{\widetilde o\in\widetilde\Ocal}\sqrt{(T_Kx)(\widetilde o)(T_Ky)(\widetilde o)}=1.
\]
For probability distributions this equality holds iff $T_Kx=T_Ky$. Since $K$ is left-invertible, $T_K$ is injective, so $x=y$. But then the clean Bhattacharyya separation is zero, contradicting $(x,y)\in\mathsf S_\eta$ with $\eta>0$. Hence $\widetilde\eta_K(\eta)>0$.

For the consequence, apply the definition with
\[
x:=\mu_c(\cdot\mid h_t,a),
\qquad
y:=\Qcal_t^{-c}(\cdot\mid h_t,a).
\]
Then
\[
B_t(c,a;h_t)\ge \eta(c)
\]
places $(x,y)$ in the nonempty set $\mathsf S_{\eta(c)}$, while
\[
\widetilde\mu_c(\cdot\mid h_t,a)=T_Kx,
\qquad
\widetilde\Qcal_t^{-c}(\cdot\mid h_t,a)=T_Ky.
\]
Therefore
\[
\widetilde B_t(c,a;h_t)\ge \widetilde\eta_K(\eta(c)).
\]
\end{proof}

\begin{proposition}[Exact preservation under decodable recodings]
\label{prop:noise-decoding}
Fix a decodable observation channel $K:\Ocal\to\Delta(\widetilde\Ocal)$ with decoder $D$, and for each program $p\in\Pcal$ define the corrupted-observation kernel
\[
\widetilde\mu_p(\widetilde o\mid \widetilde h_t,a)
:=
K(\widetilde o\mid D(\widetilde o))\,
\mu_p(D(\widetilde o)\mid D(\widetilde h_t),a).
\]
Let $\widetilde q_t(\cdot\mid \widetilde h_t)$ be the Bayes posterior on program indices computed from the family $(\widetilde\mu_p)_{p\in\Pcal}$ with the same prior weights $w(p)$, and let $\widetilde B_t(c,a;\widetilde h_t)$ be the corresponding class-against-complement Bhattacharyya separation. Then:
\begin{itemize}
\item[(i)] for every $p\in\Pcal$ and corrupted history $\widetilde h_t$,
\[
\widetilde q_t(p\mid \widetilde h_t)=q_t(p\mid D(\widetilde h_t));
\]
\item[(ii)] the behavioral and semantic equivalence relations induced by $(\widetilde\mu_p)_{p\in\Pcal}$ coincide with those induced by $(\mu_p)_{p\in\Pcal}$ on the common program index set;
\item[(iii)] for every class $c\in\Ccal$, action $a\in\Acal$, and corrupted history $\widetilde h_t$ with $0<\widetilde q_t(c\mid \widetilde h_t)<1$,
\[
\widetilde B_t(c,a;\widetilde h_t)=B_t(c,a;D(\widetilde h_t)).
\]
\end{itemize}
\end{proposition}

\begin{proof}
First note that each $\widetilde\mu_p(\cdot\mid \widetilde h_t,a)$ is a probability distribution:
\[
\sum_{\widetilde o\in\widetilde\Ocal}\widetilde\mu_p(\widetilde o\mid \widetilde h_t,a)
=
\sum_{o\in\Ocal}\mu_p(o\mid D(\widetilde h_t),a)
\sum_{\widetilde o:\,D(\widetilde o)=o}K(\widetilde o\mid o)
=1
\]
by Definition~\ref{def:decodable-channel}. Iterating the one-step definition gives, for every corrupted transcript $\widetilde o_{1:t}$,
\[
\widetilde\mu_p(\widetilde o_{1:t}\givena a_{1:t})
=
\mu_p(D(\widetilde o_{1:t})\givena a_{1:t})
\prod_{i=1}^{t}K(\widetilde o_i\mid D(\widetilde o_i)).
\]
The second factor is independent of $p$, so Bayes' rule yields
\[
\widetilde q_t(p\mid \widetilde h_t)=q_t(p\mid D(\widetilde h_t)),
\]
proving \textup{(i)}.

For \textup{(ii)}, if $\mu_{p_1}=\mu_{p_2}$ on reachable pairs, then the defining formula for $\widetilde\mu_p$ gives $\widetilde\mu_{p_1}=\widetilde\mu_{p_2}$ on reachable pairs as well. Conversely, if $\widetilde\mu_{p_1}$ and $\widetilde\mu_{p_2}$ agree on all corrupted transcripts, then summing the displayed factorization over the fiber $\{\widetilde o_{1:t}:D(\widetilde o_{1:t})=o_{1:t}\}$ gives
\[
\mu_{p_1}(o_{1:t}\givena a_{1:t})
=
\sum_{\widetilde o_{1:t}:\,D(\widetilde o_{1:t})=o_{1:t}}
\widetilde\mu_{p_1}(\widetilde o_{1:t}\givena a_{1:t})
=
\sum_{\widetilde o_{1:t}:\,D(\widetilde o_{1:t})=o_{1:t}}
\widetilde\mu_{p_2}(\widetilde o_{1:t}\givena a_{1:t})
=
\mu_{p_2}(o_{1:t}\givena a_{1:t}),
\]
where the last equality again uses Definition~\ref{def:decodable-channel}. Thus the action-conditioned chronological laws, and hence the induced behavioral and semantic partitions, coincide under the identification of the common program index set.

For \textup{(iii)}, the posterior identity from \textup{(i)} and the shared class partition from \textup{(ii)} imply
\[
\widetilde\mu_c(\widetilde o\mid \widetilde h_t,a)
=
\sum_{p\in c}\widetilde q_t(p\mid \widetilde h_t)\widetilde\mu_p(\widetilde o\mid \widetilde h_t,a)
=
K(\widetilde o\mid D(\widetilde o))\,\mu_c(D(\widetilde o)\mid D(\widetilde h_t),a),
\]
and likewise
\[
\widetilde\Qcal_t^{-c}(\widetilde o\mid \widetilde h_t,a)
=
K(\widetilde o\mid D(\widetilde o))\,\Qcal_t^{-c}(D(\widetilde o)\mid D(\widetilde h_t),a).
\]
Therefore
\begin{align*}
\sum_{\widetilde o\in\widetilde\Ocal}
\sqrt{\widetilde\mu_c(\widetilde o\mid \widetilde h_t,a)\,
\widetilde\Qcal_t^{-c}(\widetilde o\mid \widetilde h_t,a)}
&=
\sum_{\widetilde o\in\widetilde\Ocal}
K(\widetilde o\mid D(\widetilde o))
\sqrt{\mu_c(D(\widetilde o)\mid D(\widetilde h_t),a)\,
\Qcal_t^{-c}(D(\widetilde o)\mid D(\widetilde h_t),a)}\\
&=
\sum_{o\in\Ocal}
\sqrt{\mu_c(o\mid D(\widetilde h_t),a)\,
\Qcal_t^{-c}(o\mid D(\widetilde h_t),a)},
\end{align*}
again by summing over the fibers of $D$. Taking $-\log$ gives the claimed identity for $\widetilde B_t$.
\end{proof}

\begin{corollary}[Exact transfer under decodable observation recodings]
\label{cor:noise-transfer}
Fix a decodable observation channel $K:\Ocal\to\Delta(\widetilde\Ocal)$, fix a target program $p^\star\in\Pcal$, write $c^\star:=[p^\star]$, and consider the corrupted-observation family $(\widetilde\mu_p)_{p\in\Pcal}$ of Proposition~\ref{prop:noise-decoding}, with Bayes/Gibbs posterior $\widetilde q_t$ formed from the same prior weights on program indices.

\begin{itemize}
\item[(i)] If the original family satisfies the semantic separation condition of Definition~\ref{def:sep-condition} with constants $\eta(c)$, then the corrupted-observation family satisfies the same condition with the same constants.
\item[(ii)] Suppose the learner uses the same universal-prior weights $w(p)$ on the common program index set, the Bayes/Gibbs posterior $\widetilde q_t$ on corrupted histories, and the semantic action rule of Definition~\ref{def:semantic-rule} computed from the transformed family $(\widetilde\mu_p)_{p\in\Pcal}$. Then
\[
\widetilde q_t([p^\star]\mid \widetilde H_t)\longrightarrow 1
\qquad\text{almost surely.}
\]
\item[(iii)] Suppose the learner uses the same universal-prior weights $w(p)$ on the common program index set, the Bayes/Gibbs posterior $\widetilde q_t$ on corrupted histories, and the kernel semantic action rule of Definition~\ref{def:kernel-semantic-rule} computed from the transformed family $(\widetilde\mu_p)_{p\in\Pcal}$. Then
\[
\widetilde q_t([p^\star]\mid \widetilde H_t)\longrightarrow 1
\qquad\text{almost surely.}
\]
\item[(iv)] Under the corresponding bounded-ratio hypotheses of Propositions~\ref{prop:exp-rate} and~\ref{prop:kernel-exp-rate}, and Theorem~\ref{thm:exp-rate-concentration}, the expectation and high-probability rate bounds transfer with the same selector half-cap constant
\[
\hat\eta(c^\star):=\frac{1}{2}\min\{1,\eta(c^\star)\},
\]
the same kernel separation constant $\eta(c^\star)$, the same prior-mass constant $\bar w(c^\star)$, the same derived drift constants $\mu_0$ and $\mu_{0,\lambda}$, and the same likelihood-ratio constant $L=\log(1/m)$.
\end{itemize}
\end{corollary}

\begin{proof}
Part \textup{(i)} is immediate from Proposition~\ref{prop:noise-decoding}\,\textup{(iii)}: if
\[
\sup_{a\in\Acal}B_t(c,a;D(\widetilde h_t))\ge \eta(c),
\]
then
\[
\sup_{a\in\Acal}\widetilde B_t(c,a;\widetilde h_t)\ge \eta(c).
\]
For part \textup{(ii)}, Proposition~\ref{prop:noise-decoding} preserves the semantic partition, the posterior on program indices, and the Bhattacharyya separation geometry, while the prior class masses $\bar w(c)$ are unchanged because the common program index set and its partition are unchanged. Define
\[
\hat\eta(c):=\frac{1}{2}\min\{1,\eta(c)\}.
\]
The corrupted family therefore satisfies the semantic separation condition with the smaller constants $\hat\eta(c)$ as well. If
\[
\widetilde S_t^\star:=\{a\in\Acal:\widetilde B_t(c^\star,a;\widetilde H_t)\ge \hat\eta(c^\star)\},
\]
then the same promotion-support calculation as Proposition~\ref{prop:semantic-is-promotion-supporting} gives
\[
\pi_{t+1}^{\mathrm{sem}}(\widetilde S_t^\star\mid \widetilde H_t)\ge \delta_0\,\bar w(c^\star)
\qquad\text{almost surely on }\{\widetilde q_t(c^\star\mid \widetilde H_t)<1\},
\]
where $\pi^{\mathrm{sem}}$ is the semantic-rule action marginal on corrupted histories. Applying Theorem~\ref{thm:separating-support-convergence} to the corrupted family yields part \textup{(ii)}.

For part \textup{(iii)}, let
\[
\widetilde S_t^\star:=\{a\in\Acal:\widetilde B_t(c^\star,a;\widetilde H_t)\ge \eta(c^\star)\}.
\]
Define
\[
\lambda_{\min}:=\min_{a\in\Acal}\lambda(a)>0
\]
for the same full-support reference distribution $\lambda$ used in Definition~\ref{def:kernel-semantic-rule}. By part \textup{(i)} and Proposition~\ref{prop:finite-action-kernel-separation}, one has
\[
\lambda(\widetilde S_t^\star)\ge \lambda_{\min}
\qquad\text{almost surely on }\{\widetilde q_t(c^\star\mid \widetilde H_t)<1\}.
\]
Since Proposition~\ref{prop:kernel-promotion-support} depends only on $\bar w(c)$ and $\lambda$, the corrupted kernel rule satisfies
\[
\pi_{t+1}^{\mathrm{ker}}(\widetilde S_t^\star\mid \widetilde H_t)\ge \delta_0\,\bar w(c^\star)\,\lambda_{\min}
\qquad\text{almost surely on }\{\widetilde q_t(c^\star\mid \widetilde H_t)<1\},
\]
where $\pi^{\mathrm{ker}}$ is the action marginal of the corrupted kernel rule. Applying Theorem~\ref{thm:separating-support-convergence} to the corrupted family yields part \textup{(iii)}.

For part \textup{(iv)}, Proposition~\ref{prop:noise-decoding}\,\textup{(i)} gives the same posterior odds, and the displayed identities in its proof give
\[
\frac{\widetilde\mu_{c^\star}(\widetilde o\mid \widetilde h_t,a)}
{\widetilde\Qcal_t^{-c^\star}(\widetilde o\mid \widetilde h_t,a)}
=
\frac{\mu_{c^\star}(D(\widetilde o)\mid D(\widetilde h_t),a)}
{\Qcal_t^{-c^\star}(D(\widetilde o)\mid D(\widetilde h_t),a)}.
\]
Hence the same bounded-ratio constant $m$ applies. For the semantic rule, apply Theorem~\ref{thm:separating-support-rate} with the smaller separation constant
\[
\hat\eta(c^\star):=\frac{1}{2}\min\{1,\eta(c^\star)\}
\]
and with
\[
r_t\equiv \delta_0\,\bar w(c^\star),
\]
which reproduces the selector drift constant $\mu_0=2\hat\eta(c^\star)\delta_0\bar w(c^\star)$. For the kernel rule, define
\[
\lambda_{\min}:=\min_{a\in\Acal}\lambda(a)>0
\]
for the same full-support reference distribution $\lambda$ used in Definition~\ref{def:kernel-semantic-rule}, and apply Theorem~\ref{thm:separating-support-rate} with
\[
r_t\equiv \delta_0\,\bar w(c^\star)\,\lambda_{\min},
\]
which reproduces $\mu_{0,\lambda}=2\eta(c^\star)\delta_0\bar w(c^\star)\lambda_{\min}$. This yields the stated expectation and high-probability bounds with the same constants.
\end{proof}

\begin{corollary}[Left-invertible noisy transfer in the history-independent setting]
\label{cor:noise-left-invertible-history-independent}
Assume $\Ocal$ and $\widetilde\Ocal$ are finite, let $K:\Ocal\to\Delta(\widetilde\Ocal)$ be left-invertible, fix a target program $p^\star\in\Pcal$, write $c^\star:=[p^\star]$, and suppose the environment family is history-independent in the sense that
\[
\mu_p(o\mid h_t,a)=\mu_p(o\mid a)
\qquad
\text{for every }p\in\Pcal,\ h_t,\ a\in\Acal,\ o\in\Ocal.
\]
For each $p\in\Pcal$, define the corrupted-observation family by
\[
\widetilde\mu_p(\widetilde o\mid h_t,a):=\sum_{o\in\Ocal}K(\widetilde o\mid o)\mu_p(o\mid a).
\]
Then:
\begin{itemize}
\item[(i)] the behavioral and semantic equivalence relations induced by $(\widetilde\mu_p)_{p\in\Pcal}$ coincide with those induced by $(\mu_p)_{p\in\Pcal}$;
\item[(ii)] if the original family satisfies the uniform history-independent semantic separation condition of Definition~\ref{def:uniform-history-independent-separation} with constants $\eta(c)$, then the corrupted family satisfies the semantic separation condition of Definition~\ref{def:sep-condition} with constants $\widetilde\eta_K(\eta(c))$ from Proposition~\ref{prop:noise-left-invertible};
\item[(iii)] if the learner uses the same prior weights on program indices, Bayes/Gibbs belief for the corrupted family, and either the semantic action rule or the kernel semantic action rule computed from $(\widetilde\mu_p)_{p\in\Pcal}$, then the corresponding conclusions of Theorems~\ref{thm:semantic-convergence} and~\ref{thm:kernel-semantic-convergence} hold with $\eta(c)$ replaced by $\widetilde\eta_K(\eta(c))$;
\item[(iv)] if, in addition, there exists $m\in(0,1)$ such that for every action $a\in\Acal$ and every probability vector $\alpha\in\Delta(\Ccal\setminus\{c^\star\})$, the clean likelihood $\mu_{c^\star}(\cdot\mid a)$ and complement mixture $\Qcal_{\alpha}^{-c^\star}(\cdot\mid a)$ have the same support and satisfy
\[
m\le
\frac{\mu_{c^\star}(o\mid a)}{\Qcal_{\alpha}^{-c^\star}(o\mid a)}
\le \frac1m,
\]
for every $o$ in that common support, then the corresponding expectation and high-probability rate bounds transfer as well, with the same likelihood-ratio constant $m$, with the selector drift constant obtained from
\[
\widehat{\widetilde\eta}_K(c^\star):=\frac{1}{2}\min\{1,\widetilde\eta_K(\eta(c^\star))\},
\]
and with the kernel drift constant obtained by replacing $\eta(c^\star)$ by $\widetilde\eta_K(\eta(c^\star))$.
\end{itemize}
\end{corollary}

\begin{proof}
For part \textup{(i)}, history-independence means that behavioral and semantic equivalence are equality of the action-indexed one-step laws:
\[
p_1\sim_{\mathrm{behav}}p_2
\quad\Longleftrightarrow\quad
\mu_{p_1}(\cdot\mid a)=\mu_{p_2}(\cdot\mid a)\text{ for every }a\in\Acal.
\]
The same is true for the corrupted family with $\widetilde\mu_p(\cdot\mid a)=T_K\mu_p(\cdot\mid a)$. If
\[
\widetilde\mu_{p_1}(\cdot\mid a)=\widetilde\mu_{p_2}(\cdot\mid a)
\qquad\text{for every }a,
\]
then injectivity of $T_K$ gives
\[
\mu_{p_1}(\cdot\mid a)=\mu_{p_2}(\cdot\mid a)
\qquad\text{for every }a.
\]
Thus the two families induce the same behavioral equivalence, and therefore the same semantic equivalence as well.

For part \textup{(ii)}, fix a class $c\in\Ccal$ and a corrupted history $h_t$ with $0<\widetilde q_t(c\mid h_t)<1$. By part \textup{(i)}, the semantic classes of the corrupted family are the same classes $d\in\Ccal$. Because the corrupted family is again history-independent, there is a probability vector $\alpha_t^{h,c}\in\Delta(\Ccal\setminus\{c\})$ such that
\[
\widetilde\Qcal_t^{-c}(\cdot\mid h_t,a)
=
\sum_{d\in\Ccal\setminus\{c\}}\alpha_t^{h,c}(d)\,\widetilde\mu_d(\cdot\mid a)
\qquad\text{for every }a\in\Acal,
\]
namely
\[
\alpha_t^{h,c}(d):=
\frac{\widetilde q_t(d\mid h_t)}{1-\widetilde q_t(c\mid h_t)}.
\]
The clean family's corresponding class-complement mixture is
\[
\Qcal_{\alpha_t^{h,c}}^{-c}(\cdot\mid a)
=
\sum_{d\in\Ccal\setminus\{c\}}\alpha_t^{h,c}(d)\,\mu_d(\cdot\mid a).
\]
By Definition~\ref{def:uniform-history-independent-separation}, there exists an action $a^\star$ with
\[
-\log\sum_{o\in\Ocal}\sqrt{\mu_c(o\mid a^\star)\Qcal_{\alpha_t^{h,c}}^{-c}(o\mid a^\star)}
\ge \eta(c).
\]
Since
\[
\widetilde\mu_c(\cdot\mid a^\star)=T_K\mu_c(\cdot\mid a^\star),
\qquad
\widetilde\Qcal_t^{-c}(\cdot\mid h_t,a^\star)=T_K\Qcal_{\alpha_t^{h,c}}^{-c}(\cdot\mid a^\star),
\]
Proposition~\ref{prop:noise-left-invertible} gives
\[
\widetilde B_t(c,a^\star;h_t)\ge \widetilde\eta_K(\eta(c)).
\]
As $h_t$ was arbitrary, the corrupted family satisfies the semantic separation condition with constants $\widetilde\eta_K(\eta(c))$.

For part \textup{(iii)}, part \textup{(i)} preserves the target class $[p^\star]$, while part \textup{(ii)} supplies the corrupted semantic separation constant
\[
\widetilde\eta_K(\eta(c^\star))>0.
\]
For the semantic rule, define
\[
\widehat{\widetilde\eta}_K(c^\star):=\frac{1}{2}\min\{1,\widetilde\eta_K(\eta(c^\star))\}
\]
and let
\[
\widetilde S_t^\star:=\{a\in\Acal:\widetilde B_t(c^\star,a;H_t)\ge \widehat{\widetilde\eta}_K(c^\star)\}.
\]
The same promotion-support calculation as Proposition~\ref{prop:semantic-is-promotion-supporting} gives
\[
\pi_{t+1}^{\mathrm{sem}}(\widetilde S_t^\star\mid H_t)\ge \delta_0\,\bar w(c^\star)
\qquad\text{almost surely on }\{\widetilde q_t(c^\star\mid H_t)<1\},
\]
so Theorem~\ref{thm:separating-support-convergence} yields semantic recovery for the corrupted semantic rule.

For the corrupted kernel rule, let
\[
\widetilde S_t^\star:=\{a\in\Acal:\widetilde B_t(c^\star,a;H_t)\ge \widetilde\eta_K(\eta(c^\star))\}.
\]
Because $\Acal$ is finite and $\lambda$ has full support, Proposition~\ref{prop:finite-action-kernel-separation} gives
\[
\lambda(\widetilde S_t^\star)\ge \lambda_{\min},
\qquad
\lambda_{\min}:=\min_{a\in\Acal}\lambda(a)>0,
\]
and Proposition~\ref{prop:kernel-promotion-support} therefore gives
\[
\pi_{t+1}^{\mathrm{ker}}(\widetilde S_t^\star\mid H_t)\ge \delta_0\,\bar w(c^\star)\,\lambda_{\min}
\qquad\text{almost surely on }\{\widetilde q_t(c^\star\mid H_t)<1\}.
\]
Applying Theorem~\ref{thm:separating-support-convergence} yields the kernel conclusion.

For part \textup{(iv)}, fix a corrupted history $h_t$ with $0<\widetilde q_t(c^\star\mid h_t)<1$ and an action $a\in\Acal$, and define
\[
\alpha_t^{h,c^\star}(d):=
\frac{\widetilde q_t(d\mid h_t)}{1-\widetilde q_t(c^\star\mid h_t)}
\qquad
(d\in\Ccal\setminus\{c^\star\}).
\]
Then
\[
\widetilde\mu_{c^\star}(\cdot\mid h_t,a)=T_K\mu_{c^\star}(\cdot\mid a),
\qquad
\widetilde\Qcal_t^{-c^\star}(\cdot\mid h_t,a)=T_K\Qcal_{\alpha_t^{h,c^\star}}^{-c^\star}(\cdot\mid a).
\]
Suppose
\[
m\le \frac{x(o)}{y(o)}\le \frac1m
\qquad\text{for every }o\in\operatorname{supp}(x)=\operatorname{supp}(y)
\]
for $x=\mu_{c^\star}(\cdot\mid a)$ and $y=\Qcal_{\alpha_t^{h,c^\star}}^{-c^\star}(\cdot\mid a)$. Since $x$ and $y$ vanish together off their common support, it follows that
\[
m\,y(o)\le x(o)\le \frac1m\,y(o)
\qquad\text{for every }o\in\Ocal.
\]
Then for every $\widetilde o\in\widetilde\Ocal$,
\[
m\,(T_Ky)(\widetilde o)
\le
(T_Kx)(\widetilde o)
\le
\frac1m\,(T_Ky)(\widetilde o),
\]
by linearity and the nonnegativity of $K(\widetilde o\mid o)$. So the same bounded-ratio constant $m$ holds for the corrupted family. Applying Theorem~\ref{thm:separating-support-rate} with
\[
r_t\equiv \delta_0\,\bar w(c^\star)
\]
for the semantic rule and with
\[
r_t\equiv \delta_0\,\bar w(c^\star)\,\lambda_{\min}
\]
for the kernel rule yields the corresponding expectation and high-probability bounds, with the semantic-rule separation constant $\widehat{\widetilde\eta}_K(c^\star)$ and the kernel separation constant $\widetilde\eta_K(\eta(c^\star))$.
\end{proof}

\begin{remark}[Against entropy-bonus curiosity]
\label{rem:against-entropy-bonus}
Proposition~\ref{prop:noise-immunity} separates semantic separation from entropy-bonus objectives at the one-step, fixed-history level. Entropy-bonus curiosity rewards actions whose observations carry high marginal entropy $H(O\mid h_t,a)$; a noise channel $K$ of high entropy therefore inflates the bonus while, by Proposition~\ref{prop:noise-immunity}, leaving (or shrinking) the corresponding Bhattacharyya separation. Corollary~\ref{cor:noise-transfer} isolates the exact-transfer regime: if the corruption is decodable, the semantic posterior and separation geometry are preserved exactly and the main convergence theorems survive unchanged. Proposition~\ref{prop:noise-left-invertible} identifies a second, weaker positive regime: for finite observation alphabets and left-invertible channels, positive clean Bhattacharyya separation degrades to a still-positive uniform post-channel constant at a fixed latent history. Corollary~\ref{cor:noise-left-invertible-history-independent} promotes this to a full noisy-learning transfer in the history-independent setting under a uniform class-vs.-convex-complement separation hypothesis, together with the analogous uniform bounded-ratio regularity for the rate bounds. Outside these structured regimes, the class-against-complement rule is monotone under channel noise only in the pointwise sense of Proposition~\ref{prop:noise-immunity}; it never confuses extra noise for extra signal at a fixed latent history, but a full learning-level transfer is not automatic. This is one concrete sense in which the canonical action rule at $\omega_{\mathrm{behav}}$ differs in kind from the expected-free-energy action selector when the latter is read as an MI-plus-entropy bonus --- the decomposition used in the active-inference literature (cf.\ the MI-variant $\mathcal J_{t,\mathrm{raw}}^{\omega_B,\omega_A,\mathrm{MI}}$ after Definition~\ref{def:raw-two-observer-functional} and the discussion of expected free energy in Remark~\ref{rem:efe-unification}).
\end{remark}

\section{Sufficient Conditions for Semantic Convergence}
\label{sec:main}

Section~\ref{sec:functional} fixed the canonical action observer at $\omega_A=\omega_{\mathrm{behav}}$ (Corollary~\ref{cor:canonical-pair}) and Section~\ref{sec:belief} fixed the canonical belief observer at $\omega_B=\omega_{\mathrm{syn}}$; this section proves what happens once the learner acts there. Theorem~\ref{thm:separating-support-convergence} answers the boxed problem: under the universal prior, Bayes/Gibbs belief, semantic separation, and divergent cumulative policy mass on separating actions, the posterior on $[p^\star]$ converges almost surely to one. Two canonical realizations of that support condition are proved inside the section --- the class-against-complement selector (Theorem~\ref{thm:semantic-convergence}) and the selector-free Gibbs kernel lift (Theorem~\ref{thm:kernel-semantic-convergence}); Corollary~\ref{cor:compact-action-kernel} carries the kernel route to compact metric action spaces under action-continuity on the Bhattacharyya score and a full-support Borel reference measure. Matched converses show sharpness: zero separating support makes uniform recovery impossible (Corollary~\ref{cor:support-necessary}), and summable separating support still admits failure with positive probability (Theorem~\ref{thm:summable-support-insufficient}). A divergent cumulative exploration budget delivers the same guarantee on its own (Theorem~\ref{thm:exploration-floor-behavioral}). Theorem~\ref{thm:separating-support-rate} is the quantitative companion: under bounded log-likelihood ratios, the same support condition yields explicit finite-time high-probability lower bounds on the true-class posterior.

\begin{definition}[Semantic separation condition]
\label{def:sep-condition}
The interactive setting satisfies the \emph{semantic separation condition} if for every class $c\in\Ccal$ there exists $\eta(c)>0$ such that whenever $0<q_t^\star(c\mid h_t)<1$,
\[
\sup_{a\in\Acal}B_t(c,a;h_t)\ge \eta(c).
\]
\end{definition}

\begin{definition}[Uniform history-independent semantic separation]
\label{def:uniform-history-independent-separation}
Assume the environment family is history-independent in the sense that
\[
\mu_p(o\mid h_t,a)=\mu_p(o\mid a)
\qquad
\text{for every }p\in\Pcal,\ h_t,\ a\in\Acal,\ o\in\Ocal.
\]
For each semantic class $c\in\Ccal$, write $\mu_c(\cdot\mid a)$ for the common one-step law of programs in $c$ under action $a$. For any class $c$ with $\Ccal\setminus\{c\}\neq\varnothing$ and any probability vector $\alpha\in\Delta(\Ccal\setminus\{c\})$, define the class-complement mixture
\[
\Qcal_{\alpha}^{-c}(\cdot\mid a)
:=
\sum_{d\in\Ccal\setminus\{c\}}\alpha(d)\mu_d(\cdot\mid a).
\]
The history-independent family satisfies the \emph{uniform history-independent semantic separation condition} if for every class $c\in\Ccal$ with $\Ccal\setminus\{c\}\neq\varnothing$ there exists $\eta(c)>0$ such that for every $\alpha\in\Delta(\Ccal\setminus\{c\})$,
\[
\sup_{a\in\Acal}
\left(
-\log\sum_{o\in\Ocal}\sqrt{\mu_c(o\mid a)\Qcal_{\alpha}^{-c}(o\mid a)}
\right)
\ge \eta(c).
\]
\end{definition}

\begin{proposition}[Uniform history-independent separation implies semantic separation]
\label{prop:uniform-history-independent-implies-semantic}
Assume the environment family is history-independent and satisfies the uniform history-independent semantic separation condition. Then the semantic separation condition holds with the same constants $\eta(c)$.
\end{proposition}

\begin{proof}
Fix a class $c\in\Ccal$ and a history $h_t$ with $0<q_t^\star(c\mid h_t)<1$. Define
\[
\alpha_t^{h,c}(d):=
\frac{q_t^\star(d\mid h_t)}{1-q_t^\star(c\mid h_t)}
\qquad
(d\in\Ccal\setminus\{c\}).
\]
Because the family is history-independent, the class-complement predictive law is
\[
\Qcal_t^{-c}(\cdot\mid h_t,a)=\Qcal_{\alpha_t^{h,c}}^{-c}(\cdot\mid a)
\qquad\text{for every }a\in\Acal.
\]
By Definition~\ref{def:uniform-history-independent-separation}, there exists an action $a^\star$ such that
\[
-\log\sum_{o\in\Ocal}\sqrt{\mu_c(o\mid a^\star)\Qcal_{\alpha_t^{h,c}}^{-c}(o\mid a^\star)}
\ge \eta(c).
\]
The displayed quantity is exactly $B_t(c,a^\star;h_t)$, so
\[
\sup_{a\in\Acal}B_t(c,a;h_t)\ge \eta(c).
\]
Hence the semantic separation condition holds.
\end{proof}

\begin{corollary}[Uniform KL class-vs.-complement separation implies semantic separation]
\label{cor:kl-implies-semantic-separation}
Suppose that for every class $c\in\Ccal$ there exists $\kappa(c)>0$ such that whenever $0<q_t^\star(c\mid h_t)<1$,
\[
\sup_{a\in\Acal}
\KL\!\bigl(\mu_c(\cdot\mid h_t,a)\,\|\,\Qcal_t^{-c}(\cdot\mid h_t,a)\bigr)
\ge \kappa(c).
\]
Then the semantic separation condition holds with
\[
\eta(c):=\frac{\kappa(c)}{2}.
\]
\end{corollary}

\begin{proof}
Fix $c$ and a history $h_t$ with $0<q_t^\star(c\mid h_t)<1$. By hypothesis there exists an action $a^\star$ with
\[
\KL\!\bigl(\mu_c(\cdot\mid h_t,a^\star)\,\|\,\Qcal_t^{-c}(\cdot\mid h_t,a^\star)\bigr)\ge \kappa(c).
\]
Proposition~\ref{prop:chernoff-correspondence}\,\textup{(ii)} gives
\[
B_t(c,a^\star;h_t)\ge \frac{\kappa(c)}{2}=\eta(c).
\]
Thus $\sup_a B_t(c,a;h_t)\ge \eta(c)$.
\end{proof}

\begin{corollary}[Event-witness separation implies semantic separation]
\label{cor:event-witness-implies-semantic-separation}
Suppose that for every class $c\in\Ccal$ there exist constants
\[
0\le v_c<u_c\le 1
\]
such that whenever $0<q_t^\star(c\mid h_t)<1$, there exist an action $a\in\Acal$ and an event $E\subseteq\Ocal$ with
\[
\mu_c(E\mid h_t,a)\ge u_c,
\qquad
\Qcal_t^{-c}(E\mid h_t,a)\le v_c.
\]
Then the semantic separation condition holds with
\[
\eta(c):=-\log\!\bigl(\sqrt{u_cv_c}+\sqrt{(1-u_c)(1-v_c)}\bigr)>0.
\]
\end{corollary}

\begin{proof}
Fix $c$ and a history $h_t$ with $0<q_t^\star(c\mid h_t)<1$, and choose $a$ and $E$ as in the hypothesis. Write
\[
x:=\mu_c(E\mid h_t,a),
\qquad
y:=\Qcal_t^{-c}(E\mid h_t,a).
\]
Then $x\ge u_c>v_c\ge y$, so $x>y$. By Cauchy--Schwarz applied separately on $E$ and $E^c$,
\[
\sum_{o\in\Ocal}\sqrt{\mu_c(o\mid h_t,a)\Qcal_t^{-c}(o\mid h_t,a)}
\le
\sqrt{xy}+\sqrt{(1-x)(1-y)}.
\]
The function
\[
g(x,y):=\sqrt{xy}+\sqrt{(1-x)(1-y)}
\]
is decreasing in $x$ and increasing in $y$ on the region $x>y$, so on the region $x\ge u_c>v_c\ge y$ one has
\[
g(x,y)\le g(u_c,v_c)=\sqrt{u_cv_c}+\sqrt{(1-u_c)(1-v_c)}.
\]
Therefore
\[
B_t(c,a;h_t)
\ge
-\log\!\bigl(\sqrt{u_cv_c}+\sqrt{(1-u_c)(1-v_c)}\bigr)
=\eta(c).
\]
Thus $\sup_a B_t(c,a;h_t)\ge \eta(c)$.
\end{proof}

\begin{remark}[Why no pairwise-Bhattacharyya corollary]
\label{rem:no-pairwise-bhat-corollary}
We do not state a corollary that lower bounds $B_t(c,a;h_t)$ from pairwise Bhattacharyya lower bounds against each individual rival class. Such pairwise lower bounds need not survive convex mixing into the aggregate complement law $\Qcal_t^{-c}$, so they are not a reliable operational proxy for Definition~\ref{def:sep-condition}. The event-witness criterion above is preferred precisely because it is stable under convex aggregation of the complement.
\end{remark}

\begin{definition}[Kernel semantic separation condition relative to $\lambda$]
\label{def:kernel-sep-condition}
If one relaxes the standing finite-action setup and instead allows either a countable action alphabet with a full-support reference distribution $\lambda\in\Delta(\Acal)$ or a measurable action space equipped with a full-support reference probability measure $\lambda$, the interactive setting satisfies the \emph{kernel semantic separation condition relative to $\lambda$} if every class $c\in\Ccal$ admits constants $\eta_\lambda(c),\rho_\lambda(c)>0$ such that whenever $0<q_t^\star(c\mid h_t)<1$,
\[
\lambda\!\bigl(\{a\in\Acal:B_t(c,a;h_t)\ge \eta_\lambda(c)\}\bigr)\ge \rho_\lambda(c).
\]
\end{definition}

\begin{proposition}[Finite-action lift of semantic separation]
\label{prop:finite-action-kernel-separation}
Fix a full-support reference distribution $\lambda\in\Delta(\Acal)$ and define
\[
\lambda_{\min}:=\min_{a\in\Acal}\lambda(a)>0.
\]
If the semantic separation condition holds, then for every class $c\in\Ccal$ and every history $h_t$ with $0<q_t^\star(c\mid h_t)<1$,
\[
\lambda\!\bigl(\{a\in\Acal:B_t(c,a;h_t)\ge \eta(c)\}\bigr)\ge \lambda_{\min}.
\]
\end{proposition}

\begin{proof}
By Definition~\ref{def:sep-condition},
\[
\sup_{a\in\Acal}B_t(c,a;h_t)\ge \eta(c).
\]
Because $\Acal$ is finite, the supremum is a maximum, so there exists at least one action $a^\star\in\Acal$ with $B_t(c,a^\star;h_t)\ge \eta(c)$. The displayed set is therefore nonempty. Since $\Acal$ is finite and $\lambda$ has full support, every singleton has mass at least $\lambda_{\min}$, and hence every nonempty subset of $\Acal$ has $\lambda$-mass at least $\lambda_{\min}$.
\end{proof}

\begin{proposition}[Compact-action lift of semantic separation]
\label{prop:compact-action-kernel-separation}
If one relaxes the standing finite-action setup and instead lets $(\Acal,d_\Acal)$ be a compact metric space equipped with a full-support Borel probability measure $\lambda$, assume that for each class $c\in\Ccal$ there exists a modulus of continuity $\varpi_c:[0,\infty)\to[0,\infty)$ with $\varpi_c(r)\to 0$ as $r\downarrow 0$ and such that whenever $0<q_t^\star(c\mid h_t)<1$,
\[
\bigl|B_t(c,a;h_t)-B_t(c,a';h_t)\bigr|
\le
\varpi_c\!\bigl(d_\Acal(a,a')\bigr)
\qquad
\text{for all }a,a'\in\Acal.
\]
If the semantic separation condition holds, then the kernel semantic separation condition relative to $\lambda$ holds with
\[
\eta_\lambda(c):=\frac{\eta(c)}{2},
\qquad
\rho_\lambda(c):=m_\lambda(r_c),
\]
where $r_c>0$ is any radius satisfying $\varpi_c(r_c)\le \eta(c)/2$ and
\[
U_r(a):=\{a'\in\Acal:d_\Acal(a,a')<r\},
\qquad
m_\lambda(r):=\inf_{a\in\Acal}\lambda\!\bigl(U_r(a)\bigr).
\]
\end{proposition}

\begin{proof}
Fix $r>0$. Because $\lambda$ has full support, every nonempty open ball in $(\Acal,d_\Acal)$ has positive $\lambda$-mass. Compactness gives a finite cover
\[
\Acal=\bigcup_{i=1}^N U_{r/2}(a_i).
\]
For any $a\in\Acal$, choose $i$ with $a\in U_{r/2}(a_i)$. Then for every $a'\in U_{r/2}(a_i)$,
\[
d_\Acal(a',a)\le d_\Acal(a',a_i)+d_\Acal(a_i,a)<r,
\]
so $U_{r/2}(a_i)\subseteq U_r(a)$. Therefore
\[
\lambda\!\bigl(U_r(a)\bigr)\ge \lambda\!\bigl(U_{r/2}(a_i)\bigr)\ge \min_{1\le i\le N}\lambda\!\bigl(U_{r/2}(a_i)\bigr)>0.
\]
Taking the infimum over $a$ shows $m_\lambda(r)>0$ for every $r>0$.

Now fix a class $c\in\Ccal$ and choose $r_c>0$ so that $\varpi_c(r_c)\le \eta(c)/2$. Consider any history $h_t$ with $0<q_t^\star(c\mid h_t)<1$. By the semantic separation condition,
\[
\sup_{a\in\Acal}B_t(c,a;h_t)\ge \eta(c).
\]
The modulus assumption implies that $a\mapsto B_t(c,a;h_t)$ is continuous on compact $\Acal$, so the supremum is attained at some $a_c^\star\in\Acal$. Hence
\[
B_t(c,a_c^\star;h_t)\ge \eta(c).
\]
For every $a\in U_{r_c}(a_c^\star)$,
\[
B_t(c,a;h_t)\ge B_t(c,a_c^\star;h_t)-\varpi_c\!\bigl(d_\Acal(a,a_c^\star)\bigr)\ge \eta(c)-\varpi_c(r_c)\ge \frac{\eta(c)}{2}.
\]
Therefore
\[
U_{r_c}(a_c^\star)\subseteq \{a\in\Acal:B_t(c,a;h_t)\ge \eta(c)/2\},
\]
so
\[
\lambda\!\bigl(\{a\in\Acal:B_t(c,a;h_t)\ge \eta(c)/2\}\bigr)
\ge
\lambda\!\bigl(U_{r_c}(a_c^\star)\bigr)
\ge
m_\lambda(r_c)>0.
\]
This is exactly the kernel semantic separation condition relative to $\lambda$ with the displayed constants.
\end{proof}

\begin{lemma}[Conditional Borel--Cantelli]
\label{lem:conditional-bc}
Let $(\mathcal{F}_t)_{t\ge 0}$ be a filtration and $E_{t+1}\in\mathcal{F}_{t+1}$ for every $t\ge 0$. If $\sum_t\Prob(E_{t+1}\mid\mathcal{F}_t)=\infty$ almost surely, then $E_{t+1}$ i.o.\ almost surely.
\end{lemma}

\begin{proof}
See \citet[Chapter~12]{Williams1991}.
\end{proof}

\begin{lemma}[Cumulative separation drives posterior odds to zero]
\label{lem:contraction}
Fix a target program $p^\star\in\Pcal$ and write $c^\star:=[p^\star]$. Under the universal prior and Bayes/Gibbs posterior, if
\[
\sum_{t=0}^{\infty}B_t(c^\star,A_{t+1};H_t)=\infty
\qquad
\text{almost surely,}
\]
then $q_t^\star(c^\star\mid H_t)\to 1$ almost surely.
\end{lemma}

\begin{proof}
Let $\mathcal{G}_t:=\sigma(H_t,A_{t+1})$. Define the posterior odds against $c^\star$ by
\[
R_t^-:=\frac{1-q_t^\star(c^\star\mid H_t)}{q_t^\star(c^\star\mid H_t)}
=\frac{\sum_{p\notin c^\star}w(p)\mu_p(o_{1:t}\givena a_{1:t})}
      {\bar w(c^\star)\mu_{p^\star}(o_{1:t}\givena a_{1:t})}.
\]
Whenever $R_t^->0$, Definition~\ref{def:class-complement} gives
\[
R_{t+1}^-=R_t^-\cdot\frac{\Qcal_t^{-c^\star}(o_{t+1}\mid H_t,A_{t+1})}{\mu_{p^\star}(o_{t+1}\mid H_t,A_{t+1})}.
\]
Taking square roots and conditional expectation under the true law,
\[
\E\!\left[\sqrt{R_{t+1}^-}\mid\mathcal{G}_t\right]
=
\sqrt{R_t^-}\sum_{o\in\Ocal}\sqrt{\mu_{p^\star}(o\mid H_t,A_{t+1})\Qcal_t^{-c^\star}(o\mid H_t,A_{t+1})}
=
\sqrt{R_t^-}\,e^{-B_t(c^\star,A_{t+1};H_t)}.
\]
Let $S_t:=\sum_{i=0}^{t-1}B_i(c^\star,A_{i+1};H_i)$ and $M_t:=e^{S_t}\sqrt{R_t^-}$. Then $(M_t)$ is a nonnegative $\mathcal{G}_t$-martingale, hence converges almost surely. On $\{S_t\to\infty\}$, $\sqrt{R_t^-}=M_te^{-S_t}\to 0$, so $R_t^-\to 0$ and $q_t^\star(c^\star\mid H_t)\to 1$.
\end{proof}

\begin{proposition}[Full-support policies recover the behavioral observer]
\label{prop:full-support-behavioral}
Let $\pi$ be any policy satisfying
\[
\pi_{t+1}(a\mid h_t)>0
\qquad
\text{for every }t\ge 0,\ h_t,\ a\in\Acal.
\]
Then
\[
\omega_\pi=\omega_{\mathrm{behav}}.
\]
\end{proposition}

\begin{proof}
Lemma~\ref{lem:monotone-refinement} already gives $\omega_\pi\le\omega_{\mathrm{behav}}$, so only the reverse inclusion needs proof. Suppose
\[
p_1\sim_{\omega_\pi}p_2,
\]
that is, $\Prob_{p_1,\pi}=\Prob_{p_2,\pi}$ on the full history space. We show $p_1\sim_{\mathrm{behav}}p_2$.

Fix any history-action pair $(h_t,a)$ reachable under either program in the sense of the standing convention after the formal setup: assume
\[
\mu_{p_1}(o_{1:t}\givena a_{1:t})>0
\qquad\text{or}\qquad
\mu_{p_2}(o_{1:t}\givena a_{1:t})>0.
\]
Because every policy factor along the action prefix is positive, this implies
\[
\Prob_{p_1,\pi}(H_t=h_t)>0
\qquad\text{or}\qquad
\Prob_{p_2,\pi}(H_t=h_t)>0.
\]
Since the two policy laws are equal, in fact
\[
\Prob_{p_1,\pi}(H_t=h_t)=\Prob_{p_2,\pi}(H_t=h_t)>0.
\]
For each $o\in\Ocal$, the factorization identity gives
\[
\mu_{p_i}(o\mid h_t,a)
=
\frac{\Prob_{p_i,\pi}(H_{t+1}=(h_t,a,o))}
{\Prob_{p_i,\pi}(H_t=h_t)\,\pi_{t+1}(a\mid h_t)},
\qquad i\in\{1,2\},
\]
and the denominator is positive by hypothesis. Equality of the policy laws implies equality of the numerators and of the history probabilities in the denominator, while the policy factor is the same common $\pi_{t+1}(a\mid h_t)$. Hence
\[
\mu_{p_1}(o\mid h_t,a)=\mu_{p_2}(o\mid h_t,a)
\qquad\text{for every }o\in\Ocal.
\]
Because $(h_t,a)$ was an arbitrary pair reachable under either program, the two kernels agree at every reachable pair, so $p_1\sim_{\mathrm{behav}}p_2$. Thus $\omega_{\mathrm{behav}}\le\omega_\pi$, and equality follows.
\end{proof}

\begin{theorem}[Sufficient conditions for semantic recovery]
\label{thm:separating-support-convergence}
Assume \textup{(C1)--(C3)} and $\omega_B=\omega_{\mathrm{syn}}$ with Bayes/Gibbs belief \textup{(C2)}. Fix a policy $\pi$ and a target program $p^\star\in\Pcal$, write
\[
c^\star:=[p^\star],
\qquad
\mathcal F_t:=\sigma(H_t),
\qquad
E_t:=\{q_t^\star(c^\star\mid H_t)<1\},
\]
and define the separating-action set
\[
S_t^\star:=\{a\in\Acal:B_t(c^\star,a;H_t)\ge \eta(c^\star)\}.
\]
Suppose there exists a deterministic sequence $(r_t)_{t\ge 1}\subset[0,1]$ with
\[
\sum_{t=1}^\infty r_t=\infty
\]
such that for every $t\ge 0$,
\[
\pi_{t+1}(S_t^\star\mid H_t)\ge r_{t+1}
\qquad\text{almost surely on }E_t.
\]
Then
\[
q_t^\star([p^\star]\mid H_t)\longrightarrow 1
\qquad
\text{almost surely.}
\]
\end{theorem}

\begin{proof}
Define
\[
T_{t+1}:=E_t\cap\{A_{t+1}\in S_t^\star\}.
\]
Then
\[
\Prob(T_{t+1}\mid\mathcal F_t)
=
\one_{E_t}\,\pi_{t+1}(S_t^\star\mid H_t)
\ge
r_{t+1}\one_{E_t}.
\]
If $E_t$ occurs infinitely often, then in fact $E_t$ holds for every $t\ge 0$, because once $q_t^\star(c^\star\mid H_t)=1$ Bayes keeps the posterior mass outside $c^\star$ at zero forever. Hence on $E_t$ i.o.,
\[
\sum_{t=0}^\infty \Prob(T_{t+1}\mid\mathcal F_t)=\infty
\qquad\text{almost surely,}
\]
so Lemma~\ref{lem:conditional-bc} gives $T_{t+1}$ infinitely often almost surely on $E_t$ i.o. Whenever $T_{t+1}$ occurs,
\[
B_t(c^\star,A_{t+1};H_t)\ge \eta(c^\star),
\]
hence
\[
\sum_{t=0}^\infty B_t(c^\star,A_{t+1};H_t)=\infty
\qquad\text{almost surely on }E_t\text{ i.o.}
\]
and Lemma~\ref{lem:contraction} yields $q_t^\star(c^\star\mid H_t)\to 1$ on that event. If $E_t^c$ eventually, the same conclusion is immediate.
\end{proof}

\begin{theorem}[Full-support exploration floors recover the behavioral target]
\label{thm:exploration-floor-behavioral}
Assume the standing finite-action setup, \textup{(C1)--(C3)}, and $\omega_B=\omega_{\mathrm{syn}}$ with Bayes/Gibbs belief \textup{(C2)}. Fix a full-support reference distribution $\lambda\in\Delta(\Acal)$ and an exploration schedule $(\varepsilon_t)_{t\ge 1}$ satisfying
\[
\varepsilon_t\in(0,1]
\qquad\text{for every }t\ge 1,
\qquad
\sum_{t=1}^\infty \varepsilon_t=\infty.
\]
Let $\pi$ be any policy satisfying
\[
\pi_{t+1}(a\mid h_t)\ge \varepsilon_{t+1}\lambda(a)
\qquad
\text{for every }t\ge 0,\ h_t,\ a\in\Acal.
\]
Then for every $p^\star\in\Pcal$,
\[
\omega_\pi=\omega_{\mathrm{behav}},
\qquad
q_t^\star([p^\star]\mid H_t)\longrightarrow 1
\quad\text{almost surely.}
\]
\end{theorem}

\begin{proof}
The observer identity follows from Proposition~\ref{prop:full-support-behavioral}, since $\pi_{t+1}(a\mid h_t)\ge \varepsilon_{t+1}\lambda(a)>0$ for every $t,h_t,a$.

Fix $p^\star\in\Pcal$ and write $c^\star:=[p^\star]$. Let
\[
S_t^\star:=\{a\in\Acal:B_t(c^\star,a;H_t)\ge \eta(c^\star)\}.
\]
By Proposition~\ref{prop:finite-action-kernel-separation},
\[
\lambda(S_t^\star)\ge \lambda_{\min},
\qquad
\lambda_{\min}:=\min_{a\in\Acal}\lambda(a)>0,
\]
on $E_t=\{q_t^\star(c^\star\mid H_t)<1\}$. Therefore
\[
\pi_{t+1}(S_t^\star\mid H_t)\ge \varepsilon_{t+1}\lambda(S_t^\star)\ge \varepsilon_{t+1}\lambda_{\min}
\qquad\text{almost surely on }E_t.
\]
Applying Theorem~\ref{thm:separating-support-convergence} with
\[
r_t:=\varepsilon_t\lambda_{\min}
\]
gives the posterior-concentration claim.
\end{proof}

The same support condition admits a quantitative companion, phrased in terms of the true-class log-odds.

\begin{theorem}[Rates from cumulative separating-action support]
\label{thm:separating-support-rate}
Assume \textup{(C1)--(C3)} and $\omega_B=\omega_{\mathrm{syn}}$ with Bayes/Gibbs belief \textup{(C2)}. Fix a policy $\pi$, a target program $p^\star\in\Pcal$, write
\[
\begin{aligned}
c^\star&:=[p^\star],
&
Z_t&:=\log\frac{q_t^\star(c^\star\mid H_t)}{1-q_t^\star(c^\star\mid H_t)},
\\
\mathcal F_t&:=\sigma(H_t),
&
E_t&:=\{q_t^\star(c^\star\mid H_t)<1\},
\\
S_t^\star&:=\{a\in\Acal:B_t(c^\star,a;H_t)\ge \eta(c^\star)\},
\end{aligned}
\]
with the convention $Z_t=+\infty$ when $q_t^\star(c^\star\mid H_t)=1$, and suppose there exists a deterministic sequence $(r_t)_{t\ge 1}\subset[0,1]$ such that
\[
\pi_{t+1}(S_t^\star\mid H_t)\ge r_{t+1}
\qquad\text{almost surely on }E_t
\]
for every $t\ge 0$. Suppose in addition that there exists $m\in(0,1)$ such that whenever $0<q_t^\star(c^\star\mid H_t)<1$,
\[
m\le
\frac{\mu_{c^\star}(o\mid H_t,A_{t+1})}{\Qcal_t^{-c^\star}(o\mid H_t,A_{t+1})}
\le \frac{1}{m}
\]
for every $o$ with $\mu_{c^\star}(o\mid H_t,A_{t+1})>0$. Write
\[
R_t:=\sum_{i=1}^t r_i,
\qquad
\mu_{\mathrm{sep}}:=2\,\eta(c^\star),
\qquad
L:=\log(1/m).
\]
Then:
\begin{itemize}
\item[(i)] for every $t\ge 0$,
\[
\E[Z_t]\ge Z_0+\mu_{\mathrm{sep}}\,R_t;
\]
\item[(ii)] for every $t\ge 1$ and every $x>0$,
\[
\Prob\!\bigl(Z_t\le Z_0+\mu_{\mathrm{sep}}R_t-x\bigr)
\le
\exp\!\left(-\frac{x^2}{2tL^2}\right);
\]
\item[(iii)] equivalently, for every $\delta\in(0,1)$ and $t\ge 1$,
\[
\Prob\!\left(
q_t^\star(c^\star\mid H_t)\ge
\frac{1}{1+\exp\!\bigl(-Z_0-\mu_{\mathrm{sep}}R_t+L\sqrt{2t\log(1/\delta)}\bigr)}
\right)\ge 1-\delta.
\]
\end{itemize}
\end{theorem}

\begin{proof}
Let
\[
\Delta_t:=
\begin{cases}
Z_{t+1}-Z_t, & q_t^\star(c^\star\mid H_t)<1,\\
0, & q_t^\star(c^\star\mid H_t)=1.
\end{cases}
\]
On $E_t$, conditional on $\mathcal F_t$ and $A_{t+1}=a$, the class-level Bayes update gives
\[
\E[\Delta_t\mid\mathcal F_t,A_{t+1}=a]
=
\KL\!\bigl(\mu_{c^\star}(\cdot\mid H_t,a)\,\|\,\Qcal_t^{-c^\star}(\cdot\mid H_t,a)\bigr)
\ge 2\,B_t(c^\star,a;H_t),
\]
by Proposition~\ref{prop:chernoff-correspondence}\,\textup{(ii)}. On $E_t^c$, both $\Delta_t$ and $B_t(c^\star,a;H_t)$ are zero by definition. Therefore
\[
\E[\Delta_t\mid\mathcal F_t]\ge 2\,\E[B_t(c^\star,A_{t+1};H_t)\mid\mathcal F_t].
\]
On $E_t$ one has
\[
\E[B_t(c^\star,A_{t+1};H_t)\mid\mathcal F_t]
\ge
\eta(c^\star)\,\pi_{t+1}(S_t^\star\mid H_t)
\ge
\eta(c^\star)\,r_{t+1},
\]
hence
\[
\E[\Delta_t\mid\mathcal F_t]\ge \mu_{\mathrm{sep}}\,r_{t+1}
\]
on $E_t$.

For part \textup{(i)}, let $\tau:=\inf\{t\ge 0:q_t^\star(c^\star\mid H_t)=1\}$, with $\inf\varnothing=\infty$. If $\Prob(\tau\le t)>0$, then $Z_t=+\infty$ on that event and the lower bound is immediate. Otherwise $\tau>t$ almost surely, so $|\Delta_s|\le L$ and $Z_t=Z_0+\sum_{s=0}^{t-1}\Delta_s$ for all $s\le t-1$. Telescoping and the previous drift bound give
\[
\E[Z_t]\ge Z_0+\mu_{\mathrm{sep}}\sum_{s=0}^{t-1}r_{s+1}
=
Z_0+\mu_{\mathrm{sep}}R_t.
\]

For parts \textup{(ii)} and \textup{(iii)}, define the stopped increment
\[
\widetilde\Delta_s:=
\begin{cases}
\Delta_s, & \tau>s,\\
0, & \tau\le s,
\end{cases}
\]
the stopped log-odds
\[
\widehat Z_t:=Z_0+\sum_{s=0}^{t-1}\widetilde\Delta_s,
\]
and the compensated martingale
\[
M_t:=\sum_{s=0}^{t-1}\Bigl(\widetilde\Delta_s-\E[\widetilde\Delta_s\mid\mathcal F_s]\Bigr).
\]
The bounded-ratio hypothesis gives $|\widetilde\Delta_s|\le L$, so each martingale difference lies in an interval of width at most $2L$. Because $\{\tau>s\}\in\mathcal F_s$ and the drift bound above holds on $\{\tau>s\}$,
\[
\E[\widetilde\Delta_s\mid\mathcal F_s]
=
\one_{\{\tau>s\}}\E[\Delta_s\mid\mathcal F_s]
\ge
\mu_{\mathrm{sep}}\,r_{s+1}\one_{\{\tau>s\}}.
\]
Therefore
\[
\widehat Z_t
\ge
Z_0+\mu_{\mathrm{sep}}\sum_{s=0}^{t-1}r_{s+1}\one_{\{\tau>s\}}+M_t.
\]
If $Z_t\le Z_0+\mu_{\mathrm{sep}}R_t-x$, then necessarily $\tau>t$, because $Z_t=+\infty$ on $\{\tau\le t\}$. On $\{\tau>t\}$ one has $\widehat Z_t=Z_t$ and $\sum_{s=0}^{t-1}r_{s+1}\one_{\{\tau>s\}}=R_t$, so the displayed inequality implies
\[
M_t\le -x.
\]
Azuma--Hoeffding yields
\[
\Prob\!\bigl(Z_t\le Z_0+\mu_{\mathrm{sep}}R_t-x\bigr)
\le
\Prob(M_t\le -x)
\le
\exp\!\left(-\frac{x^2}{2tL^2}\right),
\]
which is part \textup{(ii)}. Part \textup{(iii)} follows from the monotonicity of $z\mapsto (1+e^{-z})^{-1}$ and the identity
\[
q_t^\star(c^\star\mid H_t)=\frac{1}{1+e^{-Z_t}}.
\]
\end{proof}

\begin{corollary}[Explicit finite-time recovery guarantee]
\label{cor:separating-support-finite-time}
Under the assumptions of Theorem~\ref{thm:separating-support-rate}, if for some $t\ge 1$ and some $\varepsilon,\delta\in(0,1)$,
\[
\mu_{\mathrm{sep}}R_t
\ge
\log\frac{1-\varepsilon}{\varepsilon}-Z_0+L\sqrt{2t\log(1/\delta)},
\]
then
\[
\Prob\!\bigl(q_t^\star([p^\star]\mid H_t)\ge 1-\varepsilon\bigr)\ge 1-\delta.
\]
\end{corollary}

\begin{proof}
The displayed hypothesis implies
\[
\frac{1}{1+\exp\!\bigl(-Z_0-\mu_{\mathrm{sep}}R_t+L\sqrt{2t\log(1/\delta)}\bigr)}\ge 1-\varepsilon,
\]
so the claim is exactly Theorem~\ref{thm:separating-support-rate}\,\textup{(iii)}.
\end{proof}

\begin{remark}[Canonical support-budget specializations]
\label{rem:canonical-support-budgets}
The three canonical realization families specialize Theorem~\ref{thm:separating-support-rate} by supplying concrete lower bounds on the separating-action mass:
\begin{itemize}
\item[(i)] for the canonical selector theorem, Proposition~\ref{prop:semantic-is-promotion-supporting} gives
\[
r_t\equiv \delta_0\,\bar w(c^\star),
\qquad
R_t=\delta_0\,\bar w(c^\star)\,t,
\]
so the drift constant becomes $\mu_0=2\eta(c^\star)\delta_0\bar w(c^\star)$ up to the selector's cap factor;
\item[(ii)] for the selector-free kernel theorem,
\[
r_t\equiv \delta_0\,\bar w(c^\star)\lambda_{\min},
\qquad
R_t=\delta_0\,\bar w(c^\star)\lambda_{\min}\,t,
\]
so the drift constant becomes $\mu_{0,\lambda}=2\eta(c^\star)\delta_0\bar w(c^\star)\lambda_{\min}$;
\item[(iii)] for full-support exploration floors and their self-referential specialization,
\[
r_t=\varepsilon_t\lambda_{\min},
\qquad
R_t=\lambda_{\min}\Gamma_t,
\qquad
\Gamma_t:=\sum_{i=1}^t\varepsilon_i.
\]
\end{itemize}
\end{remark}

\begin{theorem}[Canonical selector realization of semantic convergence]
\label{thm:semantic-convergence}
Assume:
\begin{itemize}
\item[(C1)] the prior is the universal interactive prior (Definition~\ref{def:universal-prior});
\item[(C2)] belief is the Bayes/Gibbs posterior (Lemma~\ref{lem:variational});
\item[(C3)] the semantic separation condition (Definition~\ref{def:sep-condition}) holds;
\item[(C4)] action follows the semantic action rule (Definition~\ref{def:semantic-rule}).
\end{itemize}
Then for every $p^\star\in\Pcal$,
\[
q_t^\star([p^\star]\mid H_t)\longrightarrow 1
\qquad
\text{almost surely.}
\]
\end{theorem}

\begin{proof}
Fix $p^\star\in\Pcal$ and write $c^\star:=[p^\star]$.
Define the smaller positive separation constants
\[
\hat\eta(c):=\frac{1}{2}\min\{1,\eta(c)\}
\qquad
(c\in\Ccal).
\]
Because Definition~\ref{def:sep-condition} remains true after replacing $\eta(c)$ by any smaller positive constant, the semantic separation condition also holds with the constants $\hat\eta(c)$. Fix
\[
S_t^\star:=\{a\in\Acal:B_t(c^\star,a;H_t)\ge \hat\eta(c^\star)\}.
\]
By Proposition~\ref{prop:semantic-is-promotion-supporting},
\[
\Prob(C_t=c^\star\mid H_t)\ge \delta_0\,\bar w(c^\star)
\qquad\text{for every }t.
\]
Whenever $C_t=c^\star$, the semantic-rule selector chooses an action $A_{t+1}=a_t^{c^\star}(H_t)$ satisfying
\[
B_t(c^\star,A_{t+1};H_t)\ge \hat\eta(c^\star),
\]
so
\[
\pi_{t+1}(S_t^\star\mid H_t)\ge \delta_0\,\bar w(c^\star)
\qquad\text{almost surely on }\{q_t^\star(c^\star\mid H_t)<1\}.
\]
Applying Theorem~\ref{thm:separating-support-convergence} with
\[
r_t\equiv \delta_0\,\bar w(c^\star)
\]
gives the claim.
\end{proof}

\begin{remark}[Functional-minimization reading]
\label{rem:functional-payoff}
Under the two-observer functional, (C1)--(C2) fix the belief-side minimizer of $\mathcal F_t$ via Lemma~\ref{lem:variational}; (C3) under the standing finite-action setup yields the kernel lift's uniform $\lambda$-mass floor via Proposition~\ref{prop:finite-action-kernel-separation}; and (K4) realizes the action-side minimizer of the kernel lift $\mathcal J_{t,\mathrm{ker}}^{\omega_{\mathrm{behav}}}$ directly via Proposition~\ref{prop:kernel-functional-minimizer}, with no external selector. Theorem~\ref{thm:kernel-semantic-convergence} is therefore the strongest exact variational realization of the master support theorem at the canonical pair $(\omega_B,\omega_A)=(\omega_{\mathrm{syn}},\omega_{\mathrm{behav}})$; Corollary~\ref{cor:compact-action-kernel} carries the same variational route to compact metric action spaces via Proposition~\ref{prop:kernel-functional-minimizer-compact} and Proposition~\ref{prop:compact-action-kernel-separation}. Theorem~\ref{thm:semantic-convergence} is the parallel canonical selector realization at the same pair, via $\mathcal J_t^{\omega_{\mathrm{behav}}}$ and Definition~\ref{def:semantic-rule}.
\end{remark}

\begin{remark}[Promotion-support generality]
\label{rem:promotion-support}
The proof invokes the semantic action rule only through Proposition~\ref{prop:semantic-is-promotion-supporting}. Any action rule satisfying Definition~\ref{def:promotion-supporting} and (C1)--(C3) obeys the same conclusion. Theorem~\ref{thm:semantic-convergence} thus identifies the class-against-complement rule as one instance of a broader family of observer-promoting action rules, with the near-miss construction of Theorem~\ref{thm:afe-near-miss} furnishing a failure of promotion support on the AFE side.
\end{remark}

\begin{theorem}[Kernel-functional semantic convergence]
\label{thm:kernel-semantic-convergence}
Fix a full-support reference distribution $\lambda\in\Delta(\Acal)$. Assume:
\begin{itemize}
\item[(C1)] the prior is the universal interactive prior (Definition~\ref{def:universal-prior});
\item[(C2)] belief is the Bayes/Gibbs posterior (Lemma~\ref{lem:variational});
\item[(C3)] the semantic separation condition (Definition~\ref{def:sep-condition}) holds;
\item[(K4)] action follows the kernel semantic action rule relative to $\lambda$ (Definition~\ref{def:kernel-semantic-rule}).
\end{itemize}
Then for every $p^\star\in\Pcal$,
\[
q_t^\star([p^\star]\mid H_t)\longrightarrow 1
\qquad
\text{almost surely.}
\]
\end{theorem}

\begin{proof}
Fix $p^\star\in\Pcal$ and write $c^\star:=[p^\star]$. Define
\[
S_t^\star:=\{a\in\Acal:B_t(c^\star,a;H_t)\ge \eta(c^\star)\},
\]
and let
\[
\pi_{t+1}(S\mid H_t):=\Prob(A_{t+1}\in S\mid H_t)
\]
denote the action marginal of the kernel semantic rule. Let $\lambda_{\min}:=\min_{a\in\Acal}\lambda(a)>0$. By Proposition~\ref{prop:finite-action-kernel-separation}, on $\{q_t^\star(c^\star\mid H_t)<1\}$ one has $\lambda(S_t^\star)\ge \lambda_{\min}$. Proposition~\ref{prop:kernel-promotion-support} therefore gives
\[
\pi_{t+1}(S_t^\star\mid H_t)\ge \delta_0\,\bar w(c^\star)\,\lambda_{\min}
\qquad\text{almost surely on }\{q_t^\star(c^\star\mid H_t)<1\},
\]
with $\delta_0:=e^{-\beta/\gamma}$. Applying Theorem~\ref{thm:separating-support-convergence} with
\[
r_t\equiv \delta_0\,\bar w(c^\star)\,\lambda_{\min}
\]
gives the claim.
\end{proof}

\begin{corollary}[Compact-action extension of the kernel theorem and rates]
\label{cor:compact-action-kernel}
If one relaxes the standing finite-action setup and instead lets $(\Acal,d_\Acal)$ be a compact metric space equipped with a full-support Borel probability measure $\lambda$, assume:
\begin{itemize}
\item[(i)] the measurability hypothesis of Proposition~\ref{prop:kernel-functional-minimizer-compact};
\item[(ii)] the continuity hypothesis of Proposition~\ref{prop:compact-action-kernel-separation};
\item[(iii)] the prior is the universal interactive prior and belief is the Bayes/Gibbs posterior;
\item[(iv)] action follows the exact minimizing Gibbs kernel
\[
(C_t,A_{t+1})\sim \kappa_{t,\mathrm{cts}}^{\mathrm G,\omega_{\mathrm{behav}}}(\cdot,\cdot\mid H_t).
\]
\end{itemize}
Fix a target program $p^\star\in\Pcal$ and write $c^\star:=[p^\star]$. Then:
\begin{itemize}
\item[(a)] under the semantic separation condition,
\[
q_t^\star([p^\star]\mid H_t)\longrightarrow 1
\qquad\text{almost surely;}
\]
\item[(b)] under the bounded-ratio regularity hypothesis of Proposition~\ref{prop:kernel-exp-rate}, the conclusion of that proposition and the kernel branch of Theorem~\ref{thm:exp-rate-concentration} hold with
\[
\eta_\lambda(c^\star):=\frac{\eta(c^\star)}{2},
\qquad
\rho_\lambda(c^\star):=m_\lambda(r_{c^\star}),
\]
where $r_{c^\star}>0$ is any radius satisfying $\varpi_{c^\star}(r_{c^\star})\le \eta(c^\star)/2$, and with
\[
\mu_{0,\lambda}^{\mathrm{cts}}
:=
2\,\eta_\lambda(c^\star)\,\delta_0\,\bar w(c^\star)\,\rho_\lambda(c^\star)
=
\eta(c^\star)\,\delta_0\,\bar w(c^\star)\,\rho_\lambda(c^\star).
\]
\end{itemize}
\end{corollary}

\begin{proof}
Let $c^\star:=[p^\star]$ and define
\[
S_t^\star:=\{a\in\Acal:B_t(c^\star,a;H_t)\ge \eta_\lambda(c^\star)\},
\qquad
\eta_\lambda(c^\star):=\frac{\eta(c^\star)}{2},
\qquad
\rho_\lambda(c^\star):=m_\lambda(r_{c^\star}).
\]
By Proposition~\ref{prop:compact-action-kernel-separation}, the semantic separation condition upgrades to the kernel semantic separation condition relative to $\lambda$ with the displayed constants. Let
\[
\pi_{t+1}^{\mathrm{cts}}(S\mid H_t):=\Prob(A_{t+1}\in S\mid H_t)
\]
denote the action marginal of the compact-action Gibbs kernel. Proposition~\ref{prop:kernel-promotion-support-compact} gives
\[
\pi_{t+1}^{\mathrm{cts}}(S_t^\star\mid H_t)\ge \delta_0\,\bar w(c^\star)\,\rho_\lambda(c^\star)
\qquad\text{almost surely on }\{q_t^\star(c^\star\mid H_t)<1\}.
\]
Applying Theorem~\ref{thm:separating-support-convergence} with
\[
r_t\equiv \delta_0\,\bar w(c^\star)\,\rho_\lambda(c^\star)
\]
gives part \textup{(a)}. Under the bounded-ratio hypothesis, applying Theorem~\ref{thm:separating-support-rate} with the same support floor and with $\eta(c^\star)$ replaced by $\eta_\lambda(c^\star)$ gives part \textup{(b)}, since
\[
2\,\eta_\lambda(c^\star)\,\delta_0\,\bar w(c^\star)\,\rho_\lambda(c^\star)
=
\eta(c^\star)\,\delta_0\,\bar w(c^\star)\,\rho_\lambda(c^\star)
=
\mu_{0,\lambda}^{\mathrm{cts}}.
\]
\end{proof}

\begin{corollary}[Finite-class deterministic specialization]
\label{cor:finite-maximin}
Fix a target program $p^\star\in\Pcal$ and write $c^\star:=[p^\star]$. Under (C1), (C2) and a finite semantic class set $\Ccal$ with $\eta>0$ such that whenever $q_t^\star(c^\star\mid H_t)<1$, some action $a$ satisfies $B_t(c,a;H_t)\ge\eta$ for every unresolved class $c$, define
\[
U_t(a;H_t):=
\begin{cases}
\min_{c:\,0<q_t^\star(c\mid H_t)<1}B_t(c,a;H_t), & q_t^\star(c^\star\mid H_t)<1,\\
0, & q_t^\star(c^\star\mid H_t)=1,
\end{cases}
\]
and whenever $q_t^\star(c^\star\mid H_t)<1$ let $A_{t+1}$ satisfy $U_t(A_{t+1};H_t)\ge\tfrac{1}{2}\min\{1,\sup_a U_t(a;H_t)\}$, choosing $A_{t+1}$ arbitrarily when $q_t^\star(c^\star\mid H_t)=1$. Then $q_t^\star(c^\star\mid H_t)\to 1$ almost surely.
\end{corollary}

\begin{proof}
Set
\[
\hat\eta:=\frac{1}{2}\min\{1,\eta\}>0,
\qquad
S_t^\star:=\{a\in\Acal:B_t(c^\star,a;H_t)\ge \hat\eta\}.
\]
Whenever the posterior is not concentrated, the hypothesis gives $\sup_a U_t(a;H_t)\ge\eta$, so the chosen action satisfies
\[
A_{t+1}\in S_t^\star.
\]
Thus
\[
\pi_{t+1}(S_t^\star\mid H_t)=1
\qquad\text{almost surely on }\{q_t^\star(c^\star\mid H_t)<1\}.
\]
Applying Theorem~\ref{thm:separating-support-convergence} with $r_t\equiv 1$ gives the claim.
\end{proof}

\begin{theorem}[Zero separating support makes semantic recovery impossible]
\label{cor:support-necessary}
Fix a policy $\pi$. Suppose there exist $p^\star\in\Pcal$, a history $h_{t_0}$ with $\Prob_{p^\star,\pi}(H_{t_0}=h_{t_0})>0$, an action $b\in\Acal$ with $\pi_{t_0+1}(b\mid h_{t_0})=0$, and an observation $y\in\Ocal$ with $\mu_{p^\star}(y\mid h_{t_0},b)<1$. Then no Bayesian scheme whose prior assigns positive mass to every computable program can satisfy
\[
q_t([p]\mid H_t)\longrightarrow 1
\qquad\text{almost surely under }(p,\pi)
\]
for every $p\in\Pcal$.
\end{theorem}

\begin{proof}
Theorem~\ref{thm:policy-gap} produces $p'\in R_\pi(p^\star)\setminus[p^\star]$. For every horizon $T$ and every history $h_T$ with $\Prob_{p^\star,\pi}(H_T=h_T)>0$, every program in $[p^\star]$ has the realized likelihood of $p^\star$ and every program in $[p']$ has the realized likelihood of $p'$; because $p'\in R_\pi(p^\star)$, these class likelihoods are equal under $\pi$. Hence
\[
\frac{q_T([p']\mid h_T)}{q_T([p^\star]\mid h_T)}
=
\frac{\bar w([p'])}{\bar w([p^\star])}
>0.
\]
Therefore
\[
q_T([p^\star]\mid h_T)\le \frac{1}{1+\bar w([p'])/\bar w([p^\star])}<1
\]
at every horizon on the true-history support, so concentration on $[p^\star]$ fails under $(p^\star,\pi)$. Therefore no such uniform theorem over all $p\in\Pcal$ can hold for the fixed policy $\pi$.
\end{proof}

\begin{theorem}[Summable separating support is insufficient]
\label{thm:summable-support-insufficient}
Let $(r_t)_{t\ge 1}\subset[0,1)$ satisfy
\[
\sum_{t=1}^{\infty}r_t<\infty.
\]
Then there exist a finite action alphabet $\Acal$, a finite observation alphabet $\Ocal$, two computable environments $p^\star,p'\in\Pcal$, a proper prior supported on $\{p^\star,p'\}$ with
\[
0<q_0([p^\star])<1,
\]
and a history-independent policy $\pi$ such that:
\begin{itemize}
\item[(i)] the semantic classes are singletons and there is an action $a^{\mathrm{sep}}$ that deterministically separates $[p^\star]$ from $[p']$;
\item[(ii)] for every history $h_t$ and every $t\ge 0$,
\[
\pi_{t+1}(a^{\mathrm{sep}}\mid h_t)=r_{t+1},
\qquad
\pi_{t+1}(a^{\mathrm{ref}}\mid h_t)=1-r_{t+1};
\]
\item[(iii)] with positive probability under $(p^\star,\pi)$, the posterior on $[p^\star]$ never moves:
\[
q_t([p^\star]\mid H_t)=q_0([p^\star])
\qquad\text{for every }t\ge 0.
\]
\end{itemize}
In particular, divergent cumulative separating support in Theorem~\ref{thm:separating-support-convergence} cannot be weakened to summable support.
\end{theorem}

\begin{proof}
Let
\[
\Acal:=\{a^{\mathrm{ref}},a^{\mathrm{sep}}\},
\qquad
\Ocal:=\{0,1\}.
\]
Define two history-independent computable environments by
\[
\mu_{p^\star}(\cdot\mid h,a^{\mathrm{ref}})=\delta_0,
\qquad
\mu_{p'}(\cdot\mid h,a^{\mathrm{ref}})=\delta_0,
\]
and
\[
\mu_{p^\star}(\cdot\mid h,a^{\mathrm{sep}})=\delta_1,
\qquad
\mu_{p'}(\cdot\mid h,a^{\mathrm{sep}})=\delta_0.
\]
Thus the semantic classes are singletons and the action $a^{\mathrm{sep}}$ deterministically separates $[p^\star]=\{p^\star\}$ from $[p']=\{p'\}$.

Fix any prior supported on $\{p^\star,p'\}$ with
\[
0<w(p^\star)<1,
\qquad
w(p')=1-w(p^\star).
\]
Define the history-independent policy $\pi$ by
\[
\pi_{t+1}(a^{\mathrm{sep}}\mid h_t)=r_{t+1},
\qquad
\pi_{t+1}(a^{\mathrm{ref}}\mid h_t)=1-r_{t+1}
\]
for every $t\ge 0$ and every history $h_t$.

Consider the event
\[
E:=\{A_t=a^{\mathrm{ref}}\text{ for every }t\ge 1\}.
\]
Since $\sum_t r_t<\infty$ and each $r_t<1$, the infinite product
\[
\Prob_{p^\star,\pi}(E)=\prod_{t=1}^{\infty}(1-r_t)
\]
converges to a strictly positive number. On $E$, every realized observation is $0$, and the full realized history law coincides under $p^\star$ and $p'$, because the two environments agree on $a^{\mathrm{ref}}$. Therefore for every $t\ge 0$ and every realized history $h_t$ on $E$,
\[
\mu_{p^\star}(o_{1:t}\givena a_{1:t})=\mu_{p'}(o_{1:t}\givena a_{1:t}),
\]
so Bayes' rule gives
\[
q_t([p^\star]\mid h_t)
=
\frac{w(p^\star)}{w(p^\star)+w(p')}
=
w(p^\star)
=
q_0([p^\star]).
\]
Thus the posterior on $[p^\star]$ stays equal to its initial class prior forever on the positive-probability event $E$, so it does not converge almost surely to one.
\end{proof}


\subsection{An expectation-form rate under bounded log-likelihood ratios}
\label{sec:exp-rate}

Theorem~\ref{thm:separating-support-rate} already gives an explicit finite-time support-based recovery guarantee under bounded log-likelihood ratios. This subsection records the corresponding selector and kernel realizations, whose constants specialize the master rate theorem to the two canonical action constructions. Throughout this subsection, fix a target program $p^\star\in\Pcal$ and write $c^\star:=[p^\star]$. Define the true-class log-odds
\[
Z_t\;:=\;\log\frac{q_t^\star(c^\star\mid H_t)}{1-q_t^\star(c^\star\mid H_t)},
\]
with the convention $Z_t=+\infty$ when $q_t^\star(c^\star\mid H_t)=1$. If $q_0(c^\star)=1$, then Bayes keeps $q_t^\star(c^\star\mid H_t)=1$ for all $t$ and every expectation and high-probability lower bound below is immediate. Otherwise the universal prior gives $0<q_0(c^\star)<1$, so $Z_0$ is finite. Write $\Delta_t:=Z_{t+1}-Z_t$ on the event $\{q_t^\star(c^\star\mid H_t)<1\}$ and $\mathcal F_t:=\sigma(H_t,C_0,\dots,C_{t-1})$ as in the proof of Theorem~\ref{thm:semantic-convergence}.

\begin{lemma}[Canonical selector one-step drift]
\label{lem:one-step-drift}
Under \textup{(C1)--(C4)} and on the event $\{q_t^\star(c^\star\mid H_t)<1\}$,
\[
\E[\Delta_t\mid\mathcal F_t]
\;\ge\;
2\,\E[B_t(c^\star,A_{t+1};H_t)\mid\mathcal F_t]
\;\ge\;
2\hat\eta\,\delta_0\,\bar w(c^\star),
\qquad
\hat\eta:=\tfrac{1}{2}\min\{1,\eta(c^\star)\}>0.
\]
\end{lemma}

\begin{proof}
Conditional on $\mathcal F_t$ and $A_{t+1}=a$, the class-level Bayes update gives
\[
\E[\Delta_t\mid\mathcal F_t,A_{t+1}=a]
=
\E_{o\sim\mu_{c^\star}(\cdot\mid H_t,a)}\!\left[
\log\frac{\mu_{c^\star}(o\mid H_t,a)}{\Qcal_t^{-c^\star}(o\mid H_t,a)}
\right]
=\KL\!\bigl(\mu_{c^\star}(\cdot\mid H_t,a)\,\|\,\Qcal_t^{-c^\star}(\cdot\mid H_t,a)\bigr)
\]
by Lemma~\ref{lem:odds-identity} applied to $c=c^\star$ together with the identity $\log(\alpha/(1-\alpha))=Z_t(\alpha)$. Proposition~\ref{prop:chernoff-correspondence}\,(ii) gives $\KL\ge 2B_t(c^\star,a;H_t)$, so
\[
\E[\Delta_t\mid\mathcal F_t,A_{t+1}=a]\ge 2\,B_t(c^\star,a;H_t).
\]
Taking the outer expectation over $A_{t+1}\mid\mathcal F_t$,
\[
\E[\Delta_t\mid\mathcal F_t]\ge 2\,\E[B_t(c^\star,A_{t+1};H_t)\mid\mathcal F_t].
\]
Proposition~\ref{prop:semantic-is-promotion-supporting} gives $\Prob(C_t=c^\star\mid\mathcal F_t)\ge\delta_0\bar w(c^\star)$. On $\{C_t=c^\star\}$, the semantic-rule selector satisfies $B_t(c^\star,A_{t+1};H_t)\ge\hat\eta$ by the semantic separation condition (C3), so
\[
\E[B_t(c^\star,A_{t+1};H_t)\mid\mathcal F_t]\ge\hat\eta\,\delta_0\,\bar w(c^\star).
\]
Chaining the two inequalities gives the claim.
\end{proof}

\begin{proposition}[Canonical selector rate, expectation form]
\label{prop:exp-rate}
Assume \textup{(C1)--(C4)} and, in addition, that there exists $m\in(0,1)$ such that whenever $0<q_t^\star(c^\star\mid H_t)<1$,
\[
m\le
\frac{\mu_{c^\star}(o\mid H_t,A_{t+1})}{\Qcal_t^{-c^\star}(o\mid H_t,A_{t+1})}
\le \frac{1}{m}
\]
for every $o$ with $\mu_{c^\star}(o\mid H_t,A_{t+1})>0$. Then for every $t\ge 0$,
\[
\E[Z_t]\;\ge\;Z_0+\mu_0\,t,
\qquad
\mu_0\;:=\;2\hat\eta\,\delta_0\,\bar w(c^\star)\;>\;0.
\]
In particular, $\E[Z_t]\to\infty$ linearly in $t$, and the expected true-class posterior satisfies
\[
\E\!\left[\frac{q_t^\star(c^\star\mid H_t)}{1-q_t^\star(c^\star\mid H_t)}\right]\ge e^{Z_0+\mu_0 t}
\]
by Jensen's inequality on $e^{Z_t}$.
\end{proposition}

\begin{proof}
Let $\tau:=\inf\{t\ge 0:q_t^\star(c^\star\mid H_t)=1\}$, with the convention $\inf\varnothing=\infty$. If $\Prob(\tau\le t)>0$, then $Z_t=+\infty$ on that event and the lower bound for $\E[Z_t]$ is immediate. It therefore remains to consider the case $\tau>t$ almost surely.

On $\{\tau>t\}$ one has $0<q_s^\star(c^\star\mid H_s)<1$ for every $s\le t$, and the displayed hypothesis implies
\[
\Delta_s=\log\frac{\mu_{c^\star}(O_{s+1}\mid H_s,A_{s+1})}{\Qcal_s^{-c^\star}(O_{s+1}\mid H_s,A_{s+1})}
\]
is well-defined with $|\Delta_s|\le \log(1/m)$. Hence $Z_s$ and $\Delta_s$ are integrable for $s\le t$. Lemma~\ref{lem:one-step-drift} gives $\E[\Delta_s\mid\mathcal F_s]\ge\mu_0$, hence $\E[\Delta_s]\ge\mu_0$. Since $Z_t=Z_0+\sum_{s=0}^{t-1}\Delta_s$, telescoping yields $\E[Z_t]\ge Z_0+\mu_0 t$. Jensen applied to the convex function $x\mapsto e^x$ gives the exponential lower bound on the expected odds ratio. If instead $\Prob(\tau\le t)>0$, the same inequality is trivial because the left-hand side is $+\infty$.
\end{proof}

\begin{lemma}[Kernel one-step drift]
\label{lem:one-step-drift-kernel}
Under \textup{(C1)}, \textup{(C2)}, \textup{(C3)}, and \textup{(K4)} and on the event $\{q_t^\star(c^\star\mid H_t)<1\}$,
\[
\E[\Delta_t\mid\mathcal F_t]
\;\ge\;
2\,\E[B_t(c^\star,A_{t+1};H_t)\mid\mathcal F_t]
\;\ge\;
2\,\eta(c^\star)\,\delta_0\,\bar w(c^\star)\,\lambda_{\min},
\qquad
\lambda_{\min}:=\min_{a\in\Acal}\lambda(a)>0.
\]
\end{lemma}

\begin{proof}
Conditional on $\mathcal F_t$ and $A_{t+1}=a$, the same class-level Bayes update used in Lemma~\ref{lem:one-step-drift} gives
\[
\E[\Delta_t\mid\mathcal F_t,A_{t+1}=a]
=
\KL\!\bigl(\mu_{c^\star}(\cdot\mid H_t,a)\,\|\,\Qcal_t^{-c^\star}(\cdot\mid H_t,a)\bigr)
\ge 2\,B_t(c^\star,a;H_t)
\]
by Proposition~\ref{prop:chernoff-correspondence}\,(ii). Taking the outer expectation over $A_{t+1}\mid\mathcal F_t$ yields
\[
\E[\Delta_t\mid\mathcal F_t]\ge 2\,\E[B_t(c^\star,A_{t+1};H_t)\mid\mathcal F_t].
\]
Let
\[
S_t^\star:=\{a\in\Acal:B_t(c^\star,a;H_t)\ge \eta(c^\star)\}.
\]
By Proposition~\ref{prop:finite-action-kernel-separation}, $\lambda(S_t^\star)\ge \lambda_{\min}$ on $\{q_t^\star(c^\star\mid H_t)<1\}$. Proposition~\ref{prop:kernel-promotion-support} therefore gives
\[
\Prob(C_t=c^\star,\ A_{t+1}\in S_t^\star\mid\mathcal F_t)\ge \delta_0\,\bar w(c^\star)\,\lambda_{\min}.
\]
On $\{C_t=c^\star,\ A_{t+1}\in S_t^\star\}$ one has $B_t(c^\star,A_{t+1};H_t)\ge \eta(c^\star)$, hence
\[
\E[B_t(c^\star,A_{t+1};H_t)\mid\mathcal F_t]
\ge
\eta(c^\star)\,\delta_0\,\bar w(c^\star)\,\lambda_{\min}.
\]
Chaining the two inequalities gives the claim.
\end{proof}

\begin{proposition}[Kernel-functional rate, expectation form]
\label{prop:kernel-exp-rate}
Assume \textup{(C1)}, \textup{(C2)}, \textup{(C3)}, and \textup{(K4)} and, in addition, that there exists $m\in(0,1)$ such that whenever $0<q_t^\star(c^\star\mid H_t)<1$,
\[
m\le
\frac{\mu_{c^\star}(o\mid H_t,A_{t+1})}{\Qcal_t^{-c^\star}(o\mid H_t,A_{t+1})}
\le \frac{1}{m}
\]
for every $o$ with $\mu_{c^\star}(o\mid H_t,A_{t+1})>0$. Then for every $t\ge 0$,
\[
\E[Z_t]\;\ge\;Z_0+\mu_{0,\lambda}\,t,
\qquad
\mu_{0,\lambda}\;:=\;2\,\eta(c^\star)\,\delta_0\,\bar w(c^\star)\,\lambda_{\min}\;>\;0,
\qquad
\lambda_{\min}:=\min_{a\in\Acal}\lambda(a)>0.
\]
In particular, $\E[Z_t]\to\infty$ linearly in $t$, and
\[
\E\!\left[\frac{q_t^\star(c^\star\mid H_t)}{1-q_t^\star(c^\star\mid H_t)}\right]\ge e^{Z_0+\mu_{0,\lambda} t}.
\]
\end{proposition}

\begin{proof}
Let $\tau:=\inf\{t\ge 0:q_t^\star(c^\star\mid H_t)=1\}$, with the convention $\inf\varnothing=\infty$. If $\Prob(\tau\le t)>0$, then $Z_t=+\infty$ on that event and the lower bound is immediate. It therefore remains to consider the case $\tau>t$ almost surely.

On $\{\tau>t\}$ one has $0<q_s^\star(c^\star\mid H_s)<1$ for every $s\le t$, and the displayed hypothesis implies
\[
\Delta_s=\log\frac{\mu_{c^\star}(O_{s+1}\mid H_s,A_{s+1})}{\Qcal_s^{-c^\star}(O_{s+1}\mid H_s,A_{s+1})}
\]
is well-defined with $|\Delta_s|\le \log(1/m)$. Hence $Z_s$ and $\Delta_s$ are integrable for $s\le t$. Lemma~\ref{lem:one-step-drift-kernel} gives $\E[\Delta_s\mid\mathcal F_s]\ge\mu_{0,\lambda}$, hence $\E[\Delta_s]\ge\mu_{0,\lambda}$. Since $Z_t=Z_0+\sum_{s=0}^{t-1}\Delta_s$, telescoping yields $\E[Z_t]\ge Z_0+\mu_{0,\lambda}t$. Jensen applied to the convex function $x\mapsto e^x$ gives the exponential lower bound on the expected odds ratio. If instead $\Prob(\tau\le t)>0$, the same inequality is trivial because the left-hand side is $+\infty$.
\end{proof}

\begin{theorem}[Finite-time high-probability rate bounds]
\label{thm:exp-rate-concentration}
Let
\[
L:=\log(1/m).
\]
\begin{itemize}
\item[(i)] Under the assumptions of Proposition~\ref{prop:exp-rate}, for every $t\ge 1$ and every $x>0$,
\[
\Prob\!\bigl(Z_t\le Z_0+\mu_0 t-x\bigr)
\;\le\;
\exp\!\left(-\frac{x^2}{2tL^2}\right).
\]
Equivalently, for every $\delta\in(0,1)$,
\[
\Prob\!\left(
q_t^\star(c^\star\mid H_t)\ge
\frac{1}{1+\exp\!\bigl(-Z_0-\mu_0 t+L\sqrt{2t\log(1/\delta)}\bigr)}
\right)\ge 1-\delta.
\]
\item[(ii)] Under the assumptions of Proposition~\ref{prop:kernel-exp-rate}, for every $t\ge 1$ and every $x>0$,
\[
\Prob\!\bigl(Z_t\le Z_0+\mu_{0,\lambda} t-x\bigr)
\;\le\;
\exp\!\left(-\frac{x^2}{2tL^2}\right).
\]
Equivalently, for every $\delta\in(0,1)$,
\[
\Prob\!\left(
q_t^\star(c^\star\mid H_t)\ge
\frac{1}{1+\exp\!\bigl(-Z_0-\mu_{0,\lambda} t+L\sqrt{2t\log(1/\delta)}\bigr)}
\right)\ge 1-\delta.
\]
\end{itemize}
\end{theorem}

\begin{proof}
We prove \textup{(i)}; part \textup{(ii)} is identical with $\mu_0$ replaced by $\mu_{0,\lambda}$ and Lemma~\ref{lem:one-step-drift} replaced by Lemma~\ref{lem:one-step-drift-kernel}.

Let
\[
\tau:=\inf\{t\ge 0:q_t^\star(c^\star\mid H_t)=1\},
\]
with the convention $\inf\varnothing=\infty$, and define the stopped increment
\[
\widetilde\Delta_s:=
\begin{cases}
\Delta_s, & \tau>s,\\
0, & \tau\le s.
\end{cases}
\]
By the bounded-ratio hypothesis of Proposition~\ref{prop:exp-rate}, on $\{\tau>s\}$ one has $|\Delta_s|\le L$, hence $|\widetilde\Delta_s|\le L$ almost surely. Define
\[
\widehat Z_t:=Z_0+\sum_{s=0}^{t-1}\widetilde\Delta_s
\]
and the compensated process
\[
M_t:=\sum_{s=0}^{t-1}\Bigl(\widetilde\Delta_s-\E[\widetilde\Delta_s\mid\mathcal F_s]\Bigr).
\]
Then $(M_t)_{t\ge 0}$ is a martingale. Since $\widetilde\Delta_s\in[-L,L]$ almost surely, each martingale difference
\[
M_{s+1}-M_s=\widetilde\Delta_s-\E[\widetilde\Delta_s\mid\mathcal F_s]
\]
lies in an interval of width at most $2L$.

Because $\{\tau>s\}\in\mathcal F_s$,
\[
\E[\widetilde\Delta_s\mid\mathcal F_s]
=
\one_{\{\tau>s\}}\E[\Delta_s\mid\mathcal F_s]
\ge
\mu_0\,\one_{\{\tau>s\}}
\]
by Lemma~\ref{lem:one-step-drift}. Therefore
\[
\widehat Z_t
=
Z_0+\sum_{s=0}^{t-1}\E[\widetilde\Delta_s\mid\mathcal F_s]+M_t
\ge
Z_0+\mu_0\sum_{s=0}^{t-1}\one_{\{\tau>s\}}+M_t.
\]
If $Z_t\le Z_0+\mu_0 t-x$, then necessarily $\tau>t$, because $Z_t=+\infty$ on $\{\tau\le t\}$. On $\{\tau>t\}$ one has $\widehat Z_t=Z_t$ and $\sum_{s=0}^{t-1}\one_{\{\tau>s\}}=t$, so the displayed inequality implies
\[
M_t\le -x.
\]
Azuma--Hoeffding therefore yields
\[
\Prob\!\bigl(Z_t\le Z_0+\mu_0 t-x\bigr)
\le
\Prob(M_t\le -x)
\le
\exp\!\left(-\frac{x^2}{2tL^2}\right),
\]
proving the first display in \textup{(i)}.

For the posterior form, note that $z\mapsto (1+e^{-z})^{-1}$ is increasing and
\[
q_t^\star(c^\star\mid H_t)=\frac{1}{1+e^{-Z_t}}.
\]
Set $x:=L\sqrt{2t\log(1/\delta)}$ in the first bound and rewrite the complementary event in terms of this logistic transform.
\end{proof}

\begin{corollary}[Goal-dialed convergence]
\label{cor:goal-dialed-payoff}
Fix an observer $\omega_A$ with $\omega_A\le\omega_{\mathrm{behav}}$ and a target program $p^\star\in\Pcal$. Under the observer-indexed conditions \textup{(C1$^{\omega_A}$)--(C4$^{\omega_A}$)} of Proposition~\ref{prop:goal-dialed},
\[
q_t^\star([p^\star]_{\omega_A}\mid H_t)\;\longrightarrow\;1
\qquad\text{almost surely,}
\]
and, under the corresponding observer-indexed bounded-ratio hypothesis, Proposition~\ref{prop:exp-rate} transfers with $c^\star$ replaced by $[p^\star]_{\omega_A}$, $\eta(c^\star)$ replaced by $\eta^{\omega_A}([p^\star]_{\omega_A})$, and $\bar w(c^\star)$ replaced by $\bar w^{\omega_A}([p^\star]_{\omega_A})$: the expected log-odds of $[p^\star]_{\omega_A}$ grow linearly at rate $\mu_0^{\omega_A}:=2\hat\eta^{\omega_A}\delta_0\bar w^{\omega_A}([p^\star]_{\omega_A})>0$.
\end{corollary}

\begin{proof}
Proposition~\ref{prop:goal-dialed} gives the almost-sure statement. Its proof sketch lifts Ingredients~1--5 of Theorem~\ref{thm:semantic-convergence} to $\Ccal^{\omega_A}$ in place of $\Ccal$; the one-step-drift argument of Lemma~\ref{lem:one-step-drift} uses only Lemma~\ref{lem:odds-identity} and Proposition~\ref{prop:chernoff-correspondence}, both observer-agnostic. Under the corresponding observer-indexed bounded-ratio hypothesis, the same pre-absorption telescope used in Proposition~\ref{prop:exp-rate} lifts with the $\omega_A$-indexed constants.
\end{proof}

Propositions~\ref{prop:exp-rate} and~\ref{prop:kernel-exp-rate}, together with Theorem~\ref{thm:exp-rate-concentration}, show how the master support-rate theorem specializes to the canonical observer choice $\omega_A=\omega_{\mathrm{behav}}$ in both the selector realization and the selector-free kernel lift. Corollary~\ref{cor:compact-action-kernel} shows that the same kernel-rate conclusions persist for compact metric action spaces under the continuity and full-support hypotheses of Proposition~\ref{prop:compact-action-kernel-separation}. Corollary~\ref{cor:goal-dialed-payoff} then reads the two-observer functional as an observer dial with a quantitative selector rate: coarser $\omega_A$'s alter the rate through both the larger pushforward mass $\bar w^{\omega_A}([p^\star]_{\omega_A})$ and the observer-indexed separation constant $\hat\eta^{\omega_A}$, while any $\omega_A>\omega_{\mathrm{behav}}$ is ruled out by Proposition~\ref{prop:action-cap}.

\section{Self-Referential Reduction of the Computability Gap}
\label{sec:self-ref}

Theorem~\ref{thm:semantic-convergence} and Proposition~\ref{prop:goal-dialed} drive the posterior onto $\omega_A$-classes for \emph{any} fixed action observer $\omega_A\le\omega_{\mathrm{behav}}$. The formally-stated rule is class-against-complement at $\omega_A=\omega_{\mathrm{behav}}$ (Definition~\ref{def:semantic-rule}), and the payoff class is $[p^\star]=[p^\star]_{\mathrm{behav}}$. A learner running the rule, however, cannot directly compute the $\omega_{\mathrm{behav}}$-partition: two programs are behaviorally equivalent when $\mu_{p_1}=\mu_{p_2}$ as interactive laws on the \emph{whole} policy family, and no finite history certifies this. The rule as written requires an oracle for the partition it targets.

This section isolates the tension at a finitary level rather than claiming an unconditional closure of it. The proxy object is the \emph{finite-time policy partition} of $\Pcal$ induced by the length-$t$ marginal of $\Prob_{p,\pi}$. The results below show that these partitions refine monotonically and that the self-referential rule retains the promotion-supporting Gibbs floor at the running partition. They then reduce semantic convergence of the self-referential rule to an \emph{eventual-isolation} hypothesis: if the running partition eventually isolates the true semantic class $[p^\star]$, then the main convergence argument of Section~\ref{sec:main} transfers. The remaining gap is exactly the proof of that eventual isolation.

\subsection{The finite-time policy observer}
\label{sec:finite-time-policy-observer}

\begin{definition}[Finite-time policy observer]
\label{def:finite-time-policy-observer}
Fix a policy $\pi$ and $t\ge 0$. The \emph{finite-time policy observer} $\omega_\pi^t$ is the observer of Definition~\ref{def:observer} with valuation
\[
V_{\omega_\pi^t}(p)\;:=\;\Prob_{p,\pi}(H_t\in\cdot),
\]
the law of the length-$t$ history under $(p,\pi)$. Write $[p^\star]_{\omega_\pi^t}$ for the $\omega_\pi^t$-indistinguishability class of $p^\star$.
\end{definition}

\begin{lemma}[Monotone refinement toward $\omega_\pi$]
\label{lem:monotone-refinement}
For every policy $\pi$ and every $t\ge 0$,
\[
\omega_\pi^t\;\le\;\omega_\pi^{t+1}\;\le\;\omega_\pi\;\le\;\omega_{\mathrm{behav}}\;\le\;\omega_{\mathrm{syn}},
\]
and the chain of partitions of $\Pcal$ induced by the $\omega_\pi^t$ refines to the partition induced by $\omega_\pi$:
\[
p_1\sim_{\omega_\pi}p_2\quad\Longleftrightarrow\quad p_1\sim_{\omega_\pi^t}p_2\text{ for every }t\ge 0.
\]
\end{lemma}

\begin{proof}
The length-$t$ marginal $\Prob_{p,\pi}(H_t\in\cdot)$ is the marginal on the first $t$ coordinates of $\Prob_{p,\pi}$; it is therefore a measurable function of the length-$(t+1)$ marginal, giving $\omega_\pi^t\le\omega_\pi^{t+1}$. Each length-$t$ marginal is in turn a measurable function of the full law $\Prob_{p,\pi}$, so $\omega_\pi^t\le\omega_\pi$. The nesting $\omega_\pi\le\omega_{\mathrm{behav}}\le\omega_{\mathrm{syn}}$ is the chain of Section~\ref{sec:hierarchy}. For the equivalence: $\sim_{\omega_\pi}$ is equality of the laws $\Prob_{p,\pi}$ on the full history space; by Kolmogorov extension, two laws on a countable product space are equal iff their length-$t$ marginals agree for every $t$.
\end{proof}

The finite-time policy observer is finitary in a sense the $\omega_{\mathrm{behav}}$-partition is not: for fixed $(\pi,t)$, checking whether $p_1\sim_{\omega_\pi^t}p_2$ is equality of two probability measures on the countable set of histories of length $t$, equivalently equality on all length-$t$ cylinder events. This replaces full interactive-law equality by a fixed-horizon marginal comparison. Under the standing finite-action/countable-observation setup, however, we do not claim exact decidability of this comparison.

\subsection{The self-referential semantic rule}
\label{sec:self-ref-rule}

\begin{definition}[Self-referential semantic rule]
\label{def:self-ref-rule}
Define the policy $\pi^{\mathrm{sr}}$ inductively on $t$. At $t=0$, let $\pi^{\mathrm{sr}}_1$ be the class-against-complement rule of Definition~\ref{def:semantic-rule} applied at the trivial partition $\{\Pcal\}$ (equivalently, at $\omega_{\pi^{\mathrm{sr}}}^0$). For $t\ge 1$, having defined $\pi^{\mathrm{sr}}_1,\dots,\pi^{\mathrm{sr}}_t$, let $\omega_{\pi^{\mathrm{sr}}}^t$ denote the finite-time policy observer of Definition~\ref{def:finite-time-policy-observer} computed from the prefix $(\pi^{\mathrm{sr}}_1,\dots,\pi^{\mathrm{sr}}_t)$, and set $\pi^{\mathrm{sr}}_{t+1}$ to be the class-against-complement rule of Definition~\ref{def:semantic-rule} with $\Ccal$ replaced throughout by the $\omega_{\pi^{\mathrm{sr}}}^t$-partition $\Ccal^{\omega_{\pi^{\mathrm{sr}}}^t}$.
\end{definition}

The inductive definition is well-posed as a formal recursion: $\omega_{\pi^{\mathrm{sr}}}^t$ depends only on $(\pi^{\mathrm{sr}}_1,\dots,\pi^{\mathrm{sr}}_t)$ and on the corresponding length-$t$ marginal laws. No fixed-point computation is required at the level of definition. The Gibbs class-targeting distribution $\nu_t^{\mathrm G}$ and the class-separating selector $a_t^c(h_t)$ are inherited from Definition~\ref{def:semantic-rule} applied at the running partition.

\begin{remark}[Self-referentiality vs.\ bootstrapping]
\label{rem:self-reference-scope}
The rule is self-referential in the narrow sense that the partition at step $t+1$ is derived from the rule's own behavior over $1,\dots,t$: the rule reads its own induced policy observer. This is not an open fixed-point --- the map $(\pi^{\mathrm{sr}}_1,\dots,\pi^{\mathrm{sr}}_t)\mapsto\pi^{\mathrm{sr}}_{t+1}$ is a deterministic measurable function. Paired with Lemma~\ref{lem:monotone-refinement}, the rule's partition sequence $\Ccal^{\omega_{\pi^{\mathrm{sr}}}^t}$ is nondecreasing in refinement, so the rule becomes more discriminating over time by construction.
\end{remark}

\subsection{Observer promotion under the self-referential rule}
\label{sec:observer-promotion-sr}

The key missing step is promotion of the running partition to the behavioral one. Lemma~\ref{lem:exploration-reachability} supplies only the action-support half of that story. The next proposition records the exact eventual-isolation criterion needed for promotion.

\begin{lemma}[Promotion-floor reachability]
\label{lem:exploration-reachability}
Fix a target program $p^\star\in\Pcal$. Under the self-referential rule $\pi^{\mathrm{sr}}$ of Definition~\ref{def:self-ref-rule}, for every $t\ge 0$, every class $c\in\Ccal^{\omega_{\pi^{\mathrm{sr}}}^t}$, and every history $h_t$ with $\Prob_{p^\star,\pi^{\mathrm{sr}}}(H_t=h_t)>0$,
\[
\Prob\bigl(A_{t+1}=a_t^c(h_t)\,\big|\,H_t=h_t\bigr)\;\ge\;\delta_0\,\bar w^{\omega_{\pi^{\mathrm{sr}}}^t}(c)\;>\;0.
\]
In particular, for every history $h_t$ with $\Prob_{p^\star,\pi^{\mathrm{sr}}}(H_t=h_t)>0$ and every class $c\in\Ccal^{\omega_{\pi^{\mathrm{sr}}}^t}$,
\[
\Prob_{p^\star,\pi^{\mathrm{sr}}}\bigl(H_t=h_t,\ A_{t+1}=a_t^c(h_t)\bigr)>0.
\]
Equivalently, the set of history-action pairs reachable under $\pi^{\mathrm{sr}}$ with positive probability is exactly
\[
\Bigl\{(h_t,a)\,:\,\Prob_{p^\star,\pi^{\mathrm{sr}}}(H_t=h_t)>0,\ 
a\in\{a_t^c(h_t):c\in\Ccal^{\omega_{\pi^{\mathrm{sr}}}^t}\}\Bigr\}.
\]
\end{lemma}

\begin{proof}
The self-referential rule applies Definition~\ref{def:semantic-rule} with $\Ccal^{\omega_{\pi^{\mathrm{sr}}}^t}$ in place of $\Ccal$, and Proposition~\ref{prop:semantic-is-promotion-supporting} transfers to any partition by the same argument: the class-targeting distribution $\nu_t^{\mathrm G}$ assigns mass $\ge\delta_0\bar w^{\omega_{\pi^{\mathrm{sr}}}^t}(c)$ to each class $c$, and the conditional $\Prob(A_{t+1}=a_t^c(h_t)\mid C_t=c,H_t=h_t)=1$ by the selector construction. The universal prior of Definition~\ref{def:universal-prior} assigns positive mass to every computable program, so every class in any partition has $\bar w^\omega(c)>0$.
\end{proof}

\begin{proposition}[Eventual isolation criterion]
\label{prop:observer-promotion-sr}
Fix a target program $p^\star\in\Pcal$, and let
\[
c_t^\star:=[p^\star]_{\omega_{\pi^{\mathrm{sr}}}^t}.
\]
Suppose there exists a finite time $T^\star$ such that
\[
c_t^\star=[p^\star]
\qquad\text{for every }t\ge T^\star.
\]
Then
\[
R_{\pi^{\mathrm{sr}}}(p^\star)\;=\;[p^\star].
\]
If, more generally, for every $p\in\Pcal$ there exists a finite time $T_p$ such that
\[
[p]_{\omega_{\pi^{\mathrm{sr}}}^t}=[p]
\qquad\text{for every }t\ge T_p,
\]
then
\[
\omega_{\pi^{\mathrm{sr}}}\;=\;\omega_{\mathrm{behav}}.
\]
\end{proposition}

\begin{proof}
Lemma~\ref{lem:nesting} gives $[p^\star]\subseteq R_{\pi^{\mathrm{sr}}}(p^\star)$ for every policy, so only the reverse inclusion needs proof. Let $p\notin[p^\star]$. By hypothesis, $p\notin c_t^\star$ for every $t\ge T^\star$, hence $p\not\sim_{\omega_{\pi^{\mathrm{sr}}}^t}p^\star$ for every $t\ge T^\star$. Because $\omega_{\pi^{\mathrm{sr}}}^t\le\omega_{\pi^{\mathrm{sr}}}$ by Lemma~\ref{lem:monotone-refinement}, indistinguishability under $\omega_{\pi^{\mathrm{sr}}}$ would imply indistinguishability under every $\omega_{\pi^{\mathrm{sr}}}^t$, a contradiction. Thus $p\notin R_{\pi^{\mathrm{sr}}}(p^\star)$, proving $R_{\pi^{\mathrm{sr}}}(p^\star)=[p^\star]$.

For the second statement, fix $p_1,p_2$ with $p_1\not\sim_{\mathrm{behav}}p_2$. Since $p_2\notin[p_1]$, the hypothesis at $p_1$ gives $p_2\not\sim_{\omega_{\pi^{\mathrm{sr}}}^t}p_1$ for every $t\ge T_{p_1}$. Again using $\omega_{\pi^{\mathrm{sr}}}^t\le\omega_{\pi^{\mathrm{sr}}}$, one gets $p_1\not\sim_{\omega_{\pi^{\mathrm{sr}}}}p_2$. Thus behavioral inequivalence implies policy-observer inequivalence. The converse implication is automatic from $\omega_{\pi^{\mathrm{sr}}}\le\omega_{\mathrm{behav}}$, so the two observers induce the same equivalence relation.
\end{proof}

Proposition~\ref{prop:observer-promotion-sr} isolates the missing step behind the ``stone-turning'' framing: the self-referential rule would recover the behavioral target once its running partition actually isolates the true semantic class. What remains unproved under \textup{(C1)--(C3)} alone is that this eventual isolation occurs.

\subsection{Self-referential convergence}
\label{sec:self-ref-convergence}

\begin{theorem}[Self-referential convergence under eventual isolation]
\label{thm:self-ref-convergence}
Assume \textup{(C1)--(C3)} and $\omega_B=\omega_{\mathrm{syn}}$ with Bayes/Gibbs belief \textup{(C2)}. Fix a target program $p^\star\in\Pcal$. Under the self-referential semantic rule $\pi^{\mathrm{sr}}$ of Definition~\ref{def:self-ref-rule}, let
\[
c^\star:=[p^\star],
\qquad
c_t^\star:=[p^\star]_{\omega_{\pi^{\mathrm{sr}}}^t}.
\]
If there exists a finite time $T^\star$ such that $c_t^\star=c^\star$ for every $t\ge T^\star$, then
\[
q_t^\star([p^\star]\mid H_t)\;\longrightarrow\;1
\qquad\text{almost surely.}
\]
\end{theorem}

\begin{proof}
Let $\mathcal{F}_t:=\sigma(H_t,C_0,\dots,C_{t-1})$ as in Theorem~\ref{thm:semantic-convergence}, and define
\[
E_t:=\{t\ge T^\star,\ q_t^\star(c^\star\mid H_t)<1\},
\qquad
U_{t+1}:=E_t\cap\{C_t=c^\star\}.
\]
On $\{t\ge T^\star\}$ the running class equals the semantic class, so the self-referential class-targeting distribution satisfies
\[
\Prob(U_{t+1}\mid\mathcal{F}_t)
=
\one_{E_t}\nu_t^{\mathrm G}(c^\star\mid H_t)
\ge
\delta_0\bar w(c^\star)\one_{E_t}.
\]
Set $\hat\eta:=\tfrac{1}{2}\min\{1,\eta(c^\star)\}>0$. On $\{t\ge T^\star\}$ the class-complement law of the running partition for $c_t^\star=c^\star$ is exactly the semantic class-complement law, because regrouping classes outside $c^\star$ does not change their pooled posterior predictive mixture. Hence whenever $U_{t+1}$ occurs, the selector used by the self-referential rule obeys
\[
B_t(c^\star,A_{t+1};H_t)
=
B_t^{\omega_{\pi^{\mathrm{sr}}}^t}(c^\star,A_{t+1};H_t)
\ge
\hat\eta.
\]
If $E_t$ occurs infinitely often, then $\sum_t \Prob(U_{t+1}\mid\mathcal{F}_t)=\infty$ almost surely, so Lemma~\ref{lem:conditional-bc} gives $U_{t+1}$ infinitely often almost surely. Therefore
\[
\sum_{t=0}^{\infty}B_t(c^\star,A_{t+1};H_t)=\infty
\qquad\text{almost surely on }E_t\text{ i.o.}
\]
and Lemma~\ref{lem:contraction} yields $q_t^\star(c^\star\mid H_t)\to 1$ on that event. If $E_t^c$ eventually holds, then from time $T^\star$ onward either $t<T^\star$ or $q_t^\star(c^\star\mid H_t)=1$, so the same conclusion is immediate.
\end{proof}

\begin{proposition}[Unconditional eventual isolation fails]
\label{prop:self-ref-obstruction}
There exist a finite action alphabet $\Acal$, a finite observation alphabet $\Ocal$, three computable environments $p_0,p_1,p_2\in\Pcal$, a substituted prior supported on $\{p_0,p_1,p_2\}$, and an admissible completion of the arbitrary selector choices in Definitions~\ref{def:semantic-rule} and~\ref{def:self-ref-rule} such that:
\begin{itemize}
\item[(i)] the three semantic classes are singletons and satisfy the semantic separation condition \textup{(C3)};
\item[(ii)] with $p^\star:=p_0$, the running true class under the self-referential rule never isolates $[p^\star]$: specifically,
\[
\{p_0,p_1\}\subseteq [p^\star]_{\omega_{\pi^{\mathrm{sr}}}^t}
\qquad\text{for every }t\ge 1;
\]
\item[(iii)] consequently
\[
R_{\pi^{\mathrm{sr}}}(p^\star)\supsetneq [p^\star].
\]
\end{itemize}
In particular, the eventual-isolation hypothesis of Proposition~\ref{prop:observer-promotion-sr} is not automatic for the self-referential rule.
\end{proposition}

\begin{proof}
Let
\[
\Acal:=\{a^{\mathrm{ref}},a^0,a^1,a^2\},
\qquad
\Ocal:=\{0,1,2\},
\]
with the action enumeration ordered as written, and fix $\varepsilon\in(0,e^{-2})$. Define three history-independent computable environments by
\[
\mu_{p_0}(\cdot\mid h,a^{\mathrm{ref}})=\delta_0,
\qquad
\mu_{p_1}(\cdot\mid h,a^{\mathrm{ref}})=\delta_0,
\qquad
\mu_{p_2}(\cdot\mid h,a^{\mathrm{ref}})=\varepsilon\delta_0+(1-\varepsilon)\delta_2,
\]
and
\[
\mu_{p_0}(\cdot\mid h,a^0)=\delta_0,
\qquad
\mu_{p_1}(\cdot\mid h,a^0)=\delta_1,
\qquad
\mu_{p_2}(\cdot\mid h,a^0)=\delta_1,
\]
\[
\mu_{p_0}(\cdot\mid h,a^1)=\delta_0,
\qquad
\mu_{p_1}(\cdot\mid h,a^1)=\delta_1,
\qquad
\mu_{p_2}(\cdot\mid h,a^1)=\delta_0,
\]
\[
\mu_{p_0}(\cdot\mid h,a^2)=\delta_0,
\qquad
\mu_{p_1}(\cdot\mid h,a^2)=\delta_0,
\qquad
\mu_{p_2}(\cdot\mid h,a^2)=\delta_2.
\]
Let the substituted prior $w^\flat$ assign positive mass to $p_0,p_1,p_2$ and zero mass elsewhere.
Each semantic class is a singleton because the three kernels disagree on at least one action. Moreover, (C3) holds: $a^0$ separates $p_0$ from $\{p_1,p_2\}$, $a^1$ separates $p_1$ from $\{p_0,p_2\}$, and $a^2$ separates $p_2$ from $\{p_0,p_1\}$, each with Bhattacharyya separation $+\infty$.

Choose the admissible completion of the arbitrary selector choices that picks the first action in the enumeration whenever Definition~\ref{def:semantic-rule} leaves the selector unspecified. Then $\pi_1^{\mathrm{sr}}$ plays $a^{\mathrm{ref}}$. Under this first-step policy, $p_0$ and $p_1$ induce the same law of $H_1$, while $p_2$ induces a different law, so
\[
\{p_0,p_1\}\subseteq [p_0]_{\omega_{\pi^{\mathrm{sr}}}^1}.
\]
We prove by induction on $t\ge 1$ that
\[
\pi_i^{\mathrm{sr}}\equiv a^{\mathrm{ref}}
\quad\text{for }1\le i\le t,
\qquad
\{p_0,p_1\}\subseteq [p_0]_{\omega_{\pi^{\mathrm{sr}}}^t}.
\]
The base case $t=1$ is the previous paragraph. For the induction step, assume the claim through time $t$ and write
\[
c_t^\star:=[p_0]_{\omega_{\pi^{\mathrm{sr}}}^t}.
\]
Under the induction hypothesis, the supported programs in $c_t^\star$ are exactly $p_0$ and $p_1$, while the supported complement is $p_2$. Hence at any history $h_t$ generated under the policy prefix that always plays $a^{\mathrm{ref}}$, the observer-$\omega_{\pi^{\mathrm{sr}}}^t$ class-predictive laws satisfy
\[
\mu_{c_t^\star}^{q_t^\star}(\cdot\mid h_t,a^{\mathrm{ref}})=\delta_0,
\qquad
\Qcal_t^{-c_t^\star,\omega_{\pi^{\mathrm{sr}}}^t}(\cdot\mid h_t,a^{\mathrm{ref}})=\varepsilon\delta_0+(1-\varepsilon)\delta_2,
\]
hence
\[
B_t^{\omega_{\pi^{\mathrm{sr}}}^t}(c_t^\star,a^{\mathrm{ref}};h_t)
=
-\log\sqrt{\varepsilon}>1.
\]
Therefore the capped score of the running true class at $a^{\mathrm{ref}}$ already equals $1$, so the first-satisfying selector again chooses $a^{\mathrm{ref}}$ for $c_t^\star$. On histories where the selector is unspecified because the posterior on a running class is $0$ or $1$, the chosen admissible completion also selects $a^{\mathrm{ref}}$. Induction on $t$ gives that $\pi_t^{\mathrm{sr}}$ plays $a^{\mathrm{ref}}$ at every history for every $t\ge 1$.

Under the policy that always plays $a^{\mathrm{ref}}$, the full history laws of $p_0$ and $p_1$ coincide, while the law of $p_2$ differs. Hence
\[
\{p_0,p_1\}\subseteq [p_0]_{\omega_{\pi^{\mathrm{sr}}}^t}
\qquad\text{for every }t\ge 1,
\]
proving (ii). Since $\mu_{p_0}\neq\mu_{p_1}$, one has $p_1\notin[p_0]$ but $p_1\in R_{\pi^{\mathrm{sr}}}(p_0)$, proving (iii).
\end{proof}

\subsection{Exploration-completed self-reference}
\label{sec:self-ref-exploration}

The obstruction of Proposition~\ref{prop:self-ref-obstruction} is structural for the \emph{pure} self-referential rule: a current true class can remain externally separated from its complement while never being split internally. One way to close the gap is to modify the rule rather than to strengthen the hypothesis. Section~\ref{sec:main} now supplies the policy-level route needed for that modification: Proposition~\ref{prop:full-support-behavioral} identifies the recovered observer, Theorem~\ref{thm:separating-support-convergence} gives the master semantic-recovery condition, and Theorem~\ref{thm:exploration-floor-behavioral} turns a divergent full-support exploration budget into behavioral-observer recovery plus posterior concentration. The exploration-completed self-referential rule below is the corresponding specialization.

\begin{definition}[Exploration-completed self-referential rule]
\label{def:self-ref-exploratory}
Fix a full-support reference distribution $\lambda\in\Delta(\Acal)$ and an exploration schedule $(\varepsilon_t)_{t\ge 1}$ with
\[
\varepsilon_t\in(0,1]
\qquad\text{for every }t\ge 1.
\]
The \emph{exploration-completed self-referential rule} is the policy $\pi^{\mathrm{xsr}}$ defined by
\[
\pi_{t+1}^{\mathrm{xsr}}(a\mid h_t)
:=
(1-\varepsilon_{t+1})\,\pi_{t+1}^{\mathrm{sr}}(a\mid h_t)+\varepsilon_{t+1}\,\lambda(a),
\qquad
a\in\Acal.
\]
\end{definition}

\begin{theorem}[Exploration-completed self-reference with divergent exploration budget]
\label{thm:self-ref-exploratory}
Assume the standing finite-action setup, \textup{(C1)--(C3)}, and $\omega_B=\omega_{\mathrm{syn}}$ with Bayes/Gibbs belief \textup{(C2)}. Fix a full-support reference distribution $\lambda\in\Delta(\Acal)$ and an exploration schedule $(\varepsilon_t)_{t\ge 1}$ satisfying
\[
\sum_{t=1}^{\infty}\varepsilon_t=\infty.
\]
Let action follow the exploration-completed self-referential rule $\pi^{\mathrm{xsr}}$ of Definition~\ref{def:self-ref-exploratory}. Then for every $p^\star\in\Pcal$,
\[
\omega_{\pi^{\mathrm{xsr}}}=\omega_{\mathrm{behav}},
\qquad
q_t^\star([p^\star]\mid H_t)\longrightarrow 1
\quad\text{almost surely.}
\]
\end{theorem}

\begin{proof}
The policy $\pi^{\mathrm{xsr}}$ satisfies
\[
\pi_{t+1}^{\mathrm{xsr}}(a\mid h_t)\ge \varepsilon_{t+1}\lambda(a)
\qquad
\text{for every }t\ge 0,\ h_t,\ a\in\Acal
\]
by Definition~\ref{def:self-ref-exploratory}. The conclusion is therefore the special case of Theorem~\ref{thm:exploration-floor-behavioral} with $\pi=\pi^{\mathrm{xsr}}$.
\end{proof}

\begin{theorem}[Rates for exploration-completed self-reference]
\label{thm:self-ref-exploratory-rate}
Under the assumptions of Theorem~\ref{thm:self-ref-exploratory}, fix a target program $p^\star\in\Pcal$ and write $c^\star:=[p^\star]$. Suppose in addition that there exists $m\in(0,1)$ such that whenever $0<q_t^\star(c^\star\mid H_t)<1$,
\[
m\le
\frac{\mu_{c^\star}(o\mid H_t,A_{t+1})}{\Qcal_t^{-c^\star}(o\mid H_t,A_{t+1})}
\le \frac{1}{m}
\]
for every $o$ with $\mu_{c^\star}(o\mid H_t,A_{t+1})>0$. Write
\[
\Gamma_t:=\sum_{i=1}^t \varepsilon_i,
\qquad
\mu_{\mathrm{xsr}}:=2\,\eta(c^\star)\,\lambda_{\min},
\qquad
\lambda_{\min}:=\min_{a\in\Acal}\lambda(a)>0,
\qquad
L:=\log(1/m).
\]
Then:
\begin{itemize}
\item[(i)] for every $t\ge 0$,
\[
\E[Z_t]\ge Z_0+\mu_{\mathrm{xsr}}\,\Gamma_t;
\]
\item[(ii)] for every $t\ge 1$ and every $x>0$,
\[
\Prob\!\bigl(Z_t\le Z_0+\mu_{\mathrm{xsr}}\Gamma_t-x\bigr)
\le
\exp\!\left(-\frac{x^2}{2tL^2}\right);
\]
\item[(iii)] equivalently, for every $\delta\in(0,1)$ and $t\ge 1$,
\[
\Prob\!\left(
q_t^\star(c^\star\mid H_t)\ge
\frac{1}{1+\exp\!\bigl(-Z_0-\mu_{\mathrm{xsr}}\Gamma_t+L\sqrt{2t\log(1/\delta)}\bigr)}
\right)\ge 1-\delta.
\]
\end{itemize}
In particular, because $\Gamma_t\to\infty$, the expected true-class log-odds diverge and the high-probability posterior lower bound improves along the cumulative exploration budget.
\end{theorem}

\begin{proof}
Let
\[
S_t^\star:=\{a\in\Acal:B_t(c^\star,a;H_t)\ge \eta(c^\star)\}.
\]
By Proposition~\ref{prop:finite-action-kernel-separation}, one has
\[
\lambda(S_t^\star)\ge \lambda_{\min}
\qquad\text{on }\{q_t^\star(c^\star\mid H_t)<1\},
\]
so Definition~\ref{def:self-ref-exploratory} gives
\[
\pi_{t+1}^{\mathrm{xsr}}(S_t^\star\mid H_t)\ge \varepsilon_{t+1}\lambda_{\min}
\qquad\text{almost surely on }\{q_t^\star(c^\star\mid H_t)<1\}.
\]
Applying Theorem~\ref{thm:separating-support-rate} with
\[
r_t:=\varepsilon_t\lambda_{\min},
\qquad
R_t=\lambda_{\min}\Gamma_t,
\qquad
\mu_{\mathrm{xsr}}=\mu_{\mathrm{sep}}\lambda_{\min}
=2\,\eta(c^\star)\lambda_{\min},
\]
gives the stated bounds.
\end{proof}

\begin{proposition}[One-step splitting criterion]
\label{prop:self-ref-one-step-split}
Fix a target program $p^\star\in\Pcal$, and let
\[
c_t^\star:=[p^\star]_{\omega_{\pi^{\mathrm{sr}}}^t}.
\]
Suppose $p\in c_t^\star\setminus[p^\star]$ and there exists a history $h_t$ with
\[
\Prob_{p^\star,\pi^{\mathrm{sr}}}(H_t=h_t)>0
\]
such that
\[
\mu_p(\cdot\mid h_t,a_t^{c_t^\star}(h_t))
\neq
\mu_{p^\star}(\cdot\mid h_t,a_t^{c_t^\star}(h_t)).
\]
Then
\[
p\notin [p^\star]_{\omega_{\pi^{\mathrm{sr}}}^{t+1}}.
\]
In particular, if the displayed condition holds for at least one $p\in c_t^\star\setminus[p^\star]$, then
\[
c_{t+1}^\star\subsetneq c_t^\star.
\]
\end{proposition}

\begin{proof}
Because $p\in c_t^\star$, the length-$t$ history laws under the policy prefix $(\pi_1^{\mathrm{sr}},\dots,\pi_t^{\mathrm{sr}})$ coincide for $p$ and $p^\star$. Hence
\[
\Prob_{p,\pi^{\mathrm{sr}}}(H_t=h_t)=\Prob_{p^\star,\pi^{\mathrm{sr}}}(H_t=h_t)>0.
\]
Lemma~\ref{lem:exploration-reachability} gives
\[
\pi_{t+1}^{\mathrm{sr}}(a_t^{c_t^\star}(h_t)\mid h_t)>0.
\]
Choose $o\in\Ocal$ with
\[
\mu_p(o\mid h_t,a_t^{c_t^\star}(h_t))
\neq
\mu_{p^\star}(o\mid h_t,a_t^{c_t^\star}(h_t));
\]
at least one such $o$ exists by hypothesis. For the length-$(t+1)$ history
\[
h_{t+1}:=(h_t,a_t^{c_t^\star}(h_t),o),
\]
the factorization identity gives
\begin{align*}
\Prob_{p,\pi^{\mathrm{sr}}}(H_{t+1}=h_{t+1})
&=
\Prob_{p,\pi^{\mathrm{sr}}}(H_t=h_t)\,
\pi_{t+1}^{\mathrm{sr}}(a_t^{c_t^\star}(h_t)\mid h_t)\,
\mu_p(o\mid h_t,a_t^{c_t^\star}(h_t)),\\
\Prob_{p^\star,\pi^{\mathrm{sr}}}(H_{t+1}=h_{t+1})
&=
\Prob_{p^\star,\pi^{\mathrm{sr}}}(H_t=h_t)\,
\pi_{t+1}^{\mathrm{sr}}(a_t^{c_t^\star}(h_t)\mid h_t)\,
\mu_{p^\star}(o\mid h_t,a_t^{c_t^\star}(h_t)).
\end{align*}
The first two factors agree and are positive, while the third differs. Thus the two length-$(t+1)$ history laws are different, so $p\not\sim_{\omega_{\pi^{\mathrm{sr}}}^{t+1}}p^\star$. The final claim is immediate.
\end{proof}

\begin{theorem}[Sharpened self-referential convergence under deterministic splitting]
\label{thm:self-ref-sharp}
Assume \textup{(C1)--(C3)} and $\omega_B=\omega_{\mathrm{syn}}$ with Bayes/Gibbs belief \textup{(C2)}. Fix a target program $p^\star\in\Pcal$. Under the self-referential semantic rule $\pi^{\mathrm{sr}}$ of Definition~\ref{def:self-ref-rule}, let
\[
c_t^\star:=[p^\star]_{\omega_{\pi^{\mathrm{sr}}}^t}.
\]
Suppose there exist integers $T_0\ge 0$ and $M\ge 1$ such that:
\begin{itemize}
\item[(SR1)] $c_{T_0}^\star$ is the union of at most $M$ behavioral classes;
\item[(SR2)] for every $t\ge T_0$ with $c_t^\star\neq [p^\star]$, there exist a program
\[
p_t\in c_t^\star\setminus[p^\star]
\]
and a history $h_t$ with $\Prob_{p^\star,\pi^{\mathrm{sr}}}(H_t=h_t)>0$ such that
\[
\mu_{p_t}(\cdot\mid h_t,a_t^{c_t^\star}(h_t))
\neq
\mu_{p^\star}(\cdot\mid h_t,a_t^{c_t^\star}(h_t)).
\]
\end{itemize}
Then
\[
c_t^\star=[p^\star]
\qquad\text{for every }t\ge T_0+M-1,
\]
and consequently
\[
q_t^\star([p^\star]\mid H_t)\longrightarrow 1
\qquad\text{almost surely.}
\]
\end{theorem}

\begin{proof}
Because $\omega_{\pi^{\mathrm{sr}}}^t\le\omega_{\mathrm{behav}}$ by Lemma~\ref{lem:monotone-refinement}, every running class $c_t^\star$ is a union of behavioral classes. Fix $t\ge T_0$ with $c_t^\star\neq[p^\star]$. By (SR2) and Proposition~\ref{prop:self-ref-one-step-split},
\[
c_{t+1}^\star\subsetneq c_t^\star.
\]
Since both sets are unions of behavioral classes and both contain $[p^\star]$, each strict inclusion removes at least one behavioral class distinct from $[p^\star]$.

At time $T_0$, the running true class contains at most $M$ behavioral classes by (SR1). Therefore after at most $M-1$ strict refinements, no behavioral class other than $[p^\star]$ can remain inside the running true class. Hence
\[
c_t^\star=[p^\star]
\qquad\text{for every }t\ge T_0+M-1.
\]
Applying Theorem~\ref{thm:self-ref-convergence} with $T^\star:=T_0+M-1$ gives the almost-sure posterior convergence.
\end{proof}

\begin{remark}[Finite-time proxy and remaining gap]
\label{rem:self-ref-computability}
At every finite step $t$, the self-referential rule depends only on the policy prefix $(\pi^{\mathrm{sr}}_1,\dots,\pi^{\mathrm{sr}}_t)$ and the corresponding length-$t$ marginal laws $\Prob_{p,\pi^{\mathrm{sr}}}(H_t\in\cdot)$. Under the standing finite-action/countable-observation setup, these marginals are generally countably rather than finitely supported, so the relevant objects are semicomputable cylinder probabilities rather than finite tables. The present section therefore does not claim exact computability of the running partition. Its contribution is the reduction: once the running finite-time partition isolates $[p^\star]$, Proposition~\ref{prop:observer-promotion-sr} and Theorem~\ref{thm:self-ref-convergence} transfer the semantic-convergence conclusion.
\par
Theorems~\ref{thm:separating-support-convergence} and~\ref{thm:exploration-floor-behavioral} take a different route: they replace the running-partition argument by policy-level support conditions. Theorem~\ref{thm:self-ref-exploratory} is the corresponding self-referential specialization, obtained by adding a full-support exploration schedule with divergent cumulative mass and thereby closing the observer gap without proving eventual isolation for the pure running partition itself.
\end{remark}

\begin{remark}[Relation to Proposition~\ref{prop:goal-dialed}]
\label{rem:self-ref-vs-goal-dialed}
Proposition~\ref{prop:goal-dialed} is a \emph{fixed-observer} convergence statement: for any single $\omega_A\le\omega_{\mathrm{behav}}$, the rule at $\omega_A$ drives the posterior onto $[p^\star]_{\omega_A}$. Theorem~\ref{thm:self-ref-convergence} is a \emph{conditional time-inhomogeneous transfer}: the running observer $\omega_{\pi^{\mathrm{sr}}}^t$ may refine with $t$, and convergence to $[p^\star]$ follows once that running observer eventually isolates the true behavioral class. The two results do not collide: Proposition~\ref{prop:goal-dialed} treats the observer dial as a coarsening knob for deliberate targeting at fixed $\omega_A\le\omega_{\mathrm{behav}}$; Theorem~\ref{thm:self-ref-convergence} identifies the additional isolation step required to turn the behavioral target into a self-referential proxy.
\par
Theorems~\ref{thm:separating-support-convergence} and~\ref{thm:exploration-floor-behavioral} lie alongside, not between, these two readings: they are neither fixed-observer dial results nor pure running-partition theorems, but policy-level results showing that enough cumulative mass on separating actions forces concentration, and that a full-support exploration floor upgrades this to recovery of the behavioral observer. Theorem~\ref{thm:self-ref-exploratory} is the self-referential specialization of that policy-level route.
\end{remark}

The self-referential rule turns the computability issue into a sharply defined reduction. Proposition~\ref{prop:self-ref-obstruction} shows that unconditional eventual isolation is false in general: a current true class can be stably separated from its complement while never being split internally. There are then two ways forward. The policy-level route is given by Theorems~\ref{thm:separating-support-convergence},~\ref{thm:exploration-floor-behavioral}, and~\ref{thm:separating-support-rate}: enough cumulative mass on separating actions yields concentration, a full-support exploration floor recovers the behavioral observer, and the same mechanism carries explicit expectation and high-probability bounds. Theorem~\ref{thm:self-ref-exploratory} and Theorem~\ref{thm:self-ref-exploratory-rate} are the self-referential specializations of that route. The pure self-referential route instead relies on Proposition~\ref{prop:self-ref-one-step-split} and Theorem~\ref{thm:self-ref-sharp}, which close the gap under deterministic selector-driven splitting.

\section{The Algorithmic Free Energy Principle as Boundary Example}
\label{sec:boundary-example}
\label{sec:near-miss}

Section~\ref{sec:main} identified the canonical Bayes/Gibbs-plus-class-against-complement pair at $\omega_{\mathrm{behav}}$ and proved almost-sure concentration on $[p^\star]$. The \emph{algorithmic free energy principle} --- Bayes/Gibbs belief under the universal prior paired with expected-free-energy action --- is the Bayes/Gibbs-plus-MI boundary specialization of the same two-observer framework, in its raw policy-coupled MI form, at action observer $\omega_A=\omega_{\mathrm{syn}}$, a setting strictly above the meeting point of Theorem~\ref{thm:meeting-point}. The Bayes/Gibbs architecture that delivered the belief-side unification of Section~\ref{sec:belief} therefore corresponds to a \emph{wrong-observer} instance of the framework: the action bonus is the syntactic-class mutual information rather than the behavioral-class Bhattacharyya separation, and by Proposition~\ref{prop:action-cap} this bonus is capped on behavioral twins. Theorem~\ref{thm:afe-near-miss} below exhibits the structural failure of this instance, and Theorem~\ref{thm:observer-promotion-failure} reads it at the observer level. The AFE principle is not a competitor of the main theorem but its boundary example, fixing where the admissible range ends.

\subsection{Boundary identification}
\label{sec:boundary-identification}

\begin{proposition}[Boundary identification under Bayes/Gibbs belief]
\label{prop:boundary-identity}
Fix a target program $p^\star\in\Pcal$. 
Under Bayes/Gibbs belief, the algorithmic free energy principle of Definition~\ref{def:afe-principle} matches the $\omega_A=\omega_{\mathrm{syn}}$ action-MI specialization of the two-observer framework at $\beta=1$: its belief is the unique minimizer of the belief term $\Fcal_t$ (Lemma~\ref{lem:variational}) and, with zero action cost and uniform preferences on a common finite feasible observation set, its action rule selects $a_{t+1}\in\arg\max_a \I_{q_t^\star}(P;O\mid h_t,a)$, which is exactly the action-side maximizer of the MI bonus $\I_q(C^{\omega_A};O\mid h_t,a)$ at $\omega_A=\omega_{\mathrm{syn}}$ when $q=q_t^\star$ and $C^{\omega_A}=P$. Since $\omega_{\mathrm{syn}}>\omega_{\mathrm{behav}}$, this action observer lies strictly above the meeting point of Theorem~\ref{thm:meeting-point}, and Proposition~\ref{prop:action-cap} caps the action bonus on every pair of behavioral twins. Conversely, the Bayes/Gibbs-plus-class-against-complement pair of Section~\ref{sec:main} lives at the meeting-point observer $\omega_{\mathrm{behav}}$ and, under \textup{(C1)--(C4)} --- equivalently, under the full Bayes/Gibbs-plus-class-against-complement package --- achieves $q_t^\star([p^\star]\mid H_t)\to 1$ almost surely (Theorem~\ref{thm:semantic-convergence}).
\end{proposition}

\begin{proof}
The belief identity is Lemma~\ref{lem:variational}: the Bayes/Gibbs posterior is the unique $q$-minimizer of $\Fcal_t$. For the action rule, Lemma~\ref{lem:risk-ig} gives
\[
\Gcal_t(a;h_t)
=
c_t(a;h_t)+\E_{o\sim\Qcal_t^O}[-\log\rho_t(o\mid h_t,a)]-\I_{\Qcal_t^{PO}}(P;O\mid h_t,a),
\]
so under zero action cost and action-independent preferences, $\arg\min_a\Gcal_t(a;h_t)=\arg\max_a\I(P;O\mid h_t,a)$. In the action-MI specialization at $\omega_A=\omega_{\mathrm{syn}}$ the class random variable $C^{\omega_A}$ coincides with $P$, so this is exactly the action-side maximizer of the MI bonus when the belief input is fixed at $q=q_t^\star$ and $\beta=1$. This is a conditional identification of the action side under Bayes/Gibbs belief, not a claim of joint minimization over coupled choices of $(q,\pi)$ for $\mathcal J_{t,\mathrm{raw}}^{\omega_{\mathrm{syn}},\omega_{\mathrm{syn}},\mathrm{MI}}$, since the MI term itself depends on $q$. The cap claim is Proposition~\ref{prop:action-cap} applied at $\omega_A=\omega_{\mathrm{syn}}>\omega_{\mathrm{behav}}$. The converse clause is Theorem~\ref{thm:semantic-convergence}.
\end{proof}

\begin{remark}[Reading under the functional]
\label{rem:boundary-identity}
Proposition~\ref{prop:boundary-identity} locates the AFE principle structurally. The belief-side unification of Section~\ref{sec:belief} is unchanged: algorithmic free energy is still the belief term of $\mathcal J_t^{\omega_{\mathrm{syn}},\omega_A}$ at the canonical belief observer, and the Bayes/Gibbs posterior is still its exact minimizer (Lemma~\ref{lem:variational}). The change is on the action side. Fixing $\omega_A=\omega_{\mathrm{syn}}$ and using mutual information rather than Bhattacharyya separation selects, under Bayes/Gibbs belief, an action bonus that rewards refining distinctions between \emph{programs}, which by Proposition~\ref{prop:action-cap} is capped on behavioral twins and, by Theorem~\ref{thm:afe-near-miss} below, can be deployed entirely inside $R_\pi(p^\star)\setminus[p^\star]$. The two-observer functional's admissible range $\omega_A\le\omega_{\mathrm{behav}}$ excludes this setting exactly, and the main theorem's convergence payoff (Theorem~\ref{thm:semantic-convergence}) requires an action rule at an observer inside the admissible range.
\end{remark}

\subsection{Expected free energy}

\begin{definition}[Expected free energy]
\label{def:efe}
At history $h_t$, define the joint predictive law over latent program $P$ and next observation $O$ by
\[
\Qcal_t^{PO}(p,o\mid h_t,a):=q_t^\star(p\mid h_t)\mu_p(o\mid h_t,a).
\]
With observation marginal $\Qcal_t^O(\cdot\mid h_t,a)$, preference distribution $\rho_t(\cdot\mid h_t,a)$, and action cost $c_t(a;h_t)$, the \emph{expected free energy} is
\[
\Gcal_t(a;h_t)
:=
c_t(a;h_t)
+
\E_{\Qcal_t^{PO}}\!\bigl[\log q_t^\star(P\mid h_t)-\log q_{t+1}^\star(P\mid h_t,a,O)-\log\rho_t(O\mid h_t,a)\bigr].
\]
\end{definition}

\begin{lemma}[Risk minus information gain]
\label{lem:risk-ig}
For every history $h_t$ and action $a$,
\[
\Gcal_t(a;h_t)
=
c_t(a;h_t)
+\E_{o\sim\Qcal_t^O}\bigl[-\log\rho_t(o\mid h_t,a)\bigr]
-\I_{\Qcal_t^{PO}}(P;O\mid h_t,a).
\]
\end{lemma}

\begin{proof}
The first expectation depends only on the observation marginal and equals the risk term. The second expectation equals $-\I_{\Qcal_t^{PO}}(P;O\mid h_t,a)$ by the standard mutual-information identity. Summing gives the claim.
\end{proof}

\begin{corollary}[Expected free energy as a unifying epistemic principle]
\label{cor:efe-specialization}
Fix zero action cost and uniform preferences $\rho_t$ on a common finite feasible observation set. Then expected free energy minimization is equivalent, up to an action-independent constant, to maximization of the mutual information $\I_{\Qcal_t^{PO}}(P;O\mid h_t,a)$ between the latent program and the next observation.
\end{corollary}

\begin{proof}
By Lemma~\ref{lem:risk-ig}, $\Gcal_t(a;h_t)=c_t(a;h_t)+\E_{o\sim\Qcal_t^O}[-\log\rho_t(o\mid h_t,a)]-\I_{\Qcal_t^{PO}}(P;O\mid h_t,a)$. Zero action cost annihilates the first term. Uniform $\rho_t$ on a common finite feasible set makes the second term action-independent. The displayed specialization follows.
\end{proof}

\begin{remark}
\label{rem:efe-unification}
Corollary~\ref{cor:efe-specialization} is the action-side unification. In this regime, expected-free-energy minimization is simultaneously the Bayesian experimental design criterion \citep{Lindley1956,ChalonerVerdinelli1995,MacKay1992}, the compression-based curiosity criterion \citep{Schmidhuber1991,Schmidhuber2010}, the active-inference curiosity criterion \citep{FristonEtAl2017Curiosity}, and the information-theoretic bounded-rationality criterion \citep{OrtegaBraun2011,OrtegaBraun2013}. These traditions are not approximations of each other; they coincide on the action rule they prescribe. This is what makes the near-miss proved below substantive: Theorem~\ref{thm:afe-near-miss} is a failure of generic information-seeking itself, across a family of traditions that agree on what generic information-seeking means.
\end{remark}

\subsection{The algorithmic free energy principle}

\begin{definition}[Algorithmic free energy principle]
\label{def:afe-principle}
A learner obeys the \emph{algorithmic free energy principle} if
\begin{itemize}
\item[(i)] its prior $w$ is the universal interactive prior of Definition~\ref{def:universal-prior};
\item[(ii)] its posterior $q_t^\star(\cdot\mid h_t)$ is the Bayes/Gibbs posterior of Lemma~\ref{lem:variational};
\item[(iii)] its action $A_{t+1}$ is chosen to minimize the expected free energy $\Gcal_t(\cdot\mid h_t)$ with zero action cost and uniform preferences on a common finite feasible observation set, breaking ties by a fixed measurable rule.
\end{itemize}
\end{definition}

\begin{remark}[The AFE principle as a Bayes/Gibbs-plus-MI boundary specialization]
\label{rem:afe-as-functional-minimizer}
Under the two-observer framework, Proposition~\ref{prop:boundary-identity} identifies Definition~\ref{def:afe-principle}'s triple with exact Bayes/Gibbs belief together with action-side MI maximization at $\omega_A=\omega_{\mathrm{syn}}$. Item (i) fixes the reference measure shared by $\mathcal J_t$ and its raw MI scaffold; item (ii) is the $q$-minimization of the belief term $\Fcal_t$; and item (iii) is the action-side maximization of the MI bonus in $\mathcal J_{t,\mathrm{raw}}^{\omega_{\mathrm{syn}},\omega_{\mathrm{syn}},\mathrm{MI}}$, evaluated at $q=q_t^\star$. The principle therefore has a short synopsis in this paper's vocabulary: \emph{Bayes/Gibbs plus top-observer MI action}. Theorem~\ref{thm:afe-near-miss} below is the proof that this boundary specialization fails to concentrate on $[p^\star]$.
\end{remark}

\subsection{Information decomposition relative to $[p^\star]$}

\begin{lemma}[Information decomposition]
\label{lem:info-decomp}
Fix a target program $p^\star\in\Pcal$, a history $h_t$, and an action $a$. Let $c^\star:=[p^\star]$ be the true interventional semantic class, let $\alpha_t:=q_t^\star(c^\star\mid h_t)\in(0,1)$, and let $Z:=\one\{P\in c^\star\}$ under $P\sim q_t^\star(\cdot\mid h_t)$. Then
\[
\I(P;O\mid h_t,a)
=
\I(Z;O\mid h_t,a)+(1-\alpha_t)\,\I(P;O\mid Z=0,h_t,a).
\]
\end{lemma}

\begin{proof}
Because $Z$ is a function of $P$, the chain rule gives
\[
\I(P;O\mid h_t,a)=\I(Z;O\mid h_t,a)+\I(P;O\mid Z,h_t,a).
\]
All programs in $c^\star$ induce the same next-step observation law, so $\I(P;O\mid Z=1,h_t,a)=0$. Taking the mixture over $Z$ at weights $\alpha_t,1-\alpha_t$ yields the claim.
\end{proof}

\begin{remark}
Lemma~\ref{lem:info-decomp} isolates the gap. Only the first term, information about membership in $[p^\star]$, is load-bearing for interventional semantic recovery. The second term can refine distinctions wholly inside $R_\pi(p^\star)\setminus[p^\star]$.
\end{remark}

\subsection{Near-miss: the AFE principle need not converge to $[p^\star]$}

Prior work on AIXI established that universal-prior learners can fail to identify the true environment in various senses \citep{LeikeHutter2015,Leike2016}, and that repair requires a specifically-chosen action rule such as Thompson sampling \citep{LeikeEtAl2016}. Theorem~\ref{thm:afe-near-miss} continues that line. It locates the failure of the same Bayes/Gibbs/expected-free-energy architecture, under a substituted finite-support prior, specifically at the level of the interventional semantic class $[p^\star]$, and shows that the expected-free-energy action rule is not a qualifying repair: the posterior on $[p^\star]$ can be frozen at its prior value for arbitrarily long horizons.

\begin{theorem}[Near-miss: the algorithmic free energy principle does not imply semantic convergence]
\label{thm:afe-near-miss}
Assume $\Ocal$ is infinite and $|\Acal|\ge 2$. For every $\alpha_0\in(0,1/2)$ and every horizon $T\ge 1$, there exist a computable interactive environment $p^\star\in\Pcal$ and a proper prior $w$ on $\Pcal$ supported on finitely many computable programs, with induced class prior $\bar w([p^\star])=\alpha_0$, such that the algorithmic free energy principle (Definition~\ref{def:afe-principle}) using $w$ in place of the universal interactive prior plays $a^{\mathrm{ref}}$ at every time $1,\dots,T$ on every sample path, and
\[
q_t^\star([p^\star]\mid H_t)=\alpha_0
\qquad\text{for every }t\in\{0,1,\dots,T\}.
\]
Since $T$ may be taken arbitrarily large and $\alpha_0\in(0,1/2)$ arbitrary, no rate of posterior convergence on $[p^\star]$ holds uniformly across algorithmic free energy learning problems.
\end{theorem}

The proof chooses $N+2$ distinct observations and two action modes: a ``separating'' action that discriminates $[p^\star]$ from its complement with at most $\log 2$ nats of mutual information, and a ``refining'' action that discriminates hypotheses within the complement of $[p^\star]$ with exactly $(1-\alpha_0)\log N$ nats. For $N$ large, expected free energy always prefers the refining action.

\begin{proof}
\emph{Construction.} Fix an integer $N\ge 3$ to be chosen below and the horizon $T\ge 1$. Choose distinct observations $R_1,\dots,R_N,S_0,S_1\in\Ocal$ and distinct actions $a^{\mathrm{ref}},a^{\mathrm{sep}}\in\Acal$.

Let $\Sigma_T:=\{1,\dots,N\}^T$. For each $\sigma\in\Sigma_T$, fix a computable program $p_\sigma\in\Pcal$ with action-conditioned law
\[
\mu_{p_\sigma}(o\mid h_{t-1},a)
=
\begin{cases}
\one\{o=R_{\sigma_t}\}, & a=a^{\mathrm{ref}},\ 1\le t\le T,\\
\one\{o=R_1\},         & a=a^{\mathrm{ref}},\ t>T,\\
\one\{o=S_0\},         & a\neq a^{\mathrm{ref}},
\end{cases}
\]
where the time index $t$ is read off from the length of $h_{t-1}$. Fix a computable program $p^\star\in\Pcal$ with law
\[
\mu_{p^\star}(o\mid h_{t-1},a)
=
\begin{cases}
N^{-1}\one\{o\in\{R_1,\dots,R_N\}\}, & a=a^{\mathrm{ref}},\\
\one\{o=S_1\},                        & a\neq a^{\mathrm{ref}}.
\end{cases}
\]
All such programs lie in $\Pcal$ under the universal prefix machine.

Write $c^\star:=[p^\star]$ and $c_\sigma:=[p_\sigma]$ for $\sigma\in\Sigma_T$. The classes $c^\star$ and $c_\sigma$ are distinct because the two programs disagree on $a^{\mathrm{sep}}$; classes $c_\sigma,c_{\sigma'}$ are distinct when $\sigma\neq\sigma'$ because the programs disagree at the first $t$ where $\sigma_t\neq\sigma'_t$ under $a^{\mathrm{ref}}$.

\emph{Prior.} Define a proper prior $w$ on $\Pcal$ by
\[
w(p^\star):=\alpha_0,\qquad w(p_\sigma):=\frac{1-\alpha_0}{N^T}\text{ for each }\sigma\in\Sigma_T,
\]
and $w\equiv 0$ on all other programs. This is a probability measure supported on $1+N^T$ computable programs. The induced class prior satisfies $\bar w(c^\star)=\alpha_0$ and $\bar w(c_\sigma)=(1-\alpha_0)/N^T$.

\emph{Learner.} Fix the algorithmic free energy principle of Definition~\ref{def:afe-principle} with the prior $w$ in place of the universal interactive prior, zero action cost, uniform preferences on the finite set $\{R_1,\dots,R_N,S_0,S_1\}$, and ties in action selection broken in favor of $a^{\mathrm{ref}}$.

\emph{Posterior on $c^\star$ under $a^{\mathrm{ref}}$-only transcripts.} Suppose the learner has played $a^{\mathrm{ref}}$ at times $1,\dots,t$ with observations $o_1,\dots,o_t\in\{R_1,\dots,R_N\}$ and $t\le T$. Encode $R_j\mapsto j$ to obtain $u\in\{1,\dots,N\}^t$. The number of $\sigma\in\Sigma_T$ matching $u$ on the first $t$ coordinates is exactly $N^{T-t}$, each with prior $(1-\alpha_0)/N^T$ and likelihood $1$ on the realized prefix. The program $p^\star$ has likelihood $N^{-t}$. Hence
\begin{equation}
\label{eq:nearmiss-alpha}
q_t^\star(c^\star\mid H_t)=\frac{\alpha_0\cdot N^{-t}}{\alpha_0\cdot N^{-t}+N^{T-t}\cdot (1-\alpha_0)/N^T}=\frac{\alpha_0}{\alpha_0+(1-\alpha_0)}=\alpha_0.
\end{equation}
The posterior on $c^\star$ equals its prior, on every sample path, as long as $t\le T$ and only $a^{\mathrm{ref}}$ has been played.

\emph{Conditional predictive among surviving $c_\sigma$'s.} The set of surviving $c_\sigma$'s at time $t$ is $\{\sigma\in\Sigma_T:\sigma_{1:t}=u\}$, of cardinality $N^{T-t}$. Because surviving classes have equal prior and equal likelihood, the posterior is uniform over this set. For $0\le t<T$, the conditional probability of $\sigma_{t+1}=j$ given the surviving set is
\[
\beta_t(j):=\frac{|\{\sigma\in\Sigma_T:\sigma_{1:t}=u,\,\sigma_{t+1}=j\}|}{N^{T-t}}=\frac{N^{T-t-1}}{N^{T-t}}=\frac{1}{N},
\qquad j\in\{1,\dots,N\}.
\]
So $\beta_t$ is \emph{exactly} uniform on $\{1,\dots,N\}$ for every $t\in\{0,1,\dots,T-1\}$ and every realized prefix $u$.

\emph{Mutual information at $a^{\mathrm{ref}}$.} Under $P=p^\star$ the next observation is uniform on $\{R_1,\dots,R_N\}$, so $\Hh(O\mid P=p^\star,h_t,a^{\mathrm{ref}})=\log N$. Under $P=p_\sigma$ (for any surviving $\sigma$) the next observation is $R_{\sigma_{t+1}}$ deterministically, so $\Hh(O\mid P=p_\sigma,h_t,a^{\mathrm{ref}})=0$. By \eqref{eq:nearmiss-alpha},
\[
\Hh(O\mid P,h_t,a^{\mathrm{ref}})=\alpha_0\log N+(1-\alpha_0)\cdot 0=\alpha_0\log N.
\]
The marginal predictive on $\{R_1,\dots,R_N\}$ is
\[
\lambda_t(j)=\alpha_0\cdot\frac{1}{N}+(1-\alpha_0)\cdot\beta_t(j)=\alpha_0\cdot\frac{1}{N}+(1-\alpha_0)\cdot\frac{1}{N}=\frac{1}{N},
\]
so $\Hh(O\mid h_t,a^{\mathrm{ref}})=\log N$. Therefore
\begin{equation}
\label{eq:mi-ref}
\I(P;O\mid h_t,a^{\mathrm{ref}})=\log N-\alpha_0\log N=(1-\alpha_0)\log N.
\end{equation}

\emph{Mutual information at $a^{\mathrm{sep}}$.} Under $a^{\mathrm{sep}}$, the observation is $S_1$ exactly when $P\in c^\star$ and $S_0$ otherwise. With $Z:=\one\{P\in c^\star\}$ under the posterior $q_t^\star$, the observation is a deterministic function of $Z$, so
\begin{equation}
\label{eq:mi-sep}
\I(P;O\mid h_t,a^{\mathrm{sep}})=\I(Z;O\mid h_t,a^{\mathrm{sep}})=\Hh(Z)=H_2(\alpha_0)\le\log 2,
\end{equation}
where $H_2(x):=-x\log x-(1-x)\log(1-x)$.

\emph{AFE plays $a^{\mathrm{ref}}$ through time $T$.} With zero action cost and uniform preferences on $\{R_1,\dots,R_N,S_0,S_1\}$, Lemma~\ref{lem:risk-ig} reduces expected-free-energy minimization to maximization of the mutual information $\I(P;O\mid h_t,a)$. Choose $N$ large enough that
\begin{equation}
\label{eq:nearmiss-N}
N>2^{1/(1-\alpha_0)},\qquad\text{equivalently }(1-\alpha_0)\log N>\log 2.
\end{equation}
Combining \eqref{eq:mi-ref}, \eqref{eq:mi-sep}, and \eqref{eq:nearmiss-N}, at every time $t\in\{0,\dots,T-1\}$ on which the learner has played $a^{\mathrm{ref}}$ exclusively,
\[
\I(P;O\mid h_t,a^{\mathrm{ref}})=(1-\alpha_0)\log N>\log 2\ge H_2(\alpha_0)=\I(P;O\mid h_t,a^{\mathrm{sep}}).
\]
The AFE rule selects $a^{\mathrm{ref}}$ (strictly; no tie-breaking is needed). Induction on $t$ shows that on every sample path, the learner plays $a^{\mathrm{ref}}$ at every time $1,\dots,T$.

\emph{Conclusion.} By \eqref{eq:nearmiss-alpha}, on every sample path,
\[
q_t^\star([p^\star]\mid H_t)=\alpha_0<1
\qquad\text{for every }t\in\{0,1,\dots,T\}.
\]
Since $T$ and $\alpha_0\in(0,1/2)$ were arbitrary, no rate of posterior convergence on $[p^\star]$ is uniform across AFE principle learning problems.
\end{proof}

\begin{theorem}[Observer-promotion failure of the AFE principle]
\label{thm:observer-promotion-failure}
Let $\pi^\sharp$ be the fixed policy that plays $a^{\mathrm{ref}}$ at every history. In the near-miss construction of Theorem~\ref{thm:afe-near-miss}, $\pi^\sharp$ agrees with the AFE principle's realized actions on every sample path through horizon $T$, and
\[
\omega_{\pi^\sharp}<\omega_{\mathrm{behav}}
\]
strictly in the refinement preorder.
\end{theorem}

\begin{proof}
Define $p'\in\Pcal$ by
\[
\mu_{p'}(o\mid h,a)=
\begin{cases}
\mu_{p^\star}(o\mid h,a^{\mathrm{ref}}), & a=a^{\mathrm{ref}},\\
\one\{o=S_0\}, & a\neq a^{\mathrm{ref}}.
\end{cases}
\]
Since $p^\star$'s $a^{\mathrm{ref}}$-law is computable, overriding its $a^{\mathrm{sep}}$-conditional with a fixed point-mass yields another program in $\Pcal$. Under $\pi^\sharp$ only $a^{\mathrm{ref}}$-pairs are reachable, and $\mu_{p'}$ agrees with $\mu_{p^\star}$ there, so $\Prob_{p',\pi^\sharp}=\Prob_{p^\star,\pi^\sharp}$ by the history-factorization identity; hence $p'\in R_{\pi^\sharp}(p^\star)$. The pair $(\varepsilon,a^{\mathrm{sep}})$ is jointly reachable, and $\mu_{p'}(\cdot\mid\varepsilon,a^{\mathrm{sep}})=\delta_{S_0}\neq\delta_{S_1}=\mu_{p^\star}(\cdot\mid\varepsilon,a^{\mathrm{sep}})$, so the kernel-equality convention of Section~\ref{sec:setup} gives $\mu_{p'}\neq\mu_{p^\star}$, i.e., $p'\notin[p^\star]$. Hence $R_{\pi^\sharp}(p^\star)\supsetneq[p^\star]$, which by Definition~\ref{def:observer} is $\omega_{\pi^\sharp}<\omega_{\mathrm{behav}}$.
\end{proof}

\begin{corollary}[Universal-prior case]
\label{cor:observer-promotion-universal}
The strict refinement gap $\omega_{\pi^\sharp}<\omega_{\mathrm{behav}}$ of Theorem~\ref{thm:observer-promotion-failure} is prior-independent. Replacing the near-miss's finite-support prior by the universal interactive prior does not alter this gap for the fixed policy $\pi^\sharp$.
\end{corollary}

\begin{proof}
The witness $p'$ and the policy $\pi^\sharp$ in Theorem~\ref{thm:observer-promotion-failure} are defined without reference to the prior. The same argument therefore yields the same strict refinement gap for $\pi^\sharp$ under the universal interactive prior.
\end{proof}

\begin{remark}
\label{rem:near-miss}
Theorem~\ref{thm:afe-near-miss} strengthens the near-miss from a one-step counterexample to an arbitrary-horizon convergence-failure family. Theorem~\ref{thm:observer-promotion-failure} reads this as a structural failure of action-side observer promotion: the AFE principle's realized observer $\omega_{\pi^\sharp}$ sits strictly below $\omega_{\mathrm{behav}}$, and the frozen-posterior conclusion is its Bayesian shadow. The belief rule is exact and the prior assigns positive mass to the true class; the failure is entirely on the action side, and a stricter action primitive --- the one Section~\ref{sec:main} supplies --- is required.
\end{remark}

\begin{remark}
\label{rem:near-miss-universal}
Theorem~\ref{thm:afe-near-miss} is proved for a substituted finite-support prior. Extending its exact frozen-posterior conclusion to the literal universal interactive prior would require a separate robustness argument controlling how the universal prior's additional support perturbs both the posterior values and the expected-free-energy action ranking. Corollary~\ref{cor:observer-promotion-universal} therefore isolates the prior-independent part of the failure: the strict refinement gap for the fixed policy $\pi^\sharp$ persists under the universal prior.
\end{remark}

\subsection{Observer-promotion contrast: AFE vs.\ promotion-supporting rule}
\label{sec:promotion-contrast}

The near-miss isolates the structural failure of the AFE boundary specialization. Sharpening the construction makes the contrast with the main theorem's promotion-supporting rule quantitative on a single learning problem: the two rules are run against the same prior, and the witness $p'$ of Theorem~\ref{thm:observer-promotion-failure} is put on the support.

\begin{corollary}[Observer-promotion contrast]
\label{cor:promotion-contrast}
In the near-miss construction of Theorem~\ref{thm:afe-near-miss}, replace the prior by one satisfying $w(p^\star)=\alpha_0$, $w(p')\in(0,1-\alpha_0)$ for the witness $p'$ of Theorem~\ref{thm:observer-promotion-failure}, and $w(p_\sigma)=(1-\alpha_0-w(p'))/N^T$ for each $\sigma\in\Sigma_T$, and choose $N>2^{1/(1-\alpha_0-w(p'))}$ in the construction. Then the posterior ratio $q_t^\star(p'\mid H_t)/q_t^\star(p^\star\mid H_t)$ is constant and equal to $w(p')/\alpha_0$ on every sample path for $t\le T$ under the AFE principle, while $q_t^\star(p'\mid H_t)\to 0$ almost surely under any promotion-supporting action rule satisfying (C1), (C2), and (C3).
\end{corollary}

\begin{proof}
Through horizon $T$ the AFE rule plays $\pi^\sharp$: the strengthened choice $N>2^{1/(1-\alpha_0-w(p'))}$ maintains $(1-\alpha_0-w(p'))\log N>\log 2\ge H_2(\alpha_0)$, so expected-free-energy's mutual information at $a^{\mathrm{ref}}$ strictly exceeds that at $a^{\mathrm{sep}}$ as in Theorem~\ref{thm:afe-near-miss}'s proof applied to the enlarged prior. Since $\mu_{p'}=\mu_{p^\star}$ at every reachable pair under $\pi^\sharp$, realized-history likelihoods coincide and the posterior ratio equals the prior ratio. For the promotion side, Theorem~\ref{thm:semantic-convergence} and Remark~\ref{rem:promotion-support} give $q_t^\star([p^\star]\mid H_t)\to 1$ almost surely under any promotion-supporting action rule satisfying (C1), (C2), and (C3). Because $p'\notin[p^\star]$, one has $q_t^\star(p'\mid H_t)\le 1-q_t^\star([p^\star]\mid H_t)\to 0$ almost surely.
\end{proof}

\begin{remark}
\label{rem:promotion-contrast-interpretation}
Corollary~\ref{cor:promotion-contrast} reads the knowability gap $[p^\star]\subsetneq R_{\pi^\sharp}(p^\star)$ as an action-rule test. The AFE rule respects the gap and cannot expel $p'$ from the posterior; a promotion-supporting rule probes the gap via $a^{\mathrm{sep}}$ and resolves it. Theorem~\ref{thm:policy-gap}'s set inclusion is exactly the structural object an observer-promoting action rule must close.
\end{remark}

\section{Implications}
\label{sec:implications}

The four-set hierarchy clarifies where existing learning theories stand. Each theory targets one of the four observers of Definition~\ref{def:observer}; the question for any given theory is whether its learning rule promotes the realized observer to the observer whose indistinguishability class the theory tries to identify.

Realized-history fit aims at $R_{h_t}(p^\star)$. Next-token training objectives under any fixed data distribution are instances: the learner scores well by assigning high likelihood to realized prefixes. Lemma~\ref{lem:fit-gap} is the boundary: such a learner's posterior can be on a compatibility set whose next-queried continuation is still undetermined. Predictive merging (Lemma~\ref{lem:merging}) promotes the posterior to the one-step-predictive interior of $R_{h_t}(p^\star)$ along the realized interaction.

Policy-relative prediction aims at $R_\pi(p^\star)$. Off-policy reinforcement learning and partially-observed Markov decision processes without active interventions face the boundary of Theorem~\ref{thm:policy-gap}: the posterior can concentrate on $R_\pi(p^\star)$ while leaving $[p^\star]$ unresolved.

Interventional semantic recovery aims at $[p^\star]$. It is the target of the main theorem and the theorem-level target of explanatory learning. The algorithmic free energy principle of Definition~\ref{def:afe-principle} is an exact universal-prior Bayesian learner with a generic epistemic action rule. Theorem~\ref{thm:afe-near-miss} shows that the same Bayes/Gibbs/expected-free-energy architecture, under a substituted finite-support prior, does not suffice to reach this target. Section~\ref{sec:semantic-action} identifies the stricter action primitive that does: increase the posterior odds of a semantic class against its complement. Concrete systems built around program-level structure --- Bayesian program learning and program induction \citep{LakeEtAl2015,EllisEtAl2021DreamCoder}, abstraction-and-reasoning benchmarks \citep{Chollet2024ARC}, and search over program spaces guided by large language models \citep{RomeraParedesEtAl2024FunSearch} --- can be read as instances of class-targeted learning in this sense, their semantic classes determined by the structure of the program space they search.

Canonical syntactic recovery aims at $\{p^\star\}$. Theorem~\ref{thm:factor-through-quotient} rules it out unless the designated representative is already constant on $[p^\star]$.

\subsection{A theorem-facing amortized surrogate}
\label{sec:amortized-surrogate}

The exact kernel lift $\mathcal J_{t,\mathrm{ker}}^{\omega_{\mathrm{behav}}}$ is the paper's strongest variational answer. A practical learner will typically replace the countable program space by a finite latent class model and optimize a differentiable surrogate instead. The theorem-level question is then not whether the surrogate reduces uncertainty in general, but whether minimizing it still certifies the separating-support condition of Theorem~\ref{thm:separating-support-convergence}.

Assume in this subsection that the deployed surrogate uses a finite latent class set $\widehat{\Ccal}$ and a finite observation alphabet $\Ocal$. Assume also that its parameter state is updated causally from history: for each $t$, the maps
\[
h_t\longmapsto \rho_t(\cdot\mid h_t)
\qquad\text{and}\qquad
h_t\longmapsto \widehat\mu_{\theta_t}(o\mid h_t,a,\hat c)
\]
are measurable for every $o\in\Ocal$, $a\in\Acal$, and $\hat c\in\widehat{\Ccal}$. In particular, any class-action kernel built from these quantities is an adapted policy kernel in the sense of the formal setup. Let $\rho_t(\hat c\mid h_t)$ be a differentiable full-support latent-class prior over $\widehat{\Ccal}$, and let $\widehat\mu_{\theta_t}(o\mid h_t,a,\hat c)$ be a differentiable predictive model. For each $\hat c\in\widehat{\Ccal}$ define the learned complement law
\[
\widehat Q_{t,\theta_t}^{-\hat c}(o\mid h_t,a)
:=
\sum_{\hat d\neq \hat c}
\frac{\rho_t(\hat d\mid h_t)}{1-\rho_t(\hat c\mid h_t)}
\widehat\mu_{\theta_t}(o\mid h_t,a,\hat d),
\]
the learned Bhattacharyya score
\[
\widehat B_{t,\theta_t}(\hat c,a;h_t)
:=
-\log\sum_{o\in\Ocal}
\sqrt{
\widehat\mu_{\theta_t}(o\mid h_t,a,\hat c)\,
\widehat Q_{t,\theta_t}^{-\hat c}(o\mid h_t,a)},
\]
and its cap
\[
\bar{\widehat B}_{t,\theta_t}(\hat c,a;h_t)
:=
\min\{1,\widehat B_{t,\theta_t}(\hat c,a;h_t)\}.
\]
For any differentiable predictive loss $\mathcal L_t^{\mathrm{pred}}(\theta_t;h_t)$ independent of the action kernel, define the amortized surrogate
\[
\mathcal L_{t,\mathrm{impl}}(\theta_t,\kappa;h_t)
:=
\mathcal L_t^{\mathrm{pred}}(\theta_t;h_t)
-\beta
\sum_{\hat c\in\widehat{\Ccal}}\sum_{a\in\Acal}
\kappa(\hat c,a\mid h_t)\,\bar{\widehat B}_{t,\theta_t}(\hat c,a;h_t)
+\gamma\,
\KL\!\bigl(\kappa(\cdot,\cdot\mid h_t)\,\|\,\rho_t(\cdot\mid h_t)\otimes\lambda\bigr),
\]
where $\kappa(\cdot,\cdot\mid h_t)\in\Delta(\widehat{\Ccal}\times\Acal)$ and $\lambda$ is a full-support reference distribution on $\Acal$.

\begin{proposition}[Exact action-side minimizer of the amortized surrogate]
\label{prop:amortized-surrogate-minimizer}
Fix $h_t$ and $\theta_t$. Any minimizer of $\mathcal L_{t,\mathrm{impl}}(\theta_t,\kappa;h_t)$ over $\kappa(\cdot,\cdot\mid h_t)\in\Delta(\widehat{\Ccal}\times\Acal)$ is
\[
\kappa_{t,\theta_t}^{\mathrm{impl}}(\hat c,a\mid h_t)
=
\frac{
\rho_t(\hat c\mid h_t)\lambda(a)\,
\exp\!\bigl((\beta/\gamma)\bar{\widehat B}_{t,\theta_t}(\hat c,a;h_t)\bigr)}
{\sum_{\hat d\in\widehat{\Ccal}}\sum_{b\in\Acal}
\rho_t(\hat d\mid h_t)\lambda(b)\,
\exp\!\bigl((\beta/\gamma)\bar{\widehat B}_{t,\theta_t}(\hat d,b;h_t)\bigr)}.
\]
\end{proposition}

\begin{proof}
For fixed $(\theta_t,h_t)$ the predictive term is constant in $\kappa$, so the minimization over $\kappa$ is exactly the Gibbs-kernel calculation of Proposition~\ref{prop:kernel-functional-minimizer} with $\rho_t(\cdot\mid h_t)$ in place of the exact class prior and $\bar{\widehat B}_{t,\theta_t}$ in place of the exact class score.
\end{proof}

\begin{theorem}[Semantic-recovery guarantee for the amortized surrogate]
\label{thm:amortized-surrogate}
Assume \textup{(C1)--(C3)} and $\omega_B=\omega_{\mathrm{syn}}$ with Bayes/Gibbs belief \textup{(C2)}. Fix $p^\star\in\Pcal$, write $c^\star:=[p^\star]$, and let the deployed policy be the action marginal
\[
\pi_{t+1}^{\mathrm{impl}}(a\mid h_t)
:=
\sum_{\hat c\in\widehat{\Ccal}}
\kappa_{t,\theta_t}^{\mathrm{impl}}(\hat c,a\mid h_t)
\]
of Proposition~\ref{prop:amortized-surrogate-minimizer}. Suppose the measurable adaptation hypothesis stated at the start of Section~\ref{sec:amortized-surrogate} holds. Suppose there exist a distinguished latent class $\hat c^\star\in\widehat{\Ccal}$ and constants
\[
s_0\in(0,1],
\qquad
\underline\rho>0,
\qquad
\underline\lambda>0
\]
such that for every reachable history $h_t$:
\begin{itemize}
\item[(A1)] $\rho_t(\hat c^\star\mid h_t)\ge \underline\rho$;
\item[(A2)] the learned high-score set
\[
\widehat S_t(h_t)
:=
\{a\in\Acal:\bar{\widehat B}_{t,\theta_t}(\hat c^\star,a;h_t)\ge s_0\}
\]
satisfies $\lambda(\widehat S_t(h_t))\ge \underline\lambda$;
\item[(A3)] every action in $\widehat S_t(h_t)$ is genuinely semantically separating for the true class:
\[
a\in \widehat S_t(h_t)
\implies
B_t(c^\star,a;h_t)\ge \eta(c^\star).
\]
\end{itemize}
Then for every reachable history $h_t$,
\[
\pi_{t+1}^{\mathrm{impl}}(S_t^\star\mid h_t)\ge \delta_{\mathrm{impl}},
\qquad
\delta_{\mathrm{impl}}
:=
\underline\rho\,\underline\lambda\,
\exp\!\bigl(-(\beta/\gamma)(1-s_0)\bigr),
\]
where
\[
S_t^\star
:=
\{a\in\Acal:B_t(c^\star,a;h_t)\ge \eta(c^\star)\}.
\]
Consequently,
\[
q_t^\star([p^\star]\mid H_t)\longrightarrow 1
\qquad
\text{almost surely.}
\]
\end{theorem}

\begin{proof}
Fix a reachable history $h_t$. By Proposition~\ref{prop:amortized-surrogate-minimizer},
\[
\kappa_{t,\theta_t}^{\mathrm{impl}}(\hat c^\star,\widehat S_t(h_t)\mid h_t)
=
\frac{
\sum_{a\in\widehat S_t(h_t)}
\rho_t(\hat c^\star\mid h_t)\lambda(a)\,
\exp\!\bigl((\beta/\gamma)\bar{\widehat B}_{t,\theta_t}(\hat c^\star,a;h_t)\bigr)}
{Z_t^{\mathrm{impl}}(h_t)},
\]
where
\[
Z_t^{\mathrm{impl}}(h_t)
:=
\sum_{\hat d\in\widehat{\Ccal}}\sum_{b\in\Acal}
\rho_t(\hat d\mid h_t)\lambda(b)\,
\exp\!\bigl((\beta/\gamma)\bar{\widehat B}_{t,\theta_t}(\hat d,b;h_t)\bigr).
\]
Since $\bar{\widehat B}_{t,\theta_t}\in[0,1]$, one has
\[
Z_t^{\mathrm{impl}}(h_t)\le e^{\beta/\gamma}
\sum_{\hat d\in\widehat{\Ccal}}\sum_{b\in\Acal}
\rho_t(\hat d\mid h_t)\lambda(b)
=
e^{\beta/\gamma},
\]
while assumptions \textup{(A1)} and \textup{(A2)} give
\[
\sum_{a\in\widehat S_t(h_t)}
\rho_t(\hat c^\star\mid h_t)\lambda(a)\,
\exp\!\bigl((\beta/\gamma)\bar{\widehat B}_{t,\theta_t}(\hat c^\star,a;h_t)\bigr)
\ge
\underline\rho\,\underline\lambda\,e^{(\beta/\gamma)s_0}.
\]
Therefore
\[
\kappa_{t,\theta_t}^{\mathrm{impl}}(\hat c^\star,\widehat S_t(h_t)\mid h_t)
\ge
\underline\rho\,\underline\lambda\,
\exp\!\bigl(-(\beta/\gamma)(1-s_0)\bigr)
=
\delta_{\mathrm{impl}}.
\]
Summing over $\hat c$ shows
\[
\pi_{t+1}^{\mathrm{impl}}(\widehat S_t(h_t)\mid h_t)\ge \delta_{\mathrm{impl}}.
\]
By \textup{(A3)}, $\widehat S_t(h_t)\subseteq S_t^\star$, so
\[
\pi_{t+1}^{\mathrm{impl}}(S_t^\star\mid h_t)\ge \delta_{\mathrm{impl}}
\]
for every reachable $h_t$. Applying Theorem~\ref{thm:separating-support-convergence} with $r_t\equiv \delta_{\mathrm{impl}}$ gives
\[
q_t^\star([p^\star]\mid H_t)\to 1
\qquad\text{almost surely.}
\]
\end{proof}

\begin{corollary}[Finite-time guarantee for the amortized surrogate]
\label{cor:amortized-surrogate-finite-time}
Under the assumptions of Theorem~\ref{thm:amortized-surrogate}, suppose in addition that the bounded-ratio hypothesis of Theorem~\ref{thm:separating-support-rate} holds for the deployed policy $\pi^{\mathrm{impl}}$. Then:
\begin{itemize}
\item[(i)] for every $t\ge 0$,
\[
\E[Z_t]\ge Z_0+2\,\eta(c^\star)\,\delta_{\mathrm{impl}}\,t;
\]
\item[(ii)] for every $t\ge 1$ and every $\delta\in(0,1)$,
\[
\Prob\!\left(
q_t^\star(c^\star\mid H_t)\ge
\frac{1}{1+\exp\!\bigl(-Z_0-2\,\eta(c^\star)\delta_{\mathrm{impl}}\,t+L\sqrt{2t\log(1/\delta)}\bigr)}
\right)\ge 1-\delta,
\]
where $L:=\log(1/m)$ is the bounded-ratio constant of Theorem~\ref{thm:separating-support-rate};
\item[(iii)] if for some $\varepsilon,\delta\in(0,1)$,
\[
2\,\eta(c^\star)\,\delta_{\mathrm{impl}}\,t
\ge
\log\frac{1-\varepsilon}{\varepsilon}-Z_0+L\sqrt{2t\log(1/\delta)},
\]
then
\[
\Prob\!\bigl(q_t^\star([p^\star]\mid H_t)\ge 1-\varepsilon\bigr)\ge 1-\delta.
\]
\end{itemize}
\end{corollary}

\begin{proof}
Apply Theorem~\ref{thm:separating-support-rate} and Corollary~\ref{cor:separating-support-finite-time} with $r_t\equiv \delta_{\mathrm{impl}}$.
\end{proof}

\begin{remark}
\label{rem:amortized-surrogate-scope}
The guarantee is intentionally conditional. It does \emph{not} say that stochastic gradient descent finds parameters satisfying \textup{(A1)--(A3)}, nor that the latent class $\hat c^\star$ is automatically the true semantic class. It says something narrower and theorem-facing: if the learned model keeps a designated target latent class alive, exposes a uniform positive-$\lambda$ region of high learned semantic score, and that high score is sound for genuine class-against-complement separation, then minimizing the surrogate action loss yields exactly the separating-support condition needed for semantic recovery and finite-time posterior lower bounds.
\end{remark}

\section{Discussion}
\label{sec:discussion}

The boxed problem has two logically distinct parts, and the paper resolves them at different levels. The belief rule is fully characterized: the Bayes/Gibbs posterior under the universal prior is the unique minimizer of the algorithmic free energy (Lemma~\ref{lem:variational}), and within a natural convex divergence class the KL-based functional is the only one with that property (Lemma~\ref{lem:kl-necessity}). The action rule is resolved up to a single master sufficiency condition (Theorem~\ref{thm:separating-support-convergence}), whose sharpness partitions the paper's failure results by mechanism. Theorem~\ref{thm:afe-near-miss} fails by using the wrong observer geometry: its action rule is generic information gain, correct in entropy but mismatched to class-against-complement separation. Theorems~\ref{cor:support-necessary} and~\ref{thm:summable-support-insufficient} fail in a different register --- the action rule targets the right observer, but the realized trajectory withholds enough mass from genuinely separating actions. The action objective supporting semantic recovery must therefore be class-targeted and sufficiently supported along the realized trajectory.

The observer frame places these action-rule results alongside three traditions. Pearl's do-calculus identifies when interventional distributions are determined by observational ones via graphical structure \citep{Pearl2009}; Corollary~\ref{cor:promotion-contrast} is its algorithmic analogue, naming the action-rule condition under which active play determines the interventional semantic class. \citet{Chernoff1959}'s asymptotically optimal sequential-testing rule selects, at each step, the design maximizing the separation of the most-likely hypothesis from the nearest alternative. The class-against-complement rule is its countable-hypothesis, universal-prior extension: promotion support (Definition~\ref{def:promotion-supporting}) plays the role of Chernoff's selection-probability condition, and the aggregate-complement law $\Qcal_t^{-c}$ replaces Chernoff's pairwise closest alternative when the rival class is structurally heterogeneous. \citet{Blackwell1953}'s comparison of experiments orders statistical designs by a sufficiency-type relation --- the parametric ancestor of Definition~\ref{def:observer}'s refinement preorder --- so Theorem~\ref{thm:semantic-convergence} reads as observer promotion in Blackwell's order along the learned trajectory.

The observer frame also sharpens three philosophical theses about underdetermination. \citet{Quine1960}'s claim that empirical evidence leaves logically incompatible theories equally supported is made formal by Theorem~\ref{thm:policy-gap}: whenever a fixed policy leaves some reachable history-action pair unqueried, there are concrete programs in $R_\pi(p^\star)\setminus[p^\star]$ that no amount of passive observation expels. \citet{vanFraassen1980}'s constructive empiricism holds that science aims at empirical adequacy rather than truth; Theorem~\ref{thm:factor-through-quotient} identifies the interactive analogue of empirical adequacy as $[p^\star]$, the finest target any history-based estimator can reliably recover. \citet{Putnam1980} questioned whether use could fix reference at all; Theorem~\ref{thm:semantic-convergence} shows that class-against-complement action fixes reference up to $[p^\star]$, though no farther. The observer formalism separates two kinds of underdetermination: the passive-active gap $[p^\star]\subsetneq R_\pi(p^\star)$, which active intervention closes via observer promotion, and the residue inside $[p^\star]$, which Theorem~\ref{thm:factor-through-quotient} shows no history-based observer can resolve.

Three theoretical questions remain open.

First, when do generic information gain and semantic gain coincide? Theorem~\ref{thm:afe-near-miss} shows that coincidence can fail, already in the substituted-prior family studied here. It is natural to ask for structural conditions on the hypothesis class under which expected-free-energy maximization automatically maximizes semantic gain. The answer would identify the precise boundary between generic epistemic discovery and explanatory action.

Second, what quantitative rates can be proved? Theorem~\ref{thm:separating-support-rate} and Corollary~\ref{cor:separating-support-finite-time} already give explicit finite-time sufficiency under the bounded-log-likelihood-ratio regularity hypothesis, and Propositions~\ref{prop:kernel-exp-rate} and~\ref{prop:exp-rate}, together with Theorem~\ref{thm:exp-rate-concentration}, identify the selector and kernel specializations of those support-based bounds. Corollary~\ref{cor:compact-action-kernel} shows that the same kernel-rate bounds persist for compact metric action spaces once the neighborhood-mass floor is recovered from continuity and full support. What remains open is removing or substantially weakening that regularity hypothesis while retaining comparable finite-horizon lower bounds on $q_t^\star([p^\star]\mid h_t)$. Sequential design results in statistics suggest such bounds should be possible in restricted regimes \citep{Chernoff1959,BoxHill1967}.

Third, what tractable approximations preserve the theorem-level target? Section~\ref{sec:amortized-surrogate} records a theorem-facing surrogate answer: a finite latent world model plus the differentiable loss $\mathcal L_{t,\mathrm{impl}}$ inherits the same semantic-recovery and finite-time guarantees once assumptions \textup{(A1)--(A3)} certify that the learned high-score region is nonvanishing and semantically sound. What remains open is the harder approximation question: for which model classes, training procedures, and amortized inference schemes can one prove those assumptions from the optimization itself rather than impose them as deployment-side hypotheses?

\section{Conclusion}
\label{sec:conclusion}

The four nested sets $\{p^\star\}\subseteq[p^\star]\subseteq R_\pi(p^\star)\subseteq R_{h_t}(p^\star)$ are the indistinguishability classes of four observers at increasing observational power, with all three inclusions strict in general. The belief problem is solved exactly: the universal prior and Bayes/Gibbs posterior minimize the algorithmic free energy, and at zero action cost with uniform preferences the resulting expected free energy is the information-gain criterion shared by Bayesian experimental design, compression-based curiosity, and active-inference curiosity. Even this architecture can fail: the near-miss theorem (Theorem~\ref{thm:afe-near-miss}, Theorem~\ref{thm:observer-promotion-failure}) exhibits a substituted-prior family in which the realized policy induces an observer strictly coarser than the behavioral one, so the posterior on $[p^\star]$ need not concentrate. Theorem~\ref{thm:separating-support-convergence} answers the boxed problem: under the universal prior, Bayes/Gibbs belief, semantic separation, and divergent cumulative policy mass on separating actions, the posterior on $[p^\star]$ converges almost surely to one, with explicit finite-time high-probability lower bounds under bounded log-likelihood ratios (Theorem~\ref{thm:separating-support-rate}). Three canonical realizations attain that condition under the standing finite-action setup: the class-against-complement selector (Theorem~\ref{thm:semantic-convergence}), the exact Gibbs kernel lift as the strongest exact variational realization (Theorem~\ref{thm:kernel-semantic-convergence}), and a differentiable finite-latent-class surrogate (Theorem~\ref{thm:amortized-surrogate}). A divergent cumulative exploration budget suffices on its own (Theorem~\ref{thm:exploration-floor-behavioral}), and Corollary~\ref{cor:compact-action-kernel} extends the kernel route to compact metric action spaces. The condition is sharp: zero separating support makes uniform recovery impossible (Corollary~\ref{cor:support-necessary}) and summable separating support still admits failure with positive probability (Theorem~\ref{thm:summable-support-insufficient}), so the paper isolates two distinct ways semantic recovery can fail --- wrong-observer geometry and insufficient support on separating interventions.

\emph{Observer promotion} is the organizing informal concept: each learning rule attempts to push the realized observer toward the behavioral one, and the four-observer hierarchy grades that attempt. The self-referential reduction (Section~\ref{sec:self-ref}) replaces the uncomputable behavioral observer by a finite-time proxy and recovers the main guarantee under eventual isolation (deterministic splitting or divergent exploration).

The public Lean 4 repository at \url{https://github.com/Populustremuloides/semantic_convergence} contains counterparts for all 106 manuscript items over the explicit concrete foundation used in the repo, and publishes generated coverage, proof-shape, and per-declaration axiom audits documenting the verification state. On the variational side, \path{SemanticConvergence.def_afe}, \path{SemanticConvergence.def_two_observer_functional}, and \path{SemanticConvergence.def_kernel_functional} realize the repaired generalized-KL / Gibbs belief, class, and kernel objectives directly, while \path{SemanticConvergence.lem_variational}, \path{SemanticConvergence.lem_kl_necessity}, \path{SemanticConvergence.prop_two_observer_minimizer}, and \path{SemanticConvergence.prop_kernel_functional_minimizer} certify the corresponding exact Bayes/Gibbs variational identity, countable convex-duality KL-necessity theorem, and Gibbs minimizer results on the paper-facing countable stack. In Section~\ref{sec:main}, the soft theorem is now represented explicitly: \path{SemanticConvergence.HellingerConvergenceSpine} packages the residual nonnegativity, L1-bounded martingale envelope, divergent cumulative Bhattacharyya separation, and square-root Bayes identity; \path{SemanticConvergence.thm_separating_support_convergence_hellinger_spine} proves posterior-share convergence from that package; and \path{SemanticConvergence.thm_kernel_functional_minimizer_convergence_of_hellinger_spine} composes exact kernel-functional minimization with the Hellinger convergence spine. The concrete declarations \path{SemanticConvergence.thm_separating_support_convergence}, \path{SemanticConvergence.thm_separating_support_rate}, and \path{SemanticConvergence.cor_separating_support_finite_time} remain as a stronger zero-out support/rate route with finite-horizon bounds. The manifest's stricter \texttt{fullyFirstPrinciples} flag remains false because the final internal construction of the Bayes/Gibbs-induced Hellinger spine from semantic separation and realized policy support is still an exactness lock. At the boundary and surrogate end, \path{SemanticConvergence.thm_afe_near_miss}, \path{SemanticConvergence.thm_amortized_surrogate_selector_existence}, \path{SemanticConvergence.thm_amortized_surrogate}, and \path{SemanticConvergence.cor_amortized_surrogate_finite_time} expose the repaired action-level near-miss, countable \textup{(A1)--(A3)} wrapper, deployment-side support floor, and finite-time consequence. For any interactive learner, classical or modern, the theorem reduces the question ``is this aimed at the right target?'' to a single checkable criterion: divergent cumulative policy mass on separating actions.

\bibliographystyle{plainnat}
\bibliography{semantic_convergence_interactive_learning}

\end{document}
